% Formuler les idées essentielles d’un texte à résumer — Français 1re Tech — Cours signé OnBuch+
% Programme officiel MINESEC. Compiler avec : tectonic N15-formuler-idees-essentielles-texte-resumer.tex
\newcommand{\DOCMATIERE}{Français}
\newcommand{\DOCNIVEAU}{1re Tech}
\newcommand{\DOCMODULE}{Français – classe de Première technique}
\newcommand{\DOCLECON}{N15}
\newcommand{\DOCTITRE}{Formuler les idées essentielles d’un texte à résumer}
% OnBuch+ — préambule « cours signés » ENRICHI (compilé avec tectonic / XeLaTeX)
% Système visuel « Pop » d'OnBuch+ : fond crème (jamais blanc), bordures encre
% épaisses, ombres « dures » décalées, titres Archivo Black en capitales.
% Version MATHÉMATIQUES : graphes (pgfplots/tikz), boîtes pédagogiques
% (définition, propriété, méthode, attention, le savais-tu, exercice résolu,
% exercice, TP, bilan), page de garde et signature OnBuch+ ancrée en bas.
%
% Macros attendues dans main.tex AVANT \input{preamble} :
%   \DOCMATIERE  \DOCNIVEAU  \DOCMODULE  \DOCLECON (numéro)  \DOCTITRE
\documentclass[11pt]{article}

\usepackage{fontspec}
\usepackage[french]{babel}
\usepackage[a4paper,margin=2cm,top=3.4cm,bottom=2.8cm,headheight=1.4cm,footskip=1.3cm]{geometry}
\usepackage{amsmath,amssymb,amsfonts}
\usepackage{mathtools} % \xmapsto, \coloneqq... souvent utilisés par les modèles
\usepackage{cancel} % simplifications barrées (bilans de forces)
% \abs{z} (module, valeur absolue) et \norm{u} (norme), utilisés par les modèles
\providecommand{\abs}{}\let\abs\relax\DeclarePairedDelimiter{\abs}{\lvert}{\rvert}
\providecommand{\norm}{}\let\norm\relax\DeclarePairedDelimiter{\norm}{\lVert}{\rVert}
\usepackage[table]{xcolor} % [table] : fournit aussi \rowcolors (alternance de lignes)
\usepackage{pagecolor}
\usepackage{tikz}
\usetikzlibrary{arrows.meta,positioning,calc,shapes.geometric,shapes.misc,shapes.arrows,decorations.pathmorphing,decorations.pathreplacing,decorations.markings,patterns,shadows,mindmap,trees,fit,backgrounds,matrix,intersections,angles,quotes}
\usepackage{pgfplots}
\usepackage[version=4]{mhchem} % équations nucléaires \ce{^{235}_{92}U}
\usepackage[european,siunitx]{circuitikz} % schémas électriques
\pgfplotsset{compat=1.18}
\usepgfplotslibrary{fillbetween,groupplots} % aires entre courbes (calcul intégral)
\usepackage{enumitem}
\usepackage{fancyhdr}
% [most] charge toutes les bibliothèques tcolorbox (listings, minted, raster,
% documentation...) et gonfle inutilement la mémoire TeX sur un document long
% (~300 boîtes) : on ne charge que ce qui est réellement utilisé.
\usepackage[skins,breakable]{tcolorbox}
\usepackage{graphicx}
\usepackage{colortbl}
\usepackage{booktabs}
\usepackage{tabularx}
\usepackage{multirow}
\usepackage{diagbox} % case d'en-tête diagonale des tableaux à double entrée
% Colonnes extensibles centrée / gauche / droite (C, L, R), souvent utilisées par les modèles
\newcolumntype{C}{>{\centering\arraybackslash}X}
\newcolumntype{L}{>{\raggedright\arraybackslash}X}
\newcolumntype{R}{>{\raggedleft\arraybackslash}X}
\usepackage{array}
\usepackage{multicol}
\usepackage{siunitx}
% parse-numbers=false : accepte aussi \SI{2,5 \times 10^{-3}}{...}
\sisetup{locale=FR,output-decimal-marker={,},group-minimum-digits=4,parse-numbers=false}
\DeclareSIUnit\ohm{\ensuremath{\symup{\Omega}}} % U+2126 absent de la police
\DeclareSIUnit\division{div} % réglages d'oscilloscope (ms/div, V/div)
\DeclareSIUnit\tr{tr} % tours (vitesse de rotation, ex. \SI{1500}{\tr\per\minute})
\DeclareSIUnit\molar{M} % concentration molaire (ex. \SI{3}{\molar}, \SI{1}{\micro\molar})
\DeclareSIUnit\an{an} % année (le modèle confond parfois avec \year, un compteur TeX) — ex. \SI{5}{\kilo\meter\per\an}
\DeclareSIUnit\FCFA{FCFA} % franc CFA (ex. \SI{500}{\FCFA\per\kilo\gram})
\DeclareSIUnit\degreeBrix{\textdegree Brix} % teneur en sucre d'un sirop/jus (réfractométrie)
\DeclareSIUnit\month{mois} % le modèle utilise parfois \month, un compteur TeX, comme unité de durée
\DeclareSIUnit\microliter{\micro\liter} % le modèle écrit parfois \microliter en un seul token
\DeclareSIUnit\million{millions} % dénombrement (ex. \SI{15}{\million\per\mL} spermatozoïdes)
\usepackage{qrcode}
% \degree : commande fréquemment utilisée par les modèles pour un angle ou une
% température en dehors de \SI{}{} (non fournie par les paquets chargés).
\providecommand{\degree}{^{\circ}}
% \qed (fin de démonstration), utilisé par les modèles sans amsthm
\providecommand{\qed}{\hfill$\square$}
% \barre : double barre des valeurs interdites dans un tableau de variations (tkz-tab)
\providecommand{\barre}{\ensuremath{\|}}
% environnement proof (amsthm non chargé) utilisé par les modèles
\ifdefined\proof\else\newenvironment{proof}[1][Démonstration]{\par\noindent\textit{#1.}\ }{\qed\par}\fi
\usepackage{lastpage}
\usepackage{hyperref}
\hypersetup{hidelinks}

% ============================================================
% Palette « Pop » OnBuch+ — mêmes hex que la classe P (plus_ui.dart)
% ============================================================
\definecolor{poporange}{HTML}{F59321}
\definecolor{popdark}{HTML}{E07A0C}
\definecolor{popcream}{HTML}{FFF6E8}
\definecolor{popink}{HTML}{1C1714}
\definecolor{poppurple}{HTML}{7A5AE0}
\definecolor{popgreen}{HTML}{2E9E6A}
\definecolor{poppink}{HTML}{E0457B}
\definecolor{popblue}{HTML}{2F7DE1}
\definecolor{popgold}{HTML}{C98A12}
\definecolor{popmuted}{HTML}{8A7F76}
\definecolor{popline}{HTML}{EADFD2}
\definecolor{popgray}{HTML}{8A7F76}
\colorlet{popgrey}{popgray}
% Alias tolérants (le modèle nomme parfois les couleurs différemment)
\colorlet{popwhite}{white}
\colorlet{popblack}{popink}
\colorlet{popred}{poppink}
\colorlet{popyellow}{popgold}
% Teintes claires (fonds de tableaux/encadrés) — le modèle nomme parfois
% les couleurs avec un suffixe L (ex. popblueL) sans les avoir définies.
\colorlet{poporangeL}{poporange!20}
\colorlet{popdarkL}{popdark!20}
\colorlet{poppurpleL}{poppurple!20}
\colorlet{popgreenL}{popgreen!20}
\colorlet{poppinkL}{poppink!20}
\colorlet{popblueL}{popblue!20}
\colorlet{popgoldL}{popgold!20}
\colorlet{popredL}{popred!20}
\colorlet{popgrayL}{popgray!20}
% Teintes claires pour fonds de tableaux / zones
\colorlet{popblueL}{popblue!10!white}
\colorlet{poporangeL}{poporange!14!white}
\colorlet{popgreenL}{popgreen!12!white}
\colorlet{poppurpleL}{poppurple!12!white}
\colorlet{poppinkL}{poppink!12!white}
\colorlet{popgoldL}{popgold!14!white}

\pagecolor{popcream}

% ============================================================
% Typographie
% ============================================================
\defaultfontfeatures{Ligatures=TeX}
\setmainfont{PlusJakartaSans-Medium.ttf}[Path=../../../fonts/,BoldFont=PlusJakartaSans-Bold.ttf]
% Maths en sans-serif (Fira Math) avec chiffres et lettres droites de Plus
% Jakarta, pour que les formules se fondent dans le texte.
\usepackage[math-style=ISO,bold-style=ISO]{unicode-math}
\setmathfont{FiraMath-Regular.otf}[Path=../../../fonts/]
\setmathfont{PlusJakartaSans-Medium.ttf}[Path=../../../fonts/,range={up/num,up/{Latin,latin},\mathup}]
\usepackage{icomma} % virgule décimale sans espace en mode math
% Caractères mathématiques Unicode absents des polices de texte (Plus Jakarta,
% Archivo Black) : rendus par leur équivalent mathématique, en texte comme en
% formule — sinon ils s'affichent en carré vide (ex. « Arithmétique dans ℤ »).
\usepackage{newunicodechar}
\newunicodechar{ℝ}{\ensuremath{\mathbb{R}}}
\newunicodechar{ℤ}{\ensuremath{\mathbb{Z}}}
\newunicodechar{ℂ}{\ensuremath{\mathbb{C}}}
\newunicodechar{ℕ}{\ensuremath{\mathbb{N}}}
\newunicodechar{ℚ}{\ensuremath{\mathbb{Q}}}
\newunicodechar{𝔻}{\ensuremath{\mathbb{D}}}
\newunicodechar{α}{\ensuremath{\alpha}}
\newunicodechar{β}{\ensuremath{\beta}}
\newunicodechar{γ}{\ensuremath{\gamma}}
\newunicodechar{δ}{\ensuremath{\delta}}
\newunicodechar{ε}{\ensuremath{\varepsilon}}
\newunicodechar{θ}{\ensuremath{\theta}}
\newunicodechar{λ}{\ensuremath{\lambda}}
\newunicodechar{μ}{\ensuremath{\mu}}
\newunicodechar{σ}{\ensuremath{\sigma}}
\newunicodechar{τ}{\ensuremath{\tau}}
\newunicodechar{φ}{\ensuremath{\varphi}}
\newunicodechar{ψ}{\ensuremath{\psi}}
\newunicodechar{ω}{\ensuremath{\omega}}
\newunicodechar{Σ}{\ensuremath{\Sigma}}
% majuscules produites par \MakeUppercase dans les titres (α-aminés -> Α)
\newunicodechar{Α}{\ensuremath{\alpha}}
\newunicodechar{Β}{\ensuremath{\beta}}
\newunicodechar{Γ}{\ensuremath{\Gamma}}
\newunicodechar{Δ}{\ensuremath{\Delta}}
\newunicodechar{Ε}{\ensuremath{\varepsilon}}
\newunicodechar{Θ}{\ensuremath{\Theta}}
\newunicodechar{Λ}{\ensuremath{\Lambda}}
\newunicodechar{Μ}{\ensuremath{\mu}}
\newunicodechar{Τ}{\ensuremath{\tau}}
\newunicodechar{Φ}{\ensuremath{\Phi}}
\newunicodechar{Ψ}{\ensuremath{\Psi}}
\newunicodechar{Ω}{\ensuremath{\Omega}}
\newunicodechar{↦}{\ensuremath{\mapsto}}
\newunicodechar{∈}{\ensuremath{\in}}
\newunicodechar{∉}{\ensuremath{\notin}}
\newunicodechar{∧}{\ensuremath{\wedge}}
\newunicodechar{⇔}{\ensuremath{\Leftrightarrow}}
\newunicodechar{⟺}{\ensuremath{\Longleftrightarrow}}
\newunicodechar{⇒}{\ensuremath{\Rightarrow}}
\newunicodechar{⟹}{\ensuremath{\Longrightarrow}}
\newunicodechar{∘}{\ensuremath{\circ}}
\newunicodechar{∪}{\ensuremath{\cup}}
\newunicodechar{∩}{\ensuremath{\cap}}
\newunicodechar{≡}{\ensuremath{\equiv}}
\newunicodechar{⊂}{\ensuremath{\subset}}
\newunicodechar{⊥}{\ensuremath{\perp}}
\newunicodechar{∃}{\ensuremath{\exists}}
\newunicodechar{∀}{\ensuremath{\forall}}
\newunicodechar{∛}{\ensuremath{\sqrt[3]{\,}}}
\newunicodechar{∜}{\ensuremath{\sqrt[4]{\,}}}
\newunicodechar{⃗}{\ensuremath{{}^{\to}}}
\newunicodechar{ⁿ}{\ifmmode^{n}\else\textsuperscript{\ensuremath{n}}\fi}
\newunicodechar{ˣ}{\ifmmode^{x}\else\textsuperscript{\ensuremath{x}}\fi}
\newunicodechar{ᵇ}{\ifmmode^{b}\else\textsuperscript{\ensuremath{b}}\fi}
\newunicodechar{ʳ}{\ifmmode^{r}\else\textsuperscript{\ensuremath{r}}\fi}
\newunicodechar{ᵘ}{\ifmmode^{u}\else\textsuperscript{\ensuremath{u}}\fi}
\newunicodechar{ʸ}{\ifmmode^{y}\else\textsuperscript{\ensuremath{y}}\fi}
\newunicodechar{ᵃ}{\ifmmode^{a}\else\textsuperscript{\ensuremath{a}}\fi}
\newunicodechar{ᵉ}{\ifmmode^{e}\else\textsuperscript{\ensuremath{e}}\fi}
\newunicodechar{ᵏ}{\ifmmode^{k}\else\textsuperscript{\ensuremath{k}}\fi}
\newunicodechar{ᵖ}{\ifmmode^{p}\else\textsuperscript{\ensuremath{p}}\fi}
\newunicodechar{ᴺ}{\ifmmode^{N}\else\textsuperscript{\ensuremath{N}}\fi}
\newunicodechar{ᶜ}{\ifmmode^{c}\else\textsuperscript{\ensuremath{c}}\fi}
\newunicodechar{ʰ}{\ifmmode^{h}\else\textsuperscript{\ensuremath{h}}\fi}
\newunicodechar{ᵠ}{\ifmmode^{q}\else\textsuperscript{\ensuremath{q}}\fi}
\newunicodechar{ₙ}{\ifmmode_{n}\else\textsubscript{\ensuremath{n}}\fi}
\newunicodechar{ₐ}{\ifmmode_{a}\else\textsubscript{\ensuremath{a}}\fi}
\newunicodechar{ᵢ}{\ifmmode_{i}\else\textsubscript{\ensuremath{i}}\fi}
\newunicodechar{ᵣ}{\ifmmode_{r}\else\textsubscript{\ensuremath{r}}\fi}
\newunicodechar{ₖ}{\ifmmode_{k}\else\textsubscript{\ensuremath{k}}\fi}
\newunicodechar{ᵦ}{\ifmmode_{\beta}\else\textsubscript{\ensuremath{\beta}}\fi}
\newunicodechar{ₓ}{\ifmmode_{x}\else\textsubscript{\ensuremath{x}}\fi}
\newfontfamily\popdisplay{ArchivoBlack-Regular.ttf}[Path=../../../fonts/]
\newfontfamily\poplogo{SpaceGrotesk-Bold.ttf}[Path=../../../fonts/]
\newfontfamily\popsemi{PlusJakartaSans-SemiBold.ttf}[Path=../../../fonts/]
\newfontfamily\popxbold{PlusJakartaSans-ExtraBold.ttf}[Path=../../../fonts/]
\newfontfamily\popxboldls{PlusJakartaSans-ExtraBold.ttf}[Path=../../../fonts/,LetterSpace=3.0]

\setlist[enumerate]{leftmargin=1.7em,itemsep=5pt,topsep=4pt}
\setlist[itemize]{leftmargin=1.7em,itemsep=4pt,topsep=4pt,label={\color{poporange}\textbullet}}
\setlength{\parindent}{0pt}
\setlength{\parskip}{5pt}
\renewcommand{\arraystretch}{1.25}

% pgfplots : style Pop par défaut
\pgfplotsset{
  every axis/.append style={axis line style={popink,line width=1pt},tick style={popink},
    grid style={popline,line width=0.5pt},label style={font=\small},tick label style={font=\footnotesize}},
  popcurve/.style={poporange,line width=1.8pt,no markers,smooth},
}
% Même style disponible dans un \draw TikZ simple (hors pgfplots)
\tikzset{popcurve/.style={poporange,line width=1.8pt}}
\tikzset{popgrid/.style={popline,very thin}}
\colorlet{popcurve}{poporange} % le nom du style est parfois utilisé comme couleur

% ============================================================
% Pill / squiggle
% ============================================================
\newcommand{\poppill}[3]{%
  \tikz[baseline=-0.55ex]\node[draw=popink,line width=1.2pt,rounded corners=2.6pt,%
    fill=#2,inner xsep=6.5pt,inner ysep=3pt]{%
    \popxboldls\fontsize{7}{7}\selectfont\color{#3}\MakeUppercase{#1}};%
}
\newcommand{\popsquiggle}[1]{%
  \tikz[baseline=-0.4ex]{
    \draw[#1,line width=2.6pt,line cap=round]
      plot [smooth] coordinates {(0,0.09) (0.34,0.01) (0.68,0.21) (1.02,0.01) (1.36,0.10)};
  }%
}
% Mot-clé surligné (marqueur orange sous le texte)
\newcommand{\cle}[1]{{\popxbold\color{popdark}#1}}

% ============================================================
% En-tête / pied
% ============================================================
\pagestyle{fancy}
\fancyhf{}
\renewcommand{\headrulewidth}{0pt}
\renewcommand{\footrulewidth}{0pt}
\lhead{\raisebox{-0.15ex}{\poplogo\fontsize{13}{13}\selectfont\color{popink}OnBuch\color{poporange}+}}
\rhead{\poppill{\DOCMATIERE\ \textbullet\ \DOCNIVEAU\ \textbullet\ Leçon \DOCLECON}{white}{popink}}
\lfoot{\popxbold\fontsize{7.6}{7.6}\selectfont\color{popmuted}\MakeUppercase{Cours signé OnBuch+}\ \textbullet\ {\popsemi\DOCTITRE}}
\rfoot{\tikz[baseline=-0.6ex]\node[rounded corners=8.5pt,fill=popink,minimum height=17pt,minimum width=17pt,inner xsep=5pt,inner sep=0]{\popxbold\fontsize{8}{8}\selectfont\color{popcream}\thepage\,/\,\pageref*{LastPage}};}

% ============================================================
% Titres
% ============================================================
\newcommand{\chaptitre}[2]{%
  \begin{tcolorbox}[enhanced,colback=poporange,colframe=popink,boxrule=2.6pt,arc=11pt,
      drop shadow=popink,left=15pt,right=15pt,top=12pt,bottom=11pt,before skip=3pt,after skip=18pt]
    \poppill{#2}{popink}{popcream}\par\vspace{9pt}
    {\raggedright\hyphenpenalty=10000\exhyphenpenalty=10000\popdisplay\fontsize{22}{24}\selectfont\color{popcream}\MakeUppercase{#1}\par}
  \end{tcolorbox}
}
% Section numérotée : pastille orange avec le numéro + titre Archivo
\newcounter{popsec}
\newcommand{\coursec}[1]{%
  \stepcounter{popsec}\setcounter{popsub}{0}%
  \par\vspace{16pt}\noindent
  \begin{minipage}{\linewidth}
  \tikz[baseline=-0.7ex]\node[draw=popink,line width=1.6pt,fill=poporange,rounded corners=4pt,minimum size=22pt,inner sep=0]{\popdisplay\fontsize{12}{12}\selectfont\color{popcream}\thepopsec};%
  \hspace{8pt}{\popdisplay\fontsize{15}{17}\selectfont\color{popink}\MakeUppercase{#1}}\par
  \vspace{4pt}\noindent{\color{poporange}\rule{36pt}{3.6pt}}
  \end{minipage}\par\nopagebreak\vspace{8pt}
}
\newcounter{popsub}
\newcommand{\courssub}[1]{%
  \stepcounter{popsub}%
  \par\vspace{9pt}\noindent{\popxbold\fontsize{11.5}{13}\selectfont\color{popink}{\color{poporange}\thepopsec.\thepopsub}\hspace{6pt}#1}\par\nopagebreak\vspace{4pt}
}

% ============================================================
% Boîtes pédagogiques — même recette : fond blanc, bordure épaisse colorée,
% ombre dure encre, badge plein. Titre optionnel [#1].
% ============================================================
\tcbset{popbase/.style={breakable,enhanced,colback=white,boxrule=2.3pt,arc=9pt,
  shadow={3pt}{-3pt}{0pt}{popink},left=14pt,right=14pt,top=11pt,bottom=11pt,before skip=11pt,after skip=15pt,
  fontupper=\popsemi\fontsize{10.6}{14.5}\selectfont\color{popink},
  fontlower=\popsemi\fontsize{10.6}{14.5}\selectfont\color{popink}}}
% Coche dessinée (le glyphe ✓ n'existe pas dans les polices du document)
\newcommand{\popcheck}[1]{\tikz[baseline=-0.5ex]\draw[#1,line width=1.6pt,line cap=round,line join=round](-0.13,0.0)--(-0.04,-0.09)--(0.13,0.1);}
\providecommand{\coche}{\popcheck{popgreen}}
\providecommand{\resultat}[1]{\par\smallskip\noindent\fcolorbox{popgreen}{popgreenL}{\parbox{\dimexpr\linewidth-2\fboxsep-2\fboxrule\relax}{\textbf{Résultat.} #1}}\par\smallskip}
\newcommand{\popboxhead}[3]{\poppill{#1}{#2}{white}\ifx\relax#3\relax\else\hspace{6pt}{\popxbold\fontsize{10.6}{12}\selectfont\color{popink}#3}\fi\par\vspace{7pt}}

\newtcolorbox{aretenir}[1][]{popbase,colframe=popgreen,before upper={\popboxhead{À retenir}{popgreen}{#1}}}
\newtcolorbox{exemplebox}[1][]{popbase,colframe=poporange,before upper={\popboxhead{Exemple}{poporange}{#1}}}
\newtcolorbox{definition}[1][]{popbase,colframe=popblue,before upper={\popboxhead{Définition}{popblue}{#1}}}
\newtcolorbox{propriete}[1][]{popbase,colframe=poppurple,before upper={\popboxhead{Propriété}{poppurple}{#1}}}
\newtcolorbox{methode}[1][]{popbase,colframe=popgold,before upper={\popboxhead{Méthode}{popgold}{#1}}}
\newtcolorbox{attention}[1][]{popbase,colframe=poppink,before upper={\popboxhead{Attention}{poppink}{#1}}}
\newtcolorbox{experience}[1][]{popbase,colframe=popblue,colback=popblueL,before upper={\popboxhead{Expérience}{popblue}{#1}}}
% « Le savais-tu ? » : carte violette pleine (contexte, culture, Cameroun)
\newtcolorbox{savaistu}[1][]{popbase,colback=poppurple,colframe=popink,
  fontupper=\popsemi\fontsize{10.4}{14.2}\selectfont\color{white},
  before upper={\poppill{Le savais-tu\,?}{white}{poppurple}\ifx\relax#1\relax\else\hspace{6pt}{\popxbold\color{white}#1}\fi\par\vspace{7pt}}}
% Exercice résolu : énoncé en haut, \tcblower puis la solution sur fond crème
\newcounter{popexo}\newcounter{popexr}
\newtcolorbox{exoresolu}[1][]{popbase,colframe=poporange,colbacklower=popcream,
  segmentation style={popink,line width=1.2pt,dashed},
  before upper={\stepcounter{popexr}\popboxhead{Exercice résolu \thepopexr}{poporange}{#1}},
  before lower={\poppill{Solution}{popink}{popcream}\par\vspace{6pt}}}
% Exercice (non résolu en place) — la correction est dans la partie Corrigés
\newtcolorbox{exercice}[1][]{popbase,colframe=popink,
  before upper={\stepcounter{popexo}\popboxhead{Exercice \thepopexo}{popink}{#1}}}
% Corrigé
\newtcolorbox{corrige}[1][]{popbase,colframe=popgreen,colback=popgreenL,
  before upper={\popboxhead{Corrigé}{popgreen}{#1}}}
% Objectifs / prérequis / bilan
\newtcolorbox{objectifs}{popbase,colframe=popink,colback=poporangeL,
  before upper={\popboxhead{Objectifs de la leçon}{popink}{}}}
\newtcolorbox{prerequis}{popbase,colframe=popink,colback=popblueL,
  before upper={\popboxhead{Prérequis}{popblue}{}}}
\newtcolorbox{bilan}{popbase,colframe=popink,colback=white,boxrule=2.8pt,
  before upper={\popboxhead{Fiche bilan}{poporange}{L'essentiel à connaître pour l'examen}}}
% Figure encadrée avec légende
\newtcolorbox{popfigure}[1][]{enhanced,breakable=false,colback=white,colframe=popline,boxrule=1.4pt,arc=8pt,
  left=10pt,right=10pt,top=10pt,bottom=8pt,before skip=10pt,after skip=12pt,halign=center,
  lower separated=false,halign lower=center,
  fontlower=\popsemi\fontsize{9}{11.5}\selectfont\color{popmuted},
  #1}
\newcommand{\legende}[1]{\par\vspace{6pt}{\popsemi\fontsize{9}{11.5}\selectfont\color{popmuted}#1}}

% ============================================================
% Signature OnBuch+
% ============================================================
\newcommand{\poplogoinline}[1]{{\poplogo\fontsize{#1}{#1}\selectfont\color{popink}OnBuch\color{poporange}+}}
\newcommand{\signatureonbuch}{%
  \begin{tcolorbox}[enhanced,colback=popink,colframe=popink,boxrule=2.2pt,arc=10pt,
      drop shadow=poporange,left=14pt,right=14pt,top=10pt,bottom=10pt,before skip=0pt,after skip=0pt]
    \tikz[baseline=-0.7ex]\node[circle,fill=popgreen,draw=popcream,line width=1.3pt,minimum size=22pt,inner sep=0]{\popcheck{white}};%
    \hspace{9pt}\begin{minipage}[c]{0.62\linewidth}
      {\popxbold\fontsize{12}{13}\selectfont\color{popcream}Signé\ \textbullet\ {\poplogo OnBuch\color{poporange}+}}\par\vspace{2pt}
      {\raggedright\popsemi\fontsize{8}{10.5}\selectfont\color{popline}\MakeUppercase{Équipe pédagogique OnBuch+}\\\MakeUppercase{Conforme au programme officiel MINESEC}\par}
    \end{minipage}\hfill
    {\poplogo\fontsize{22}{22}\selectfont\color{popcream}OnBuch\color{poporange}+}
  \end{tcolorbox}%
}

% ============================================================
% Page de garde
% #1 titre · #2 matière · #3 niveau · #4 module · #5 n° leçon · #6 durée ·
% #7 description / objectifs (texte ou itemize)
% ============================================================
\newcommand{\pagegarde}[7]{%
  \thispagestyle{empty}
  \newgeometry{a4paper,margin=1.7cm,top=1.6cm,bottom=1.4cm}%
  \begin{tikzpicture}[remember picture,overlay]
    \fill[poporange!18] ([xshift=-3.2cm,yshift=-3.6cm]current page.north east) circle (4.6cm);
    \fill[poppurple!14] ([xshift=2.4cm,yshift=7.6cm]current page.south west) circle (3.8cm);
  \end{tikzpicture}%
  {\poplogo\fontsize{30}{30}\selectfont\color{popink}OnBuch\color{poporange}+}\hfill
  \raisebox{0.6ex}{\poppill{Cours signé \textbullet\ Premium}{popink}{popcream}}\par
  \vspace{1pt}\popsquiggle{poppurple}\par
  \vspace{8pt}
  \begin{tcolorbox}[enhanced,colback=poporange,colframe=popink,boxrule=2.8pt,arc=13pt,
      drop shadow=popink,left=18pt,right=18pt,top=12pt,bottom=13pt,after skip=10pt]
    \poppill{#2 \textbullet\ #3}{popink}{popcream}\hspace{5pt}\poppill{#4}{white}{popink}\par\vspace{8pt}
    {\popdisplay\fontsize{14}{14}\selectfont\color{popink}LEÇON #5}\par\vspace{4pt}
    {\raggedright\hyphenpenalty=10000\exhyphenpenalty=10000\popdisplay\fontsize{25}{27}\selectfont\color{popcream}\MakeUppercase{#1}\par}
    \vspace{8pt}
    {\popxbold\fontsize{9.5}{9.5}\selectfont\color{popink}\MakeUppercase{Durée conseillée : #6}}
  \end{tcolorbox}
  \begin{tcolorbox}[enhanced,colback=white,colframe=popink,boxrule=2.2pt,arc=10pt,
      drop shadow=popink,left=16pt,right=16pt,top=10pt,bottom=10pt,after skip=8pt]
    \poppill{Dans cette leçon}{poppurple}{white}\par\vspace{5pt}
    \popsemi\fontsize{9.3}{12.6}\selectfont\color{popink}#7
  \end{tcolorbox}
  \noindent\begin{minipage}[t]{0.487\linewidth}
    \begin{tcolorbox}[enhanced,colback=popgreen,colframe=popink,boxrule=2.2pt,arc=10pt,
        drop shadow=popink,left=11pt,right=11pt,top=10pt,bottom=10pt,height=3.1cm,valign=center]
      \tcbox[colback=white,colframe=popink,boxrule=1.3pt,arc=5pt,boxsep=0pt,nobeforeafter,
        left=3pt,right=3pt,top=3pt,bottom=3pt,valign=center]{\qrcode[height=1.9cm]{https://play.google.com/store/apps/details?id=com.luvvix.onbuch&hl=fr}}%
      \hspace{8pt}\begin{minipage}[c]{0.5\linewidth}
        {\popdisplay\fontsize{11}{12.5}\selectfont\color{white}\MakeUppercase{Télécharge OnBuch}}\par\vspace{4pt}
        {\raggedright\popsemi\fontsize{8.4}{11}\selectfont\color{white}Gratuit \textbullet\ Google Play\\Tuteur IA, annales, concours, résultats\par}
      \end{minipage}
    \end{tcolorbox}
  \end{minipage}\hfill
  \begin{minipage}[t]{0.487\linewidth}
    \begin{tcolorbox}[enhanced,colback=poppurple,colframe=popink,boxrule=2.2pt,arc=10pt,
        drop shadow=popink,left=11pt,right=11pt,top=10pt,bottom=10pt,height=3.1cm,valign=center]
      \tikz[baseline=-0.7ex]\node[circle,fill=white,draw=popink,line width=1.3pt,minimum size=1.6cm,inner sep=0]{\popdisplay\fontsize{24}{24}\selectfont\color{poppurple}+};%
      \hspace{8pt}\begin{minipage}[c]{0.6\linewidth}
        {\popdisplay\fontsize{11}{12.5}\selectfont\color{white}\MakeUppercase{Passe à OnBuch+}}\par\vspace{4pt}
        {\raggedright\popsemi\fontsize{8.4}{11}\selectfont\color{white}Tous les cours signés, Tuteur IA illimité, fiches PDF — onglet OnBuch+ de l'app\par}
      \end{minipage}
    \end{tcolorbox}
  \end{minipage}
  \vspace{8pt}
  \popcontenu
  \vfill
  \signatureonbuch
  \restoregeometry
  \newpage
}

% Rangée « Ce cours contient » (page de garde)
\newcommand{\poptile}[3]{% #1 couleur #2 gros texte #3 légende
  \begin{tcolorbox}[enhanced,colback=#1,colframe=popink,boxrule=2pt,arc=9pt,shadow={2.5pt}{-2.5pt}{0pt}{popink},
      left=6pt,right=6pt,top=9pt,bottom=9pt,halign=center,valign=center,height=1.75cm,width=\linewidth,nobeforeafter]
    {\popdisplay\fontsize{12.5}{13}\selectfont\color{white}\MakeUppercase{#2}}\par\vspace{3pt}
    {\centering\popsemi\fontsize{7.8}{9.5}\selectfont\color{white}#3\par}
  \end{tcolorbox}}
\newcommand{\popcontenu}{%
  \par\noindent{\popxboldls\fontsize{7.5}{7.5}\selectfont\color{popmuted}CE COURS SIGNÉ CONTIENT}\par\vspace{7pt}
  \noindent\begin{minipage}[t]{0.235\linewidth}\poptile{popblue}{Cours}{complet, illustré, pas à pas}\end{minipage}\hfill
  \begin{minipage}[t]{0.235\linewidth}\poptile{poporange}{Méthodes}{et exercices résolus}\end{minipage}\hfill
  \begin{minipage}[t]{0.235\linewidth}\poptile{poppink}{Exercices}{type examen + corrigés}\end{minipage}\hfill
  \begin{minipage}[t]{0.235\linewidth}\poptile{popgreen}{Fiche bilan}{l'essentiel à retenir}\end{minipage}\par}

% Sommaire
\newtcolorbox{sommaire}{enhanced,colback=white,colframe=popink,boxrule=2.2pt,arc=10pt,
  drop shadow=popink,left=18pt,right=18pt,top=13pt,bottom=13pt,before skip=6pt,after skip=20pt,
  before upper={\poppill{Sommaire}{poppurple}{white}\par\vspace{9pt}\popsemi\fontsize{10.6}{14.5}\selectfont\color{popink}}}

% Signature de fin de cours — poussée tout en bas de la dernière page
\newcommand{\findecours}{%
  \par\vfill\nopagebreak
  \begin{center}\popsquiggle{poporange}\end{center}\vspace{4pt}
  \signatureonbuch
}
\AtBeginDocument{\def\llbracket{\mathopen{[\mkern-3.5mu[}}\def\rrbracket{\mathclose{]\mkern-3.5mu]}}} % intervalles d'entiers (glyphe ⟦ absent de la police)
% Glyphes absents de Fira Math : redéfinis à partir de glyphes disponibles
\AtBeginDocument{%
  \def\setminus{\mathbin{\backslash}}\let\smallsetminus\setminus
  \def\vdots{\mathord{\vbox{\baselineskip3.2pt\lineskiplimit0pt\kern4pt\hbox{.}\hbox{.}\hbox{.}}}}%
  \def\ddots{\mathinner{\mkern1mu\raise6.5pt\hbox{.}\mkern2mu\raise3.6pt\hbox{.}\mkern2mu\raise0.7pt\hbox{.}\mkern1mu}}%
  \def\triangle{\mathord{\scalebox{0.9}{$\symup{\Delta}$}}}%
  \def\nabla{\mathord{\raisebox{\depth}{\scalebox{1}[-1]{$\symup{\Delta}$}}}}%
  \def\checkmark{\popcheck{popdark}}%
  \def\star{\mathbin{\ast}}\def\bigstar{\mathord{\ast}}%
  \def\diamond{\mathbin{\scalebox{0.6}{$\Diamond$}}}%
  \def\bigcup{\mathop{\mathchoice{\raisebox{-0.3em}{\scalebox{1.7}{$\cup$}}}{\scalebox{1.25}{$\cup$}}{\cup}{\cup}}\displaylimits}%
  \def\bigcap{\mathop{\mathchoice{\raisebox{-0.3em}{\scalebox{1.7}{$\cap$}}}{\scalebox{1.25}{$\cap$}}{\cap}{\cap}}\displaylimits}%
  \def\lesseqqgtr{\mathrel{\lessgtr}}\def\bigoplus{\mathop{\oplus}\displaylimits}%
}
\providecommand{\ssi}{si et seulement si} % abréviation fréquente
\AtBeginDocument{\providecommand{\wideparen}{\overparen}} % arc de cercle

%% FIGFIX-BEGIN v5 — correctifs globaux des figures (docs/onbuch-generation/tools/figfix)
\usepackage{adjustbox}\usepackage{etoolbox}\tcbuselibrary{raster}
% toute figure TikZ/pgfplots plus large que la ligne est réduite à la largeur disponible
\BeforeBeginEnvironment{tikzpicture}{\begin{adjustbox}{max width=\linewidth}}
\AfterEndEnvironment{tikzpicture}{\end{adjustbox}}
% légendes pgfplots lisibles par-dessus les courbes
\pgfplotsset{every axis/.append style={legend style={fill=white,fill opacity=0.92,text opacity=1,draw=black!15,inner sep=3pt}}}
% étiquettes d'arêtes (syntaxe edge["..."]) avec fond blanc
\tikzset{every edge quotes/.append style={fill=white,inner sep=1.2pt}}
%% FIGFIX-END
\begin{document}

\pagegarde{Formuler les idées essentielles d’un texte à résumer}{Français}{1re Tech}{Français – classe de Première technique}{N15}{10 séances (fiche de progression harmonisée)}{Cette leçon apprend à identifier la structure argumentative d'un texte, à en dégager les idées essentielles et à les reformuler fidèlement dans un résumé. Elle s'appuie sur la lecture méthodique du Lion et la perle et mobilise les outils de la langue (tonalités, paronymie, antonymie) pour affiner la compréhension et l'expression.
\vspace{4pt}
\begin{multicols}{2}\small
\begin{itemize}
\item À la fin de cette leçon, je sais identifier les caractéristiques d'un texte argumentatif (thèse, arguments, exemples, connecteurs).
\item À la fin de cette leçon, je sais repérer la structure argumentative d'un texte (introduction, développement, conclusion).
\item À la fin de cette leçon, je sais appliquer les techniques de réduction : suppression, sélection, généralisation, fusion.
\item À la fin de cette leçon, je sais reformuler les idées essentielles d'un texte avec mes propres mots sans déformer la pensée de l'auteur.
\item À la fin de cette leçon, je sais reconnaître les tonalités comique, dramatique et satirique dans un texte littéraire.
\item À la fin de cette leçon, je sais distinguer les paronymes et utiliser les antonymes pour nuancer un propos.
\end{itemize}\end{multicols}}

\begin{sommaire}
\begin{multicols}{2}
\begin{enumerate}
\item Le texte argumentatif : caractéristiques
\item La structure argumentative et les techniques de réduction
\item Lecture méthodique du Lion et la perle (n° 4 et 5)
\item Les tonalités : comique, dramatique, satirique
\item Reformuler les idées essentielles
\item Paronymie, antonymie et exposé oral
\item Activité d'intégration
\item Méthodes et exercices résolus
\item Exercices
\item Corrigés des exercices
\item Fiche bilan
\end{enumerate}
\end{multicols}
\end{sommaire}

\begin{objectifs}
\begin{itemize}
\item À la fin de cette leçon, je sais identifier les caractéristiques d'un texte argumentatif (thèse, arguments, exemples, connecteurs).
\item À la fin de cette leçon, je sais repérer la structure argumentative d'un texte (introduction, développement, conclusion).
\item À la fin de cette leçon, je sais appliquer les techniques de réduction : suppression, sélection, généralisation, fusion.
\item À la fin de cette leçon, je sais reformuler les idées essentielles d'un texte avec mes propres mots sans déformer la pensée de l'auteur.
\item À la fin de cette leçon, je sais reconnaître les tonalités comique, dramatique et satirique dans un texte littéraire.
\item À la fin de cette leçon, je sais distinguer les paronymes et utiliser les antonymes pour nuancer un propos.
\item À la fin de cette leçon, je sais préparer et présenter un exposé oral structuré devant la classe.
\item À la fin de cette leçon, je sais rédiger et justifier une démarche de résumé conforme aux exigences du Baccalauréat technique.
\end{itemize}
\end{objectifs}

\begin{prerequis}
\begin{itemize}
\item Notion de résumé vue en classe de Seconde (réduction quantitative d'un texte narratif)
\item Types de textes et schéma de la communication (4e/3e)
\item Connecteurs logiques et organisation du paragraphe
\item Lectures méthodiques n° 1 à 3 du Lion et la perle
\item Champs lexicaux et relations de sens (synonymie, hyperonymie)
\end{itemize}
\end{prerequis}

\newpage

\coursec{Le texte argumentatif : caractéristiques}

\courssub{Situation d'accroche et problématique}

Lors des débats à la radio CRTV ou dans les meetings politiques à Yaoundé, les discours durent parfois une heure : pourtant, un journaliste doit en rendre compte en trois minutes, et un élève doit résumer un article du \emph{Cameroon Tribune} en quelques lignes à l'examen. Comment distinguer l'essentiel de l'accessoire ? Comment dire la même chose avec ses propres mots sans trahir l'auteur ? C'est tout l'art de la \cle{contraction de texte}, compétence décisive au Probatoire et au Baccalauréat technique.

Mais avant de réduire un texte, il faut d'abord comprendre \emph{comment il est construit}. Or la plupart des textes proposés à l'examen (éditoriaux, articles de presse, discours, extraits d'essais) appartiennent à une même famille : le \cle{texte argumentatif}. Toute cette leçon repose donc sur une question centrale : \textbf{qu'est-ce qu'un texte argumentatif et comment est-il organisé ?}

\courssub{Définition du texte argumentatif : convaincre et persuader}

Imaginons une scène de la vie quotidienne. À Maroua, une jeune élève veut obtenir de son père l'autorisation de poursuivre ses études au lycée. Elle ne se contente pas de dire : « Je veux aller au lycée. » Elle donne des \emph{raisons} : « Les filles instruites aident mieux leur famille ; le lycée n'est qu'à deux kilomètres ; j'ai obtenu d'excellents résultats. » Sans le savoir, elle produit un raisonnement argumentatif : elle défend une position à l'aide de raisons.

\begin{definition}[Le texte argumentatif]
Un \cle{texte argumentatif} est un texte qui défend ou rejette une \cle{thèse} (une opinion, une position) au moyen d'\cle{arguments} (des raisons) étayés d'\cle{exemples} (des cas précis), dans le but de faire adhérer un destinataire à cette position.
\end{definition}

L'auteur d'un texte argumentatif poursuit deux objectifs voisins mais distincts :

\begin{itemize}
\item \textbf{Convaincre} : s'adresser à la \emph{raison} du destinataire. L'auteur utilise des arguments logiques, des faits vérifiables, des statistiques, des démonstrations. Le lecteur adhère parce qu'il est intellectuellement persuadé.
\item \textbf{Persuader} : s'adresser aux \emph{émotions} du destinataire. L'auteur utilise des images frappantes, des témoignages touchants, un vocabulaire valorisant ou péjoratif, l'ironie. Le lecteur adhère parce qu'il est émotionnellement touché.
\end{itemize}

\begin{popfigure}
\begin{center}
\begin{tikzpicture}[node distance=0.9cm,
  box/.style={draw=popblue, thick, rounded corners, fill=popblueL, text width=5.6cm, align=center, minimum height=1.6cm, font=\small},
  lab/.style={font=\small\bfseries, text=popdark}]
  \node[lab] (titre) at (0,0) {Faire adhérer le destinataire à une thèse};
  \node[box, below left=1.1cm and -0.6cm of titre, align=center] (conv){\textbf{CONVAINCRE}\\ s'adresser à la raison\\ faits, chiffres, logique, démonstration};
  \node[box, below right=1.1cm and -0.6cm of titre, draw=poporange, fill=poporangeL, align=center] (pers){\textbf{PERSUADER}\\ s'adresser aux émotions\\ images, témoignages, ironie, vocabulaire affectif};
  \draw[-{Stealth}, thick, popblue] (titre.south) -- (conv.north);
  \draw[-{Stealth}, thick, poporange] (titre.south) -- (pers.north);
\end{tikzpicture}
\end{center}
\legende{Tableau comparatif : convaincre fait appel à la raison, persuader fait appel à l'émotion. Un même texte combine souvent les deux procédés.}
\end{popfigure}

\begin{exemplebox}[Convaincre ou persuader ?]
Compare ces deux phrases tirées d'un débat sur la scolarisation des filles :
\begin{itemize}
\item « Selon les statistiques du ministère, une fille scolarisée jusqu'au secondaire fait reculer la pauvreté de sa famille. » $\rightarrow$ l'auteur \emph{convainc} (chiffre, fait vérifiable).
\item « Fermer l'école à une fille, c'est éteindre une lumière qui aurait éclairé tout un village. » $\rightarrow$ l'auteur \emph{persuade} (image, émotion).
\end{itemize}
\end{exemplebox}

\courssub{Les éléments constitutifs du texte argumentatif}

Un texte argumentatif n'est pas un amas de phrases : c'est une \emph{construction} ordonnée, comparable à une maison. La thèse est le toit que tout l'édifice soutient ; les arguments sont les piliers ; les exemples sont les briques qui consolident chaque pilier.

\begin{aretenir}[Les quatre éléments du texte argumentatif]
\begin{enumerate}
\item La \cle{thèse} : la position de l'auteur, l'idée qu'il défend (\emph{thèse soutenue}) ou qu'il combat (\emph{thèse rejetée}).
\item Les \cle{arguments} : les raisons générales qui justifient la thèse.
\item Les \cle{exemples} : les cas précis, faits, anecdotes ou chiffres qui illustrent et prouvent chaque argument.
\item La \cle{conclusion} : le rappel de la thèse, parfois accompagné d'un appel à l'action ou d'une ouverture.
\end{enumerate}
\end{aretenir}

\begin{popfigure}
\begin{center}
\begin{tikzpicture}[
  these/.style={draw=popdark, very thick, rounded corners, fill=popgoldL, text width=6.5cm, align=center, font=\small\bfseries},
  arg/.style={draw=popblue, thick, rounded corners, fill=popblueL, text width=4.2cm, align=center, font=\small},
  ex/.style={draw=popgreen, thick, rounded corners, fill=popgreenL, text width=4.2cm, align=center, font=\small},
  concl/.style={draw=popdark, very thick, rounded corners, fill=poppinkL, text width=6.5cm, align=center, font=\small\bfseries},
  fleche/.style={-{Stealth}, thick, popdark}]
  \node[these, align=center] (t) at (0,0){THÈSE\\ « Il faut scolariser toutes les jeunes filles »};
  \node[arg, align=center] (a1) at (-3.4,-2.2){ARGUMENT 1\\ L'instruction libère la femme};
  \node[arg, align=center] (a2) at (3.4,-2.2){ARGUMENT 2\\ L'instruction enrichit la famille};
  \node[ex] (e1) at (-3.4,-4.2) {Exemple : les femmes diplômées de Garoua créent des entreprises};
  \node[ex] (e2) at (3.4,-4.2) {Exemple : les mères instruites font réussir leurs enfants};
  \node[concl, align=center] (c) at (0,-6.2){CONCLUSION\\ rappel de la thèse + appel à l'action};
  \draw[fleche] (t.south) -- (a1.north);
  \draw[fleche] (t.south) -- (a2.north);
  \draw[fleche] (a1.south) -- (e1.north);
  \draw[fleche] (a2.south) -- (e2.north);
  \draw[fleche] (e1.south) |- (c.west);
  \draw[fleche] (e2.south) |- (c.east);
\end{tikzpicture}
\end{center}
\legende{Schéma en arbre de la structure argumentative : la thèse est soutenue par des arguments, chacun étayé par un exemple ; l'ensemble aboutit à la conclusion.}
\end{popfigure}

Précisons la différence entre la \emph{thèse soutenue} et la \emph{thèse rejetée}. Dans un débat, deux positions s'affrontent. L'auteur peut choisir de défendre directement sa position (thèse soutenue : « le téléphone portable nuit aux études ») ou de combattre la position adverse (thèse rejetée : « non, le téléphone portable n'est pas un outil de travail scolaire »). Souvent, un bon texte argumentatif présente d'abord la thèse adverse pour mieux la réfuter : c'est la stratégie de la \emph{réfutation}.

\courssub{Les marques linguistiques de l'argumentation}

Comment reconnaître un texte argumentatif à coup sûr ? Grâce à des \emph{indices de langue} visibles, comme des panneaux indicateurs sur une route.

\textbf{1. Les connecteurs logiques.} Ce sont des mots de liaison qui organisent le raisonnement :

\begin{center}
\begin{tabularx}{\linewidth}{|l|X|}
\hline
\rowcolor{popblueL} \textbf{Fonction} & \textbf{Connecteurs} \\
\hline
Énumération des arguments & d'abord, ensuite, enfin, de plus, en outre \\
\hline
Explication, justification & en effet, car, parce que, puisque, c'est pourquoi \\
\hline
Opposition, concession & mais, cependant, pourtant, néanmoins, certes... mais \\
\hline
Conséquence, conclusion & donc, ainsi, par conséquent, c'est pourquoi, finalement \\
\hline
Illustration & par exemple, c'est le cas de, notamment, c'est ainsi que \\
\hline
\end{tabularx}
\end{center}

\textbf{2. Les modalisateurs.} Ce sont les mots qui expriment le degré de certitude ou le jugement de l'auteur : \emph{il faut, il est évident que, sans doute, certes, heureusement, malheureusement, peut-être}. Ils révèlent l'engagement de celui qui écrit.

\textbf{3. Les pronoms personnels.} Le \emph{je} engage l'auteur (« je soutiens que... ») ; le \emph{nous} associe le lecteur à la cause (« nous devons agir ») ; le \emph{on} généralise. Le \emph{vous} interpelle directement le destinataire.

\begin{methode}[Trouver la thèse d'un texte]
\begin{enumerate}
\item Lis le texte entier une première fois pour saisir le sujet général.
\item Repère la phrase qui résume la position de l'auteur : elle est souvent marquée par « il faut », « nous pensons que », « il est urgent de », « je suis convaincu que ».
\item Vérifie que tous les arguments du texte soutiennent bien cette phrase : si oui, c'est la thèse.
\item Formule la thèse en une phrase claire, avec tes propres mots si l'exercice le demande.
\end{enumerate}
\end{methode}

\courssub{Argumentation directe et argumentation indirecte}

L'argumentation peut emprunter deux voies :

\begin{itemize}
\item L'\cle{argumentation directe} : l'auteur expose clairement sa thèse et ses arguments, en son nom propre. Genres concernés : l'\emph{essai}, l'\emph{éditorial}, l'\emph{article de presse}, le \emph{discours}, le \emph{pamphlet}, la \emph{lettre ouverte}.
\item L'\cle{argumentation indirecte} : l'auteur fait passer son message à travers une fiction, des personnages, une histoire. Le lecteur doit dégager lui-même la leçon. Genres concernés : le \emph{conte}, la \emph{fable}, le \emph{roman} (comme \emph{Le Lion et la perle}), le \emph{théâtre}, la \emph{satire}.
\end{itemize}

\begin{exemplebox}[Les deux voies sur un même sujet]
Pour dénoncer la corruption, un éditorialiste écrira : « La corruption détruit notre économie ; il faut la combattre » (argumentation directe). Un fabuliste racontera l'histoire d'un rat devenu chef qui dévore les réserves du village et finit chassé par tous (argumentation indirecte). Le message est le même ; le chemin diffère.
\end{exemplebox}

\courssub{Fait ou opinion ? Argument ou exemple ?}

Pour analyser un texte argumentatif, deux distinctions sont indispensables.

\textbf{Le fait et l'opinion.} Un \cle{fait} est une réalité vérifiable, observable, que personne ne peut contester (« Maroua est le chef-lieu de la région de l'Extrême-Nord »). Une \cle{opinion} est un jugement personnel, discutable, qui engage celui qui l'émet (« Maroua est la plus belle ville du Cameroun »). Le texte argumentatif mélange les deux : il s'appuie sur des faits pour défendre des opinions.

\textbf{L'argument et l'exemple.} L'\emph{argument} est une raison \emph{générale} ; l'\emph{exemple} est un cas \emph{particulier} qui illustre cette raison. L'argument répond à la question « pourquoi ? », l'exemple répond à la question « comment le voit-on ? ».

\begin{attention}[Erreur fréquente à l'examen]
Ne confonds pas un \emph{exemple} (cas précis) avec un \emph{argument} (raison générale). Dans la phrase « L'instruction des filles favorise la santé des enfants : à Ngaoundéré, par exemple, les mères instruites font vacciner tous leurs enfants », l'argument est « l'instruction des filles favorise la santé des enfants » ; « le cas des mères de Ngaoundéré » n'est qu'un exemple. À la question « citez deux arguments », un exemple ne rapporte aucun point.
\end{attention}

\courssub{Application : repérage de la thèse dans un éditorial}

Lisons ce court éditorial inspiré du \emph{Cameroon Tribune}, consacré à la scolarisation des jeunes filles à Maroua, puis appliquons la méthode.

\begin{exemplebox}[Éditorial : « L'école pour toutes les filles »]
« À Maroua, trop de jeunes filles quittent encore l'école avant la fin du primaire. Certains parents estiment qu'une fille n'a pas besoin d'études pour réussir sa vie. \textbf{Cette idée est fausse et dangereuse : il faut scolariser toutes les jeunes filles.} D'abord, \textbf{l'instruction donne à la femme les moyens de son indépendance} : ainsi, les jeunes diplômées de l'Extrême-Nord créent des commerces et emploient d'autres femmes. Ensuite, \textbf{une mère instruite fait mieux réussir ses enfants} : les enquêtes montrent que les élèves dont la mère a fréquenté le secondaire obtiennent de meilleurs résultats. Il est donc urgent que les autorités et les familles unissent leurs efforts pour garder nos filles à l'école. »
\end{exemplebox}

\begin{exoresolu}[Dégager la thèse et les arguments de l'éditorial]
\textbf{Énoncé.} 1. Dégagez la thèse défendue par l'éditorialiste. 2. Relevez les deux arguments qui la soutiennent et l'exemple qui illustre le premier argument. 3. L'argumentation est-elle directe ou indirecte ? Justifiez.
\tcblower
\textbf{Solution détaillée.}
\begin{enumerate}
\item \textbf{Thèse} : « Il faut scolariser toutes les jeunes filles » (l'éditorialiste rejette la thèse adverse selon laquelle une fille n'a pas besoin d'études). Le marqueur « il faut » et la formule « cette idée est fausse et dangereuse » signalent la position de l'auteur.
\item \textbf{Argument 1} : l'instruction donne à la femme les moyens de son indépendance (raison générale). \textbf{Exemple} : les jeunes diplômées de l'Extrême-Nord qui créent des commerces (cas précis, introduit par « ainsi »). \textbf{Argument 2} : une mère instruite fait mieux réussir ses enfants (étayé par les résultats des enquêtes).
\item Il s'agit d'une \textbf{argumentation directe} : l'auteur expose sa position clairement, en son nom, dans un éditorial, avec des connecteurs logiques visibles (« d'abord », « ensuite », « donc »).
\end{enumerate}
\end{exoresolu}

\begin{savaistu}[La thèse, une question qui rapporte]
Au Probatoire et au Baccalauréat technique, la question « dégagez la thèse de l'auteur » ou « quelle est la position de l'auteur ? » revient presque chaque année dans l'épreuve de Français et vaut souvent 2 points. Une réponse exacte, formulée en une phrase complète, suffit : inutile de recopier tout le texte ! Entraîne-toi à repérer les marqueurs « il faut », « nous pensons que », « il est temps de ».
\end{savaistu}

\begin{exercice}[À toi de jouer]
Voici trois extraits. Pour chacun, indique s'il s'agit d'un fait ou d'une opinion, puis dis si le texte relève de l'argumentation directe ou indirecte.
\begin{enumerate}
\item « Le Cameroun compte dix régions. »
\item « Nous sommes convaincus que le sport scolaire doit être obligatoire : il discipline les élèves et améliore leur santé. »
\item « Le vieux baobab, voyant les jeunes arbres se plaindre du vent, leur dit : "Celui qui plie survit ; celui qui refuse de plier se brise." »
\end{enumerate}
\end{exercice}

\begin{corrige}[Exercice : À toi de jouer]
\begin{enumerate}
\item Fait (vérifiable, incontestable) ; ce n'est pas une argumentation, c'est une simple information.
\item Opinion défendue par deux arguments (« il discipline », « il améliore la santé ») : argumentation directe (marqueur « nous sommes convaincus que », les deux-points annonçant les raisons).
\item Opinion transmise par une parole de personnage dans un récit : argumentation indirecte (la sagesse du baobab fait passer le message « il faut savoir s'adapter »).
\end{enumerate}
\end{corrige}

\begin{aretenir}[L'essentiel de la section]
Un texte argumentatif défend ou rejette une thèse pour \emph{convaincre} (par la raison) et/ou \emph{persuader} (par l'émotion). Il s'organise autour de quatre éléments : thèse, arguments, exemples, conclusion. On le reconnaît à ses marques linguistiques (connecteurs logiques, modalisateurs, pronoms personnels). Il peut être direct (essai, éditorial) ou indirect (conte, fable, roman). Savoir distinguer fait et opinion, argument et exemple, est la première étape de toute contraction de texte réussie.
\end{aretenir}

\coursec{La structure argumentative et les techniques de réduction}

Dans la section précédente, nous avons découvert ce qui fait la spécificité du \cle{texte argumentatif} : une thèse défendue, des arguments, des exemples, un locuteur qui cherche à convaincre ou à persuader. Nous savons désormais \emph{reconnaître} un tel texte. Mais à l'examen du Probatoire et du Baccalauréat technique, on ne vous demandera pas seulement de le reconnaître : on vous demandera de le \cle{résumer}, c'est-à-dire d'en restituer les idées essentielles dans un nombre de mots imposé. Pour y parvenir, deux compétences sont indispensables : comprendre comment le texte est \emph{construit} (sa structure argumentative) et maîtriser les \emph{techniques de réduction} qui permettent de le raccourcir sans le trahir. C'est l'objet de cette section.

\courssub{La structure argumentative type d'un texte}

\textbf{Intuition.} Imaginez un vendeur du marché Mokolo qui veut vous convaincre d'acheter ses arachides. Il ne commence pas par crier ses prix : il attire d'abord votre attention (« Viens goûter, mon frère ! »), puis il avance ses raisons (« C'est frais, ça vient de Mbanga, le prix est doux »), et il conclut (« Tu prends combien de bouteilles ? »). Tout texte argumentatif fonctionne de la même manière : il \emph{accroche}, il \emph{argumente}, il \emph{conclut}.

\begin{definition}[La structure argumentative]
La \cle{structure argumentative} est l'organisation logique d'un texte qui défend une thèse. Elle comprend classiquement trois parties :
\begin{itemize}
\item une \cle{introduction} : elle présente le \textbf{sujet} et pose la \textbf{problématique} (la question à laquelle le texte va répondre) ;
\item un \cle{développement} : il expose les \textbf{arguments} (raisons qui soutiennent la thèse), hiérarchisés du plus faible au plus fort ou du plus simple au plus complexe, chacun étant généralement illustré par des \textbf{exemples} ;
\item une \cle{conclusion} : elle rappelle la thèse défendue et peut ouvrir sur une perspective nouvelle.
\end{itemize}
\end{definition}

\begin{aretenir}[Le squelette d'un texte argumentatif]
\begin{center}
\begin{tabularx}{\linewidth}{|l|X|X|}
\hline
\rowcolor{popblueL} \textbf{Partie} & \textbf{Fonction} & \textbf{Indices repérables} \\
\hline
Introduction & Présenter le sujet, poser le problème & Questions, « on peut se demander », « il convient d'examiner » \\
\hline
Développement & Enchaîner les arguments & « d'abord », « ensuite », « en outre », « de plus », « enfin » \\
\hline
Conclusion & Affirmer la thèse, ouvrir & « ainsi », « en définitive », « c'est pourquoi » \\
\hline
\end{tabularx}
\end{center}
\end{aretenir}

\courssub{Lire le plan d'un texte : articulations logiques et phrases charnières}

Pour résumer un texte, il faut d'abord en voir le \emph{plan}. Or le plan n'est presque jamais annoncé explicitement : il est signalé par des mots-outils.

\begin{definition}[Articulations logiques et phrases charnières]
Les \cle{articulations logiques} (ou connecteurs) sont les mots et expressions qui relient les idées entre elles : \emph{d'abord, ensuite, en effet, cependant, par conséquent, enfin}. Une \cle{phrase charnière} est une phrase qui annonce le passage d'une étape du raisonnement à une autre : « Mais ce premier argument ne suffit pas », « Venons-en maintenant à la question essentielle ».
\end{definition}

\begin{methode}[Repérer le plan d'un texte argumentatif]
\begin{enumerate}
\item \textbf{Première lecture} globale, sans stylo, pour saisir le sujet et la thèse générale.
\item \textbf{Deuxième lecture} active : souligner en rouge les connecteurs logiques, en bleu les phrases qui expriment la thèse et chaque argument.
\item \textbf{Numéroter} dans la marge les grandes étapes du raisonnement (1, 2, 3...).
\item \textbf{Rédiger} le schéma du texte en trois ou quatre lignes : « L'auteur affirme que... D'abord il montre que... Ensuite il prouve que... Il conclut que... ».
\end{enumerate}
\end{methode}

\begin{exemplebox}[Un plan repéré grâce aux connecteurs]
Texte : « \emph{D'abord}, l'école forme des citoyens capables de lire et d'écrire. \emph{Ensuite}, elle transmet les valeurs de la nation. \emph{Cependant}, certains prétendent qu'elle ne prépare pas à la vie pratique. \emph{En réalité}, sans instruction de base, aucun apprentissage technique n'est possible. \emph{C'est pourquoi} l'école demeure irremplaçable. »

Plan : thèse = l'école est irremplaçable ; argument 1 = elle instruit ; argument 2 = elle transmet des valeurs ; objection réfutée = le reproche d'impraticabilité ; conclusion = réaffirmation de la thèse.
\end{exemplebox}

\begin{attention}[Piège fréquent]
Ne confondez pas les \textbf{exemples} avec les \textbf{arguments}. L'exemple illustre, l'argument prouve. Dans un résumé, on garde les arguments et on supprime presque toujours les exemples. Souligner tout le texte revient à ne rien souligner : soyez sélectif.
\end{attention}

\courssub{La contraction de texte : définition et principe général}

\begin{definition}[La contraction de texte]
La \cle{contraction de texte} (ou résumé) est la réduction d'un texte à une \textbf{fraction donnée} de sa longueur initiale — souvent le \textbf{quart} — sans en déformer le sens, sans ajouter d'idée personnelle et en respectant la logique du texte d'origine. Un texte de 300 mots résumé au quart donne un texte d'environ 75 mots.
\end{definition}

L'image la plus parlante est celle de l'\textbf{entonnoir} : le texte entre large en haut, et il ressort étroit en bas, mais c'est toujours la même substance qui coule, simplement débarrassée de ce qui n'était pas essentiel.

\begin{popfigure}
\begin{center}
\begin{tikzpicture}
% Entonnoir
\draw[fill=popblueL, draw=popblue] (0,4) -- (6,4) -- (4.2,2.4) -- (1.8,2.4) -- cycle;
\draw[fill=poporangeL, draw=poporange] (1.8,2.4) -- (4.2,2.4) -- (3.5,1.2) -- (2.5,1.2) -- cycle;
\draw[fill=popgreenL, draw=popgreen] (2.5,1.2) -- (3.5,1.2) -- (3.1,0) -- (2.9,0) -- cycle;
% Étiquettes
\node[align=center] at (3,3.2) {\textbf{Texte de départ}\\ 300 mots};
\node[align=center, font=\small] at (3,1.8) {Sélection / Généralisation\\ $\sim$ 150 mots};
\node[align=center, font=\small] at (3,0.6) {\footnotesize Fusion / Reformulation};
\node[below, font=\small] at (3,0) {\textbf{Résumé : 75 mots}};
% Flèches latérales
\draw[-{Stealth}, thick, popdark] (6.4,3.2) -- (6.4,1.8) node[midway, right, font=\small, align=left] {suppression\\ sélection\\ généralisation};
\draw[-{Stealth}, thick, popdark] (6.4,1.6) -- (6.4,0.4) node[midway, right, font=\small, align=left] {fusion\\ reformulation};
\end{tikzpicture}
\end{center}
\legende{Schéma en entonnoir de la contraction : un texte de 300 mots passe par la sélection des idées essentielles (150 mots), puis par la reformulation, pour aboutir à un résumé d'environ 75 mots, soit le quart du texte initial.}
\end{popfigure}

\begin{methode}[Les quatre étapes du résumé]
\begin{enumerate}
\item \textbf{Lire deux fois} le texte : une fois pour comprendre, une fois pour analyser.
\item \textbf{Souligner les idées essentielles} : la thèse, les arguments, la conclusion (et non les exemples).
\item \textbf{Réduire} en appliquant les quatre techniques : suppression, sélection, généralisation, fusion.
\item \textbf{Compter les mots} et vérifier que le total obtenu respecte la consigne (tolérance de $\pm$ 10 \%).
\end{enumerate}
\end{methode}

\courssub{Les quatre techniques de réduction}

\textbf{Technique n° 1 : la suppression.} C'est la technique la plus simple : on \emph{élimine} purement et simplement ce qui n'est pas indispensable à la compréhension de la pensée de l'auteur. On supprime :
\begin{itemize}
\item les \textbf{répétitions} (la même idée dite deux fois avec des mots différents) ;
\item les \textbf{digressions} (les passages qui s'écartent du sujet) ;
\item les \textbf{détails} et précisions secondaires (dates, chiffres, descriptions pittoresques) ;
\item les \textbf{énumérations} longues ;
\item les \textbf{exemples} qui ne font qu'illustrer un argument déjà exprimé.
\end{itemize}

\textbf{Technique n° 2 : la sélection.} C'est l'inverse de la suppression : au lieu de regarder ce qu'on enlève, on regarde ce qu'on \emph{garde}. On conserve les \cle{mots-clés} (les noms et verbes porteurs de sens), les phrases qui expriment la \textbf{thèse} et les \textbf{arguments}, ainsi que les connecteurs indispensables à la logique du texte. Tout le reste peut disparaître.

\textbf{Technique n° 3 : la généralisation.} On remplace une énumération ou un terme précis par un \cle{terme générique}. Ainsi « mangues, goyaves, avocats et papayes » devient simplement « fruits » ; « marteaux, tournevis, pinces et scies » devient « outils ». La généralisation fait gagner beaucoup de mots sans rien perdre de l'idée.

\textbf{Technique n° 4 : la fusion.} On \emph{condense} plusieurs phrases en une seule, en repérant l'idée commune. « Le taxi roulait vite. Il slalomait entre les motos. Il klaxonnait sans cesse. » devient : « Le taxi roulait vite et dangereusement. » Trois phrases, une seule : l'économie est considérable.

\begin{popfigure}
\begin{center}
\begin{tikzpicture}
\node[font=\bfseries\large, popdark] at (5,5.6) {Les quatre techniques de réduction};
% Ligne 1
\node[draw=popblue, fill=popblueL, rounded corners, minimum width=2.6cm, minimum height=1cm, font=\bfseries] at (1.3,4.5) {Suppression};
\node[draw=popline, fill=white, rounded corners, minimum width=6.4cm, minimum height=1cm, align=left, font=\small, text width=6cm] at (6.3,4.5) {« Le jeune élève, âgé de dix-sept ans, travaille avec sérieux. » $\rightarrow$ « L'élève travaille. »};
% Ligne 2
\node[draw=poporange, fill=poporangeL, rounded corners, minimum width=2.6cm, minimum height=1cm, font=\bfseries] at (1.3,3.2) {Sélection};
\node[draw=popline, fill=white, rounded corners, minimum width=6.4cm, minimum height=1cm, align=left, font=\small, text width=6cm] at (6.3,3.2) {On garde : thèse + arguments + mots-clés ;\\ on abandonne : exemples, détails, citations.};
% Ligne 3
\node[draw=popgreen, fill=popgreenL, rounded corners, minimum width=2.6cm, minimum height=1cm, font=\bfseries] at (1.3,1.9) {Généralisation};
\node[draw=popline, fill=white, rounded corners, minimum width=6.4cm, minimum height=1cm, align=left, font=\small, text width=6cm] at (6.3,1.9) {« mangues, goyaves, avocats, papayes »\\ $\rightarrow$ « fruits »};
% Ligne 4
\node[draw=poppurple, fill=poppurpleL, rounded corners, minimum width=2.6cm, minimum height=1cm, font=\bfseries] at (1.3,0.6) {Fusion};
\node[draw=popline, fill=white, rounded corners, minimum width=6.4cm, minimum height=1cm, align=left, font=\small, text width=6cm] at (6.3,0.6) {« Il pleuvait. La route était glissante. »\\ $\rightarrow$ « La pluie rendait la route glissante. »};
\end{tikzpicture}
\end{center}
\legende{Tableau synoptique des quatre techniques de réduction, avec pour chacune un exemple avant / après.}
\end{popfigure}

\begin{attention}[La règle d'or du résumé]
Ne jamais ajouter d'idée personnelle ni de jugement dans un résumé. Le résumé n'est pas un commentaire : vous ne dites pas si l'auteur a raison ou tort, vous restituez \emph{sa} pensée, fidèlement et objectivement. Toute phrase commençant par « je pense que », « heureusement », « malheureusement » est une faute dans un résumé. De même, ne copiez pas de longues phrases du texte : reformulez avec vos propres mots.
\end{attention}

\begin{savaistu}[Le nombre de mots, un critère de notation]
À l'examen, le respect du nombre de mots imposé est un critère de notation \emph{explicite}. Une tolérance de $\pm$ 10 \% est généralement admise : pour un résumé demandé en 100 mots, tout texte compris entre 90 et 110 mots est accepté ; en dessous ou au dessus, des points sont retirés. Indiquez toujours le nombre de mots à la fin de votre copie, entre parenthèses : cela prouve au correcteur que vous avez vérifié.
\end{savaistu}

\courssub{Application guidée : réduire un paragraphe de moitié}

Appliquons maintenant les quatre techniques à un paragraphe de 120 mots sur les marchés de Mokolo à Yaoundé, que nous devons réduire de moitié (60 mots), puis au quart (30 mots).

\begin{exoresolu}[Réduire un paragraphe de 120 mots à 30 mots]
\textbf{Énoncé.} Réduisez le texte suivant d'abord de moitié (environ 60 mots), puis au quart (environ 30 mots).

« Le marché de Mokolo, situé au cœur de la ville de Yaoundé, capitale politique du Cameroun, est l'un des plus grands marchés de toute la sous-région d'Afrique centrale. Chaque matin, dès cinq heures, avant même le lever du soleil, des centaines de commerçants, hommes et femmes de tous âges, arrivent de tous les quartiers de la ville et même des villages environnants pour installer leurs étals. On y trouve des mangues, des goyaves, des avocats, des papayes, des bananes plantains, mais aussi des tissus colorés, des chaussures, des casseroles et bien d'autres marchandises. Malheureusement, ce marché souffre d'un encombrement permanent : les allées sont étroites, boueuses en saison des pluies, et les pickpockets y font parfois régner l'insécurité. Il faut donc moderniser ce marché, car il reste le poumon économique de la ville. » (120 mots)
\tcblower
\textbf{Solution détaillée.}

\textbf{Étape 1 : lecture et repérage.} Thèse : le marché de Mokolo doit être modernisé car il est essentiel à l'économie de Yaoundé. Arguments : 1) c'est un grand marché très actif ; 2) il souffre de l'encombrement et de l'insécurité.

\textbf{Étape 2 : suppression.} On élimine les détails (« capitale politique du Cameroun », « avant même le lever du soleil », « hommes et femmes de tous âges »), l'exemple des pickpockets et le jugement « malheureusement ».

\textbf{Étape 3 : généralisation.} « mangues, goyaves, avocats, papayes, bananes plantains » $\rightarrow$ « fruits » ; « tissus, chaussures, casseroles » $\rightarrow$ « marchandises diverses ».

\textbf{Étape 4 : fusion.} « Les allées sont étroites. Elles sont boueuses en saison des pluies. » $\rightarrow$ « les allées sont étroites et boueuses ».

\textbf{Résumé à 60 mots :} « Le marché de Mokolo, au cœur de Yaoundé, est l'un des plus grands d'Afrique centrale. Chaque matin, des centaines de commerçants y installent leurs étals de fruits et de marchandises diverses. Mais ce marché souffre d'un encombrement permanent : les allées sont étroites et boueuses, et l'insécurité y règne parfois. Il faut donc le moderniser, car il reste le poumon économique de la ville. » (60 mots)

\textbf{Résumé à 30 mots :} « Grand marché de Yaoundé, Mokolo attire chaque jour des centaines de commerçants. Mais encombrement et insécurité y règnent : il faut le moderniser, car il demeure le poumon économique de la ville. » (30 mots)

\textbf{Vérification :} la thèse (moderniser Mokolo) et les deux arguments (importance économique, encombrement) sont conservés ; aucune idée personnelle n'a été ajoutée ; le nombre de mots respecte la consigne.
\end{exoresolu}

\begin{exemplebox}[Un autre passage chiffré : la circulation à Douala]
Texte de départ (120 mots) : « À Douala, capitale économique du Cameroun, la circulation est devenue un véritable cauchemar pour les habitants. Chaque matin, entre six heures et neuf heures, des embouteillages monstres paralysent les grands axes : le pont de la Wouri, le rond-point Deïdo, l'axe Akwa-Bonabéri. Les motos-taxis, les taxis jaunes, les camions et les bus se disputent la chaussée dans un concert assourdissant de klaxons. Les causes de cette situation sont connues : l'augmentation rapide du nombre de véhicules, le mauvais état de certaines routes et l'indiscipline des conducteurs. Les conséquences sont lourdes : retards au travail, fatigue, pollution de l'air. Il est urgent d'améliorer les transports en commun. »

Résumé au quart (30 mots) : « À Douala, les embouteillages quotidiens, dus à la multiplication des véhicules, au mauvais état des routes et à l'indiscipline, provoquent retards et pollution : il est urgent de développer les transports en commun. » (30 mots)
\end{exemplebox}

\begin{exercice}[À vous de jouer]
Réduisez de moitié (environ 50 mots) le paragraphe suivant, en appliquant les quatre techniques et en indiquant le nombre de mots obtenu :

« Le sport scolaire, qu'il s'agisse du football, du handball, de l'athlétisme ou de la course de fond, joue un rôle capital dans la vie des élèves des lycées techniques. D'abord, il permet aux jeunes de se maintenir en excellente santé, car courir, sauter et lancer développent les muscles et le souffle. Ensuite, il apprend le respect des règles et de l'arbitre, qualité indispensable dans la vie professionnelle. De plus, lors des compétitions inter-établissements, il crée entre les élèves une amitié solide qui dure parfois toute la vie. C'est pourquoi les pouvoirs publics devraient doter tous les lycées de terrains et d'équipements sportifs modernes. » (100 mots)
\end{exercice}

\begin{corrige}[Exercice : le sport scolaire]
\textbf{Étapes :} suppression des exemples (football, handball, athlétisme, course de fond $\rightarrow$ généralisation : « le sport » ; courir, sauter, lancer $\rightarrow$ « l'effort physique ») ; conservation de la thèse (doter les lycées d'équipements sportifs) et des trois arguments (santé, respect des règles, amitié).

\textbf{Proposition de résumé (50 mots) :} « Le sport scolaire joue un rôle capital pour les élèves : il assure leur santé physique, leur enseigne le respect des règles et forge des amitiés solides. C'est pourquoi l'État doit doter tous les lycées techniques d'installations sportives modernes et adaptées. » (48 mots — toute formulation voisine, comprise entre 45 et 55 mots et conservant la thèse et les trois arguments, est acceptable.)
\end{corrige}

\begin{aretenir}[L'essentiel de la section]
\begin{itemize}
\item Un texte argumentatif suit une structure en trois temps : \textbf{introduction} (sujet + problématique), \textbf{développement} (arguments hiérarchisés), \textbf{conclusion} (thèse réaffirmée).
\item Le plan se repère grâce aux \textbf{articulations logiques} et aux \textbf{phrases charnières}.
\item Quatre techniques permettent de réduire : \textbf{suppression}, \textbf{sélection}, \textbf{généralisation}, \textbf{fusion}.
\item Le résumé restitue fidèlement la pensée de l'auteur, sans idée personnelle, dans le nombre de mots imposé ($\pm$ 10 \%).
\end{itemize}
\end{aretenir}

Maintenant que nous savons découper un texte argumentatif et le réduire méthodiquement, nous allons mettre ces compétences à l'épreuve d'un texte littéraire complet : la lecture méthodique du \emph{Lion et la perle} de Wole Soyinka, qui fera l'objet de la section suivante.

\coursec{Lecture méthodique du Lion et la perle (n° 4 et 5)}

Après avoir étudié la structure argumentative et les techniques de réduction qui permettent de condenser un texte sans en trahir la pensée, nous appliquons désormais ces outils à l'analyse fine d'une œuvre au programme : \emph{Le Lion et la perle} de Wole Soyinka. La lecture méthodique, telle que le programme la conçoit, n'est pas une simple explication de texte linéaire ; c'est une enquête rigoureuse qui, à partir d'extraits choisis, met au jour les procédés d'écriture, les stratégies argumentatives des personnages et les enjeux idéologiques de l'œuvre. C'est précisément ce travail de repérage et d'interprétation qui alimente, ensuite, la contraction de texte : on ne peut réduire que ce qu'on a d'abord compris en profondeur.

\courssub{Rappel de l'œuvre : tradition, modernité et le triangle amoureux}

\begin{definition}[Le Lion et la perle — fiche d'identité]
\emph{Le Lion et la perle} (\emph{The Lion and the Jewel}), publié en 1963, est une comédie de l'écrivain nigérian Wole Soyinka (prix Nobel de littérature 1986). La pièce se déroule à Ilujinle, village yoruba fictif, et met en scène un conflit générationnel et culturel à travers trois figures archétypales :
\begin{itemize}
    \item \textbf{Baroka (le Lion) :} le Bale (chef traditionnel), rusé, polygame, gardien des coutumes, mais fin stratège qui sait utiliser la modernité à son profit.
    \item \textbf{Lakunle (l'instituteur) :} jeune homme « moderne », occidentalisé, qui rejette les coutumes (dot, polygamie) au nom du « progrès » et de l'égalité, mais dont le discours est souvent pédant, maladapté et stérile.
    \item \textbf{Sidi (la Perle) :} la belle villageoise, objet du désir des deux hommes, symbole de la jeunesse africaine tiraillée entre l'attrait du moderne (photographie, célébrité) et la sécurité valorisante de la tradition.
\end{itemize}
\emph{Sadiku}, l'épouse senior de Baroka, joue le rôle d'intermédiaire et de confidente, manipulée par son mari mais aussi actrice de la ruse finale.
\end{definition}

\begin{popfigure}
\begin{tikzpicture}[
    node distance=2.5cm and 3cm,
    pers/.style={circle, draw=popdark, thick, fill=popcream, minimum width=2.2cm, align=center, font=\small},
    rel/.style={->, >=Stealth, thick, popdark, font=\footnotesize, sloped},
    bidir/.style={<->, >=Stealth, thick, poporange, font=\footnotesize, sloped}
]
% Nœuds
\node[pers, align=center] (baroka){Baroka\\ (Le Lion)\\\textit{Bale, Tradition, Ruse}};
\node[pers, right=of baroka, align=center] (sidi){Sidi\\ (La Perle)\\\textit{Jeunesse, Beauté,}\\\textit{ Enjeu du conflit}};
\node[pers, below=of sidi, align=center] (lakunle){Lakunle\\ (L'instituteur)\\\textit{Modernité, Progrès,}\\\textit{ Discours livresque}};
\node[pers, left=of lakunle, align=center] (sadiku){Sadiku\\ (L'épouse senior)\\\textit{Intermédiaire,}\\\textit{ Manipulatrice / Manipulée}};
% Relations
\draw[bidir] (baroka) -- node[above] {Désir / Conquête} (sidi);
\draw[bidir] (lakunle) -- node[right, align=center]{Amour / Mariage\\ \textit{sans dot}} (sidi);
\draw[rel, bend left=20] (baroka) to node[above] {Ordonne / Manipule} (sadiku);
\draw[rel, bend left=20] (sadiku) to node[below] {Rapporte / Conseille} (baroka);
\draw[rel, dashed, poppurple] (sadiku) to node[left] {Persuade / Trahit} (sidi);
\draw[rel, dashed, poppurple] (lakunle) to node[below] {Méprise / Ignore} (sadiku);
% Légende
\node[draw=popline, fill=popcream, rounded corners, below left=-0.5cm and -0.5cm of lakunle.south west, anchor=north west, font=\footnotesize, text width=5cm] {
    \textbf{Légende :}\\
    \textcolor{poporange}{\textbf{---}} Relation centrale (amour/rivalité)\\
    \textbf{---} Relation de pouvoir / manipulation\\
    \textcolor{poppurple}{\textbf{--- ---}} Relation d'influence / communication
};
\end{tikzpicture}
\legende{Schéma des relations entre les personnages principaux : le triangle Baroka–Sidi–Lakunle médiatisé par Sadiku.}
\end{popfigure}

\courssub{Lecture méthodique n° 4 : le stratagème de Baroka — la fausse impuissance (Acte 2, scène « Le Bureau de Baroka »)}

Cette lecture porte sur l'extrait où Baroka reçoit Sidi dans son bureau, après que Sadiku a annoncé la « nouvelle » de son impuissance. C'est le cœur dramaturgique de la pièce : la victoire de la ruse traditionnelle sur la naïveté moderniste.

\begin{methode}[La lecture méthodique en trois temps]
Pour analyser cet extrait (ou tout extrait au programme), suivez cette démarche :
\begin{enumerate}
    \item \textbf{Lecture 1 — Repérage global :} Qui parle ? À qui ? Où ? Quand ? Quel est le sujet apparent ? Quel ton domine (ironique, solennel, intime) ?
    \item \textbf{Lecture 2 — Analyse des procédés :} Repérez les figures de style, le lexique (champs lexicaux), la syntaxe (phrases courtes/longues, interrogatives, impératives), les registres de langue (soutenu, familier, proverbes), la structure argumentative (thèses, arguments, exemples).
    \item \textbf{Lecture 3 — Interprétation et synthèse :} Quel est le projet énonciatif du locuteur ? Comment la forme sert le fond ? Quels enjeux idéologiques, culturels, dramatiques ?
\end{enumerate}
\end{methode}

\begin{exemplebox}[Analyse du discours de Baroka sur le timbre et la machine — Argumentation par l'émotion et le symbole]
\textbf{Extrait (résumé) :} Baroka montre à Sidi le timbre-poste à son effigie qu'il a fait fabriquer, puis évoque la machine à timbrer qu'il a fait saboter pour retarder la modernité.

\emph{Repérage des procédés (Lecture 2) :}
\begin{itemize}
    \item \textbf{Champ lexical de la tradition et de l'ancestralité :} « ancêtres », « racines », « sang », « terre », « coutume ».
    \item \textbf{Champ lexical de la modernité dépréciée :} « machine », « fer », « froid », « bruit », « étranger », « papier ».
    \item \textbf{La métaphore filée du timbre :} Le timbre à l'effigie de Baroka symbolise l'appropriation de l'outil moderne (la poste, l'image reproductible) par le pouvoir traditionnel. Baroka \emph{est} l'image, il \emph{contrôle} la diffusion.
    \item \textbf{L'argumentation par l'absurde et l'ironie :} « J'ai payé pour que la machine ne marche pas. » C'est un argument \emph{a fortiori} : sa puissance se mesure à sa capacité à arrêter le progrès.
    \item \textbf{L'adaptation à l'auditoire (Sidi) :} Il utilise un langage imagé, sensuel, flatteur (« la perle », « mes yeux »), loin du jargon pédagogique de Lakunle. Il séduit par l'intelligence du cœur, pas par la leçon de morale.
\end{itemize}

\emph{Interprétation (Lecture 3) :}
Baroka ne refuse pas la modernité ; il la \emph{domestique}. Son argumentaire vise à prouver que la tradition n'est pas figée : elle est vivante, assez forte pour avaler le moderne et le rendre insignifiant (la machine cassée). Il persuade Sidi non par la logique formelle, mais par la \emph{preuve par l'acte} (le timbre existe, la machine est arrêtée) et la \emph{séduction esthétique}. C'est une argumentation \emph{ethos}-centrée (crédibilité du chef, prestige) et \emph{pathos}-centrée (flatterie, beauté), là où Lakunle ne mobilise que le \emph{logos} (raison abstraite, droits de la femme).
\end{exemplebox}

\begin{attention}[Piège fréquent : confondre résumé d'intrigue et analyse argumentative]
\emph{Ne pas écrire :} « Baroka ment à Sadiku en disant qu'il est impuissant pour que Sidi vienne et il la séduit. »
\emph{Mais écrire :} « Baroka construit une \textbf{stratégie communicationnelle} en trois temps : 1) La désinformation contrôlée (le secret confié à Sadiku, garantie de diffusion). 2) La mise en scène de soi (le bureau, le lutteur, le timbre) qui établit son \emph{ethos} de puissance virile et intellectuelle. 3) L'argumentation \emph{ad hominem} envers Sidi : il valide sa beauté (reconnaissance de sa valeur marchande symbolique) et dévalorise Lakunle par l'ironie ("le livreur de lettres"), retournant la modernité (la photo, le timbre) en instrument de sa propre gloire. »
\end{attention}

\courssub{Repérage de l'argumentation des personnages : Baroka vs Lakunle}

Pour préparer la contraction, il faut isoler les \cle{nœuds argumentatifs} de chaque personnage. Le tableau ci-dessous synthétise les thèses, arguments et procédés dominants.

\begin{center}
\begin{tabularx}{\linewidth}{|l|X|X|}\hline
\rowcolor{popblueL}
\textbf{Personnage} & \textbf{Baroka (Tradition vivante / Ruse)} & \textbf{Lakunle (Modernité livresque / Rigidité)} \\ \hline
\textbf{Thèse centrale} & La tradition est dynamique ; le chef moderne est celui qui maîtrise les deux mondes. & La tradition est rétrograde ; le progrès impose l'égalité et le refus des coutumes « barbares ». \\ \hline
\textbf{Arguments types} & \begin{itemize}\item Pragmatisme : « Je paye pour arrêter la machine. »\item Symbolique : Le timbre, la photo = appropriation.\item Séduction : Valorisation de la femme (Sidi) comme reine, pas comme égale abstraite.\item Ironie : Moquerie douce de Lakunle (« l'homme qui veut payer la dot avec des mots »).\end{itemize} & \begin{itemize}\item Principes abstraits : « Dot = achat d'esclave », « Monogamie = civilisation ».\item Références livresques : Bible, dictionnaire, « livres d'école ».\item Culpabilisation : Sidi est « ignorante », « païenne ».\item Refus du compromis : Rejet total de la coutume (pas de dot = pas de mariage).\end{itemize} \\ \hline
\textbf{Procédés rhétoriques} & \begin{itemize}\item \textbf{Métaphore / Image} (timbre, lion, perle).\item \textbf{Proverbes yoruba} (autorité ancestrale).\item \textbf{Parole performative} (il agit en parlant).\item \textbf{Registre oral / familier / sensuel}.\end{itemize} & \begin{itemize}\item \textbf{Logos formel} (syllogismes maladroits).\item \textbf{Néologismes / Anglicismes} (« progress », « machine »).\item \textbf{Ton professoral, donneur de leçons}.\item \textbf{Registre soutenu, déconnecté du réel villageois}.\end{itemize} \\ \hline
\textbf{Efficacité dramatique} & \textbf{Gagnante :} Il obtient Sidi, le village l'acclame, il intègre la photo dans son palais. & \textbf{Perdante :} Il reste seul, ridicule, incapable d'agir concrètement (ne paie pas la dot, ne « modernise » rien). \\ \hline
\end{tabularx}
\end{center}

\courssub{Les tonalités : comique, dramatique, satirique}

L'étude du \emph{Lion et la perle} offre un terrain d'application privilégié pour l'analyse des \cle{tonalités}, notion de stylistique rhétorique au programme. La tonalité est l'attitude affective et intellectuelle de l'énonciateur (ou de l'auteur à travers ses personnages) vis-à-vis du sujet traité ; elle se manifeste par le choix lexical, la syntaxe, les figures et les registres.

\begin{definition}[Les trois tonalités majeures de la pièce]
\begin{itemize}
    \item \textbf{La tonalité comique :} Elle naît du décalage, de la disproportion ou de l'incongruité. Elle vise à faire rire pour mieux critiquer.
        \begin{itemize}
            \item \emph{Comique de situation :} Lakunle qui refuse de porter le paquet de Sidi « parce que c'est le travail des femmes » puis se plaint de la fatigue ; la course-poursuite du photographe.
            \item \emph{Comique de caractère :} Lakunle, « homme moderne » qui cite le dictionnaire pour draguer, mais ne sait pas danser ni payer la dot. Baroka, « vieux lion » qui fait de la musculation avec un lutteur pour prouver sa virilité.
            \item \emph{Comique de mots :} Les néologismes de Lakunle (« progress », « machine », « civilised ») et ses métaphores maladroites (« la femme est le sel de la cuisine »).
        \end{itemize}
    \item \textbf{La tonalité dramatique :} Elle exprime la tension, le conflit, l'angoisse ou la gravité d'un choix irréversible.
        \begin{itemize}
            \item Le face-à-face Sidi / Baroka au bureau : enjeu de la perte de virginité, du pouvoir, de l'identité.
            \item Le retour de Sidi (Acte 3) : l'annonce « La dot est payée » scelle un destin. La violence symbolique de l'injure de Lakunle (« vieux singe ») marque la rupture définitive.
        \end{itemize}
    \item \textbf{La tonalité satirique :} Elle utilise l'ironie, la caricature et le ridicule pour dénoncer des travers sociaux, politiques ou culturels. C'est la tonalité \emph{auctoriale} dominante chez Soyinka.
        \begin{itemize}
            \item \emph{Satire de la modernité superficielle :} Lakunle incarne une modernité « de surface » (vocabulaire, vêtements) qui ne change pas les rapports de force réels.
            \item \emph{Satire du pouvoir traditionnel :} Baroka n'est pas un saint ; sa ruse, sa polygamie, son autoritarisme sont montrés sans complaisance (le sabotage de la machine, la manipulation de Sadiku).
            \item \emph{Satire de la condition féminine :} Sidi est objet de convoitise (photo, timbre, dot), mais c'est elle qui a le dernier mot (le choix, la danse). L'auteur satirise l'aliénation \emph{et} l'instrumentalisation.
        \end{itemize}
\end{itemize}
\end{definition}

\begin{exemplebox}[Identification de la tonalité dans un extrait]
\textbf{Extrait :} Lakunle : « Je ne paierai pas la dot. C'est une coutume barbare, dégradante, humiliante. / Sidi : Alors tu ne m'épouseras pas. / Lakunle : Nous nous marierons devant le tribunal. Ce sera un mariage chrétien, monogame, civilisé. »

\begin{itemize}
    \item \textbf{Tonalité comique :} Le contraste entre le vocabulaire juridique abstrait (« tribunal », « monogame », « civilisé ») et la réalité villageoise (Sidi ne comprend pas, le public rit de l'inadéquation).
    \item \textbf{Tonalité satirique :} La caricature du « jeune instruit » qui pense changer une société par un décret lexical. L'ironie dramatique : le spectateur sait que ce projet échouera.
    \item \textbf{Tonalité dramatique (latente) :} Le refus de la dot menace l'honneur de Sidi et de sa famille ; le rire bascule en tension.
\end{itemize}
\end{exemplebox}

\courssub{Lecture méthodique n° 5 : le dénouement et le choix de Sidi (Acte 3, scène finale)}

La dernière lecture méthodique porte sur la scène du mariage et la danse finale. Sidi a choisi Baroka. Lakunle, d'abord révolté, finit par danser avec une jeune fille, suggérant l'acceptation (ou la résignation) du statu quo.

\begin{exoresolu}[Dégagement des idées essentielles de l'extrait final pour contraction]
\textbf{Énoncé :} À partir de l'extrait suivant (fin de la pièce), dégagez les 3 idées essentielles à retenir pour une contraction de texte.

\emph{Extrait (synthèse) :} Sidi revient du palais de Baroka, triomphante. Elle annonce qu'elle a vu la machine à timbrer (elle ne marchait pas). Elle refuse Lakunle (« Je ne veux pas de lui »). Elle court préparer le festin de noce. Lakunle, déçu, se tourne vers une jeune fille qui passe et l'invite à danser.

\tcblower
\textbf{Solution détaillée — Démarche de sélection :}

\textbf{Étape 1 : Segmenter le texte en séquences logiques.}
\begin{enumerate}
    \item Retour de Sidi / Annonce de la supercherie (machine arrêtée).
    \item Rejet explicite de Lakunle / Choix de Baroka.
    \item Préparation des noces / Danse finale (communauté réunie).
    \item Attitude de Lakunle (repli / adaptation).
\end{enumerate}

\textbf{Étape 2 : Identifier l'idée-force de chaque séquence (noyau dur).}
\begin{itemize}
    \item Séquence 1 : \textbf{La ruse de Baroka est confirmée} ; la modernité technique est neutralisée par le pouvoir traditionnel.
    \item Séquence 2 : \textbf{Sidi opte pour la tradition valorisante} (reine, dot payée) contre la modernité dévalorisante (égalité abstraite, pas de dot).
    \item Séquence 3 : \textbf{La communauté valide l'ordre établi} par la fête rituelle ; la tradition sort renforcée.
    \item Séquence 4 : \textbf{Lakunle, marginalisé, s'adapte} (danse) sans avoir changé le système.
\end{itemize}

\textbf{Étape 3 : Fusionner et hiérarchiser pour la contraction (3 idées max).}
\begin{enumerate}
    \item \textbf{La victoire de la ruse traditionnelle :} Baroka, par une feinte d'impuissance, neutralise la modernité technique (machine) et conquiert Sidi.
    \item \textbf{Le choix lucide de Sidi :} Elle préfère le statut de « dernière épouse » d'un chef puissant (reconnaissance sociale, dot) à l'union « moderne » sans dot proposée par Lakunle, perçue comme une déchéance.
    \item \textbf{La restauration de l'ordre communautaire :} La fête finale scelle l'harmonie retrouvée ; Lakunle lui-même y participe, signe que le « progrès » n'a pas de prise sur le réel villageois.
\end{enumerate}

\textbf{Conseil de rédaction :} Dans la contraction, ces trois idées deviennent les trois paragraphes (ou les trois phrases complexes) du résumé. On élimine : les détails de la danse, les noms des musiciens, les répliques anecdotiques de Sadiku, les descriptions du festin.
\end{exoresolu}

\courssub{Dégagement des idées essentielles : de l'analyse à la contraction}

Le passage de la lecture méthodique à la contraction repose sur une opération intellectuelle précise : la \cle{hiérarchisation de l'information}. Tout ce qui est repéré en lecture 2 (procédés, images, arguments) sert à \emph{comprendre} ; seul le résultat de la lecture 3 (thèses, enjeux, issue dramatique) sert à \emph{écrire} la contraction.

\begin{propriete}[Règle de transposition : Analyse $\rightarrow$ Contraction]
\begin{itemize}
    \item \textbf{Les procédés stylistiques} (métaphores, ironie, proverbes) $\rightarrow$ deviennent des \textbf{adverbes ou adjectifs qualificatifs} dans le résumé (ex: « \emph{ironiquement} », « \emph{habilement} », « \emph{par une ruse} »).
    \item \textbf{Les arguments des personnages} $\rightarrow$ deviennent le \textbf{contenu factuel} du résumé (ex: « Baroka \emph{fait croire à son impuissance} »).
    \item \textbf{La structure dramatique} (début, péripétie, dénouement) $\rightarrow$ dicte le \textbf{plan du résumé} (chronologie respectée ou ordre logique).
    \item \textbf{Le point de vue du narrateur / auteur} (ici, l'ironie de Soyinka sur Lakunle) $\rightarrow$ doit \textbf{transparaître sans être commenté} (on écrit « Lakunle \emph{échoue} » et non « Soyinka \emph{montre que} Lakunle échoue »).
\end{itemize}
\end{propriete}

\courssub{Paronymie et antonymie : précision lexicale au service de l'argumentation}

La maîtrise du lexique est indispensable pour reformuler sans trahir la pensée de l'auteur. Deux phénomènes de sémantique lexicale sont au programme : la \cle{paronymie} et l'\cle{antonymie}.

\begin{definition}[Paronymie et antonymie]
\begin{itemize}
    \item \textbf{Paronymie :} Relation entre deux mots (paronymes) qui présentent une grande ressemblance formelle (phonétique ou graphique) mais des sens différents. La confusion est une faute de langue fréquente.
    \item \textbf{Antonymie :} Relation d'opposition de sens entre deux mots (antonymes). On distingue :
        \begin{itemize}
            \item \emph{Antonymes complémentaires (contradictoires) :} L'affirmation de l'un implique la négation de l'autre, sans terme moyen (ex: \emph{vivant / mort}, \emph{vrai / faux}, \emph{tradition / modernité} dans une vision binaire).
            \item \emph{Antonymes gradués (contraires) :} Opposition sur une échelle avec des degrés intermédiaires (ex: \emph{chaud / froid} $\rightarrow$ tiède ; \emph{riche / pauvre} $\rightarrow$ aisé ; \emph{rusé / naïf} $\rightarrow$ prudent).
            \item \emph{Antonymes réciproques (corrélatifs) :} L'un implique l'autre dans une relation inverse (ex: \emph{acheter / vendre}, \emph{donner / recevoir}, \emph{manipuler / être manipulé}).
        \end{itemize}
\end{itemize}
\end{definition}

\begin{exemplebox}[Paronymes et antonymes dans Le Lion et la perle]
\begin{center}
\begin{tabularx}{\linewidth}{|l|X|X|}\hline
\rowcolor{popblueL}
\textbf{Notion} & \textbf{Exemples (Paire)} & \textbf{Utilisation / Distinction} \\ \hline
\textbf{Paronymie} & \emph{Inférer / Infirmer} & Lakunle \emph{infère} (déduit) que Sidi l'aime ; Baroka \emph{infirme} (contredit/démontre le faux) la rumeur d'impuissance. \\
& \emph{Prospectif / Perspectif} & Le projet de Lakunle est \emph{prospectif} (tourné vers l'avenir) ; la vision de Baroka est \emph{perspective} (qui a du relief, de la profondeur) ou \emph{perspectiviste}. \\
& \emph{Éminent / Imminent} & Baroka est un chef \emph{éminent} (illustre) ; le danger n'est pas \emph{imminent} (proche) pour lui. \\
& \emph{Concret / Concomitant} & La dot est un acte \emph{concret} (réel) ; la danse est \emph{concomitante} (simultanée) au mariage. \\ \hline
\textbf{Antonymie (Graduée)} & \emph{Tradition / Modernité} & Non pas exclusives (Baroka les combine), mais en tension. Degrés : \emph{archaïsme — tradition vivante — modernité adaptée — modernisme aveugle}. \\
& \emph{Ruse / Force} & Baroka use de \emph{ruse} (intelligence détournée) ; Lakunle croit en la \emph{force} de la loi/raison. Degré intermédiaire : \emph{stratégie}. \\
& \emph{Parole / Acte} & Lakunle : parole stérile. Baroka : parole performative (acte). Réciproques : \emph{Promettre / Tenir}. \\ \hline
\textbf{Antonymie (Complémentaire)} & \emph{Marié / Célibataire} & Sidi bascule de l'un à l'autre sans état intermédiaire (dans la coutume). \\
& \emph{Vierge / Non-vierge} & Enjeu dramatique majeur : statut binaire aux conséquences sociales irréversibles. \\ \hline
\end{tabularx}
\end{center}
\end{exemplebox}

\begin{attention}[Piège : confusion paronymique en contraction]
Dans une contraction, chaque mot compte. Confondre \emph{inférer} et \emph{infirmer}, \emph{prospectif} et \emph{perspectif}, \emph{concret} et \emph{concomitant} change le sens de l'analyse.
\emph{Exemple :} « Baroka \emph{infère} sa puissance » (il la déduit ? Non). $\rightarrow$ « Baroka \emph{affirme} / \emph{démontre} / \emph{manifeste} sa puissance. »
\emph{Exemple :} « La modernité est \emph{concomitante} à la tradition » (simultanée ? Pas forcément). $\rightarrow$ « La modernité est \emph{intégrée} / \emph{domestiquée} par la tradition. »
\end{attention}

\begin{exercice}[Entraînement lexical : Paronymie et antonymie]
\begin{enumerate}
    \item Choisissez le paronyme correct entre parenthèses :
    \begin{enumerate}
        \item La stratégie de Baroka est (prospective / perspective) : il anticipe les coups.
        \item Lakunle (infère / infirme) la validité de la coutume par des arguments juridiques.
        \item Le sort de Sidi n'est pas (éminent / imminent) au début de la pièce.
        \item La dot est un geste (concret / concomitant) qui scelle l'union.
    \end{enumerate}
    \item Qualifiez le type d'antonymie (graduée, complémentaire, réciproque) pour les paires suivantes, en justifiant par la pièce :
    \begin{enumerate}
        \item \emph{Baroka / Lakunle} (rivaux)
        \item \emph{Manipuler / Être manipulé} (Sadiku / Baroka)
        \item \emph{Païen / Chrétien} (selon Lakunle)
    \end{enumerate}
\end{enumerate}
\end{exercice}

\begin{corrige}[Exercice n° 15-4 — Lexique]
\begin{enumerate}
    \item \textbf{a.} \emph{Prospective} (tournée vers l'avenir, anticipation). \textbf{b.} \emph{Infirme} (contredit, nie la validité). \textbf{c.} \emph{Imminent} (qui va arriver bientôt ; \emph{éminent} = illustre). \textbf{d.} \emph{Concret} (réel, tangible ; \emph{concomitant} = simultané).
    \item \textbf{a.} \emph{Graduée} (rivaux sur une échelle de pouvoir/âge/stratégie ; Baroka domine mais Lakunle existe). \textbf{b.} \emph{Réciproque} (si Baroka manipule Sadiku, Sadiku est manipulée par Baroka ; relation corrélative). \textbf{c.} \emph{Complémentaire} (pour Lakunle, on est l'un ou l'autre, pas de milieu ; vision binaire dogmatique).
\end{enumerate}
\end{corrige}

\courssub{L'exposé oral : préparation et présentation}

Le programme exige la maîtrise de l'\cle{exposé oral} (préparation, structure, prise de parole). Le débat « Dot vs Amour » (voir \emph{Expérience} ci-dessous) est un entraînement direct. Voici la méthode standard pour un exposé réussi (5 à 10 minutes).

\begin{methode}[Réussir son exposé oral : 5 étapes]
\begin{enumerate}
    \item \textbf{Analyser le sujet et délimiter :} Questionner les termes (ex: « La dot : achat ou alliance ? »). Formuler une \textbf{problématique} claire (ex: « La dot est-elle un obstacle à la liberté de la femme ou le fondement de sa dignité sociale ? »).
    \item \textbf{Documenter et structurer :} Rechercher arguments (juridiques, anthropologiques, économiques, littéraires). Organiser en \textbf{plan dialectique} (Thèse / Antithèse / Synthèse) ou \textbf{thématique} (Arguments pour / Arguments contre / Dépassement). Prévoir une \textbf{introduction} (accroche, définition, problématique, annonce du plan) et une \textbf{conclusion} (réponse, ouverture).
    \item \textbf{Rédiger les notes (pas le texte intégral) :} Fiches bristol ou diapositives : mots-clés, chiffres, citations courtes (Soyinka, Code civil, proverbes), structure logique. \textbf{Interdiction de lire.}
    \item \textbf{Répéter (chronométrer) :} S'enregistrer, travailler la voix (débit, articulation, intonation), le regard (balayage de l'auditoire), la gestuelle (ouverte, ancrée), la gestion du stress (respiration).
    \item \textbf{Présenter et interagir :} Annoncer le plan. Signaler les transitions (« Premier argument... », « En revanche... », « Pour conclure... »). Gérer le temps. Accueillir les questions (reformuler, répondre synthétiquement, assumer l'inconnu).
\end{enumerate}
\end{methode}

\begin{aretenir}[Grille d'auto-évaluation de l'exposé (Probatoire)]
\begin{tabularx}{\linewidth}{|X|c|c|c|}\hline
\rowcolor{popblueL}
\textbf{Critère} & \textbf{Non acquis} & \textbf{En cours} & \textbf{Acquis} \\ \hline
\textbf{Contenu :} Pertinence, exactitude, richesse arguments/exemples & & & \\ \hline
\textbf{Structure :} Intro/Plan/Transitions/Conclusion claires & & & \\ \hline
\textbf{Expression orale :} Débit, articulation, vocabulaire précis (paronymes !), syntaxe correcte & & & \\ \hline
\textbf{Communication :} Regard, posture, adaptation au public, temps respecté & & & \\ \hline
\textbf{Argumentation :} Thèses claires, arguments étayés, registers variés (logos/ethos/pathos) & & & \\ \hline
\end{tabularx}
\end{aretenir}

\courssub{Lien avec le contexte camerounais : coutumes villageoises et modernité urbaine}

La lecture de \emph{Le Lion et la perle} trouve une résonance immédiate dans le vécu des élèves camerounais. Le conflit Baroka/Lakunle n'est pas une pièce de musée ; il se joue chaque jour dans les chefferies, les familles et les administrations.

\begin{savaistu}[Le Lion et la perle au Cameroun : réalités vécues]
\begin{itemize}
    \item \textbf{La dot (le « prix de la mariée ») :} Sujet de débat permanent. Lakunle incarne le jeune diplômé urbain qui refuse de « payer » sa femme (argument : « on n'achète pas un être humain »). Baroka incarne le chef de famille villageois pour qui la dot scelle l'alliance entre deux lignées, honore la femme et ses parents. Au Cameroun, le Code civil (Loi 81-02) reconnaît le mariage coutumier et la dot, mais la jurisprudence veille à ce qu'elle ne soit pas « excessive ». La tension est vive entre \emph{droit moderne} (consentement, égalité) et \emph{droit coutumier} (alliance, transfert de filiation).
    \item \textbf{La polygamie :} Baroka l'assume (épouses senior, junior). Lakunle la rejette au nom de la monogamie chrétienne/laïque. Au Cameroun, le mariage coutumier est potentiellement polygamique ; le mariage civil est monogame. Beaucoup de Camerounais naviguent entre les deux statuts (mariage civil + coutumier), créant des situations juridiques complexes que la pièce préfigure.
    \item \textbf{L'appropriation de la modernité par la tradition :} Baroka fait frapper un timbre à son effigie ; aujourd'hui, les chefferies traditionnelles ont des sites web, des pages Facebook, des comptes WhatsApp pour gérer les conflits fonciers. Le « Lion » moderne ne refuse pas le smartphone ; il l'utilise pour asseoir son autorité. C'est la leçon centrale de Soyinka : \textbf{la tradition n'est pas l'opposé de la modernité, elle en est un usager stratégique}.
    \item \textbf{La place de la femme (Sidi) :} Ni soumise passive, ni libérée à l'occidentale. Sidi négocie sa valeur : elle veut la dot (reconnaissance), elle veut le chef (puissance), elle refuse l'instituteur qui ne lui offre que des mots. C'est une figure de \emph{féminisme situé} : l'autonomie par l'inscription dans sa culture, pas par la rupture.
\end{itemize}
\end{savaistu}

\begin{experience}[Activité de classe : Débat argumenté « Dot vs Amour »]
\textbf{Objectif :} Mettre en pratique l'analyse argumentative (thèses, arguments, registres) et l'exposé oral sur un sujet de société camerounais.
\textbf{Protocole :}
\begin{enumerate}
    \item Constituer deux groupes : « Groupe Baroka » (défense de la dot comme institution sociale) et « Groupe Lakunle » (défense du mariage sans dot comme droit moderne).
    \item Chaque groupe prépare 3 arguments structurés (Thèse + Argument + Exemple/Preuve) en 10 min (préparation d'exposé).
    \item Débat contradictoire (15 min) : un orateur par groupe, tour de parole 1 min (simulation d'exposé court).
    \item Synthèse du professeur : repérage des registres utilisés (juridique, affectif, économique, religieux), des fallaces (homme de paille, généralisation abusive), et de l'efficacité rhétorique.
\end{enumerate}
\textbf{Observation attendue :} Les élèves « Baroka » mobilisent souvent le \emph{pathos} (respect des parents, honneur de la fille) et l'\emph{ethos} (coutume = sagesse des ancêtres). Les élèves « Lakunle » mobilisent le \emph{logos} (Code civil, CEDAW, égalité homme/femme) mais peinent à incarner leurs arguments dans le réel affectif.
\textbf{Interprétation :} Cela confirme l'analyse de la pièce : l'argumentation « moderne » purement logique perd face à l'argumentation « traditionnelle » ancrée dans le symbolique et l'affectif, sauf si elle sait s'hybrider (ex: « Je paie la dot \emph{et} je respecte l'égalité dans le foyer »).
\end{experience}

\begin{popfigure}
\begin{tikzpicture}[
    timeline/.style={very thick, popblue},
    event/.style={circle, fill=poporange, draw=popdark, thick, minimum size=8pt},
    labelbox/.style={rectangle, draw=popdark, fill=popcream, rounded corners=3pt, font=\small, align=center, inner sep=3pt},
    datelabel/.style={font=\small\bfseries, popdark, anchor=north}
]
% Axe principal
\draw[timeline] (0,0) -- (16,0);
% Événements (5 moments forts)
\foreach \x/\label/\desc in {
    0.5/1/La danse de la perte de la virginité\\ (Lakunle refuse de payer la dot),
    3.5/2/La venue du photographe\\ (Sidi devient « star » locale),
    7/3/Le stratagème de Baroka\\ (Fausse impuissance annoncée à Sadiku),
    10.5/4/La séduction du bureau\\ (Sidi vaincue par le timbre/la ruse),
    14/5/Le mariage et la danse finale (Triomphe de la tradition)} {
    \node[event] (e) at (\x,0) {};
    \node[datelabel] at (\x, -0.6) {\label};
    \node[labelbox, anchor=south] at (\x, 1.2) {\desc};
    % Tiges
    \draw[poporange, thick] (\x,0) -- (\x, 0.6);
}
% Flèches de sens
\draw[->, >=Stealth, poppurple, thick] (0.5, 2.5) -- (3.5, 2.5) node[midway, above, font=\small, poppurple] {Provocation / Modernité};
\draw[->, >=Stealth, popgreen, thick] (7, 2.5) -- (10.5, 2.5) node[midway, above, font=\small, popgreen] {Ruse / Appropriation};
\draw[->, >=Stealth, popblue, thick] (10.5, 2.5) -- (14, 2.5) node[midway, above, font=\small, popblue] {Consécration / Ordre rétabli};
% Titre
\node[font=\bfseries\large, popdark] at (8, 3.8) {Frise chronologique : Les 5 moments forts de \emph{Le Lion et la perle}};
\end{tikzpicture}
\legende{Frise des cinq séquences dramatiques majeures : de la provocation moderniste (1-2) à la ruse traditionaliste (3-4) jusqu'à la consécration rituelle (5).}
\end{popfigure}

\begin{aretenir}[Synthèse pour la contraction : les 3 piliers de l'analyse Soyinka]
Pour réussir la contraction d'un extrait du \emph{Lion et la perle}, gardez en mémoire ce triptyque :
\begin{enumerate}
    \item \textbf{Le conflit est structurel, pas anecdotique :} Il oppose deux \emph{visions du monde} (cosmovision yoruba vs universalisme occidental), pas deux simples prétendants.
    \item \textbf{La langue est le champ de bataille :} Baroka parle \emph{yoruba} (proverbes, images, oralité) ; Lakunle parle \emph{anglais « bookish »} (abstractions, néologismes). Le choix de Sidi est aussi un choix linguistique.
    \item \textbf{L'ironie de l'auteur est le guide :} Soyinka ne prend pas parti pour Baroka (il critique la ruse, la polygamie, le pouvoir personnel) ni pour Lakunle (il critique l'aliénation, l'incompétence). Il montre l'\emph{impasse} d'une modernité sans ancrage et d'une tradition sans éthique. La contraction doit restituer cette \emph{complexité}, pas un manichéisme.
\end{enumerate}
\end{aretenir}

\begin{exercice}[Entraînement à la lecture méthodique n° 5 — Extrait du dénouement]
Lisez l'extrait ci-dessous (fin de l'acte 3, Sidi entre en scène après la nuit de noce) et répondez aux questions.

\emph{Extrait :} « Sidi entre, rayonnante. Elle porte le pagne de mariée. Elle ignore Lakunle, va droit à Sadiku, l'embrasse. — « Maman, la machine ne marchait pas ! » [...] Lakunle, furieux : « Tu as couché avec ce vieux singe ! » Sidi, riant : « Et je suis sa femme. La dot est payée. » Elle danse. Lakunle, seul, puis rejoignant une jeune fille : « Viens, dansons. » »

\begin{enumerate}
    \item Repérez les \textbf{deux registres de langue} qui s'affrontent dans les répliques de Lakunle et de Sidi. Citez un exemple pour chacun.
    \item Identifiez l'\textbf{argument décisif} de Sidi (« La dot est payée »). Quelle en est la portée \textbf{juridique} et \textbf{symbolique} dans la société yoruba (et camerounaise) ?
    \item Analysez le \textbf{dernier geste de Lakunle} (danser avec une autre). Victoire, défaite ou compromis ? Justifiez par le texte.
    \item \textbf{Contraction (50 mots ± 10\%) :} Résumez cet extrait en appliquant la méthode des 3 idées essentielles (voir \emph{Exercice résolu} ci-dessus).
\end{enumerate}
\end{exercice}

\begin{corrige}[Exercice n° 15-3 — Lecture méthodique n° 5]
\begin{enumerate}
    \item \textbf{Registres :} Lakunle : \emph{registre injurieux / familier dégradé} (« vieux singe » = insulte zoomorphique, déni d'humanité). Sidi : \emph{registre affirmatif / juridique-coutumier} (« La dot est payée » = performatif, acte de droit coutumier ; « Maman » = respect codifié).
    \item \textbf{Argument « La dot est payée » :} Juridiquement : le mariage coutumier est valide, irréversible, opposable à tous (Lakunle y compris). Symboliquement : Sidi n'est pas « achetée » (thèse de Lakunle), elle est \emph{honorée} ; la dot scelle l'alliance lignagère, elle donne à Sidi un statut social (épouse légitime, future mère de lignée) que le « mariage d'amour sans dot » de Lakunle ne lui offrait pas.
    \item \textbf{Geste de Lakunle :} \textbf{Compromis amer / Résignation active.} Il ne part pas (exil), il ne tue pas (violence), il ne convainc pas (échec du logos). Il \emph{danse} : il intègre le rituel communautaire qu'il méprisait. C'est une adaptation mimétique : il fait « comme si » il acceptait, tout en restant en marge (avec une « jeune fille », pas Sidi). La tradition a absorbé sa contestation.
    \item \textbf{Contraction (modèle ~50 mots) :}
    \emph{Revenue du palais de Baroka, Sidi triomphe : la machine moderne a failli, la dot coutumière est payée, scellant son mariage. Elle rejette Lakunle, dont l'injure (« vieux singe ») masque l'impuissance. La danse finale unit la communauté ; Lakunle, marginalisé, y participe à son tour, signe que le progrès n'a pas désarmé la tradition.} (52 mots)
\end{enumerate}
\end{corrige}

\coursec{Les tonalités : comique, dramatique, satirique}

Après avoir analysé la structure argumentative du texte et pratiqué la contraction lors des séances précédentes, nous poursuivons notre lecture méthodique du \emph{Lion et la perle} (lectures n° 4 et 5) en nous attachant maintenant à la dimension \textbf{émotionnelle} et \textbf{critique} de l'œuvre. Une même scène peut faire rire, émouvoir ou indigner : c'est le jeu des \cle{tonalités} (ou registres). Les identifier, c'est comprendre l'intention profonde de Wole Soyinka : divertir pour mieux instruire, rire pour mieux corriger.

\courssub{Qu'est-ce qu'une tonalité (ou registre) ?}

La \cle{tonalité} (souvent appelée \cle{registre} en classe de Première) désigne l'\textbf{effet émotionnel et intellectuel dominant} que le texte produit sur le lecteur ou le spectateur. Elle résulte d'un faisceau de choix linguistiques (vocabulaire, syntaxe, figures de style) et de choix dramatiques (situation, gestes, réactions des personnages). Ce n'est pas une étiquette décorative : c'est la \textbf{porte d'entrée du sens}. Dans \emph{Le Lion et la perle}, Soyinka ne se contente pas de raconter une querelle de mariage ; il utilise le rire (comique), l'angoisse (dramatique) et la moquerie critique (satirique) pour interroger l'avenir de l'Afrique face à la modernité.

\begin{definition}[Tonalité / Registre]
La tonalité est l'attitude affective et intellectuelle de l'auteur (ou de l'énonciateur) à l'égard de son sujet, manifestée par des procédés linguistiques et scéniques précis, et destinée à provoquer une réaction déterminée chez le récepteur (rire, larmes, indignation, réflexion, peur, etc.). Les trois tonalités majeures au programme sont le \textbf{comique}, le \textbf{dramatique} et le \textbf{satirique}.
\end{definition}

\begin{attention}[Ne pas confondre genre et tonalité]
\emph{Le Lion et la perle} est une \textbf{comédie} (genre théâtral : dénouement heureux, personnages types). Mais \textbf{à l'intérieur} de cette comédie, on trouve des passages dramatiques (la peur de Sidi, la ruse de Baroka) et une visée satirique globale. Une œuvre d'un genre donné peut traverser plusieurs tonalités.
\end{attention}

\courssub{Le comique : rire pour mieux voir}

Le \cle{comique} vise à provoquer le rire ou le sourire. Il naît d'un \textbf{décalage} entre ce qui est attendu et ce qui advient, entre la norme et la transgression. On distingue traditionnellement quatre formes que Soyinka manie avec virtuosité :

\begin{itemize}
    \item \textbf{Le comique de mots (verbal) :} jeux sur la polysémie, calembours, quiproquos linguistiques. \emph{Exemple :} Lakunle utilise un anglais ampoulé, « académique », qui sonne faux face au pidgin ou à l'anglais simple des villageois. Ses métaphores mal maîtrisées (« \emph{the sweetness of the honey} » pour parler d'amour) créent un effet de ridicule.
    \item \textbf{Le comique de gestes (visuel / burlesque) :} chutes, mimiques, gesticules excessifs. \emph{Exemple :} La danse mimée du « \emph{Surveyor} » (l'arpenteur) par Lakunle au début de la pièce : il singe la modernité technique avec une rigidité grotesque.
    \item \textbf{Le comique de situation :} quiproquo, renversement inattendu, situation absurde. \emph{Exemple :} Baroka, le Bale, qui feint l'impuissance devant Sadiku pour tester sa discrétion (scène de la chambre à coucher). Le spectateur sait, pas Sadiku : c'est de l'\cle{ironie dramatique} source de comique.
    \item \textbf{Le comique de caractère :} traits exagérés, obsessionnels, qui définissent un « type ». \emph{Exemple :} Lakunle, l'instituteur « moderniste » qui refuse la dot par principe idéologique mais veut la femme ; Baroka, le rusé traditionaliste qui collectionne les femmes comme des timbres.
\end{itemize}

\begin{exemplebox}[La danse de Baroka : comique de situation et de caractère]
Dans la scène finale (Matin, près de l'arbre de l'odéon), Baroka lutte avec son lutteur. C'est une scène de \textbf{comique de geste} (le vieux chef qui lutte, l'effort physique simulé) et de \textbf{comique de situation} : il vient de « conquérir » Sidi par la ruse, il affiche une vitalité triomphante qui contredit la rumeur d'impuissance qu'il a lui-même répandue. Le rire du public bascule alors vers l'admiration teintée d'effroi : le « Lion » a rugi. Le comique masque ici une démonstration de pouvoir.
\end{exemplebox}

\courssub{Le dramatique : la tension au cœur de l'action}

Le \cle{dramatique} (ou \cle{pathétique}) suscite l'\textbf{émotion intense}, la pitié, la peur, l'angoisse face à un destin qui échappe au personnage. Il repose sur le \textbf{conflit} (interne ou externe), l'\textbf{urgence} et l'\textbf{irréversible}. Dans la pièce, le drame ne tue pas (c'est une comédie), mais il \textbf{menace} l'équilibre psychologique des personnages.

\begin{itemize}
    \item \textbf{Le dilemme de Sidi :} Elle est tiraillée entre deux modèles de masculinité et deux visions du monde. Choisir Lakunle, c'est accepter la modernité occidentale (monogamie, pas de dot, « amour romantique ») mais risquer de perdre son ancrage social et sa dignité traditionnelle (la dot est le prix de sa valeur). Choisir Baroka, c'est accepter la tradition (polygamie, dot, pouvoir patriarcal) mais assurer son statut social et la pérennité du lignage.
    \item \textbf{La scène de la séduction (Soir, chambre de Baroka) :} C'est le sommet du \textbf{dramatique}. Sidi, piégée, découvre la supercherie. La tension monte : Baroka avance, Sidi recule. Le dialogue se fait serré, les silences pèsent. Le pathétique naît de la \textbf{prise de conscience brutale} de Sidi : sa beauté (la Perle) ne la protège pas, elle l'expose. Elle est « prise au piège » comme l'animal dont parle le titre.
    \item \textbf{Le monologue de Sadiku (après l'annonce de l'impuissance) :} Elle danse, triomphe, croit avoir le pouvoir. Le spectateur ressent une \emph{ironie dramatique} teintée de pathétique : on sait qu'elle est manipulée, sa joie est vaine, ce qui la rend touchante et pitoyable.
\end{itemize}

\begin{propriete}[Mécanisme du pathétique au théâtre]
Le pathétique théâtral repose sur trois piliers :
1. \textbf{L'enjeu vital :} honneur, liberté, vie, amour, identité.
2. \textbf{L'impuissance relative :} le personnage ne maîtrise pas tous les paramètres (ici, Sidi ignore la ruse).
3. \textbf{L'empathie du spectateur :} identification à la victime (Sidi) ou angoisse pour le manipulateur (Baroka).
\end{propriete}

\courssub{Le satirique : la critique par le rire}

Le \cle{satirique} est l'arme majeure de Soyinka. Ce n'est pas seulement « faire rire », c'est \textbf{rire \emph{de} quelque chose pour le corriger}. La satire vise des \textbf{travers collectifs} (coutumes archaïques, mimétisme culturel, corruption du pouvoir, naïveté juvénile) en les grossissant jusqu'à l'absurde.

\begin{definition}[La satire]
La satire est un discours (littéraire, graphique, oral) qui \textbf{critique les travers d'une société, d'un groupe ou d'individus en les tournant en dérision}. Elle utilise le rire comme arme correctrice : \emph{castigat ridendo mores} (elle corrige les mœurs par le rire). Elle suppose une \textbf{norme implicite} (ce qui est bien, juste, vrai) que le texte montre bafouée.
\end{definition}

\begin{exemplebox}[L'ironie de Soyinka sur Lakunle et la dot]
Lakunle refuse de payer la dot (\emph{bride price}), la qualifiant de « barbarie », d'« achat de femme », d'« esclavage ». Il veut épouser Sidi « à la moderne », par amour pur. \textbf{L'ironie satirique} éclate quand on voit que :
\begin{enumerate}
    \item Lakunle \emph{veut} Sidi (désir charnel, possession) tout en niant le cadre coutumier qui légitime cette union.
    \item Il refuse de payer \emph{maintenant}, mais il bénéficiera du travail de Sidi (champs, maison, enfants) \emph{gratuitement}. Il veut les avantages de la tradition (femme, descendance) sans en payer le prix symbolique.
    \item Son langage « moderne » cache un conservatisme profond : il veut une femme soumise, « moderne » seulement en surface (robe, danse), mais traditionnelle dans le foyer.
\end{enumerate}
Soyinka raille ici le \textbf{modernisme superficiel} : une modernité de façade, vide de sens, qui nie l'économie réelle du mariage africain.
\end{exemplebox}

\begin{aretenir}[Les cibles de la satire dans \emph{Le Lion et la perle}]
\begin{itemize}
    \item \textbf{Le traditionalisme aveugle / gérontocratique :} Baroka incarne le pouvoir qui se perpétue par la ruse, la polygamie, le mépris de la jeunesse et de la femme (objet de collection). La statue du « \emph{Bale} » qu'il veut faire ériger symbolise l'égocentrisme du pouvoir.
    \item \textbf{Le modernisme superficiel / aliéné :} Lakunle incarne l'intellectuel déraciné, qui singe l'Occident (vocabulaire, vêtements, idées) sans en comprendre la substance, et qui méprise sa propre culture.
    \item \textbf{La vanité féminine / la naïveté :} Sidi, aveuglée par sa beauté (les photos du magazine), perd sa lucidité. Sadiku, par bavardage, trahit le secret du chef.
\end{itemize}
\end{aretenir}

\courssub{Procédés de la satire : l'arsenal de l'ironie}

Pour tourner en dérision, Soyinka utilise des procédés rhétoriques précis qu'il faut savoir repérer :

\begin{itemize}
    \item \textbf{L'\cle{ironie} :} Dire le contraire de ce que l'on pense, ou faire dire au personnage le contraire de la réalité. \emph{Exemple :} Lakunle qui dit « Je ne paierai pas la dot, c'est barbare » alors qu'il cherche à obtenir la femme sans contrepartie.
    \item \textbf{La \cle{caricature} :} Grossissement outrancier d'un trait physique ou moral. \emph{Exemple :} La description de la « machine à dupliquer » (le cyclostyle) par Lakunle comme un miracle de la science, ou la bedaine de Baroka, sa force de lutteur, sa virilité exhibée.
    \item \textbf{L'\cle{exagération} (hyperbole) :} Dépasser la mesure pour souligner le ridicule. \emph{Exemple :} Les métaphores filées de Lakunle (« \emph{the honey of my tongue, the sugar of my lips} ») qui transforment une déclaration d'amour en cours de chimie culinaire.
    \item \textbf{L'\cle{antiphrase} :} Figure d'ironie où l'on affirme le contraire de ce que l'on veut dire, sans marqueur explicite. \emph{Exemple :} Sadiku traitant Baroka de « grand homme », de « lion » au moment où elle croit le détruire par la rumeur d'impuissance.
    \item \textbf{Le \cle{paradoxe} :} Affirmation apparemment contradictoire qui révèle une vérité cachée. \emph{Exemple :} « Le vieux est plus fort que le jeune » (Baroka vs Lakunle / Lutteur). La tradition (vieille) vainc la modernité (jeune) par la ruse, pas par la force brute.
\end{itemize}

\begin{methode}[Identifier la tonalité dominante d'un extrait]
\begin{enumerate}
    \item \textbf{Contextualiser :} Qui parle ? À qui ? Où ? Quand dans l'intrigue ? Quel enjeu ?
    \item \textbf{Relever les indices lexicaux :} Champs lexicaux (joie/tristesse, noble/trivial, concret/abstrait), connotations, registres de langue (soutenu, familier, technique).
    \item \textbf{Repérer les figures de style et procédés :} Ironie, hyperbole, métonymie, questions rhétoriques, images, répétitions, silences didascalies (gestes, ton).
    \item \textbf{Analyser la situation dramatique :} Conflit, quiproquo, révélation, suspense, renversement.
    \item \textbf{Nommer l'effet produit :} Rire franc (comique) ? Rire jaune / réflexion critique (satirique) ? Angoisse / pitié / émotion forte (dramatique) ?
    \item \textbf{Conclure :} « La tonalité dominante est \textbf{X} car [procédés] produisent [effet], ce qui sert l'intention de l'auteur : [critique / divertissement / émotion]. » \textit{Nota bene :} Souvent, une tonalité \textbf{domine} mais une seconde \textbf{teinte} la scène (ex : comique satirique, dramatique ironique). Il faut le signaler.
\end{enumerate}
\end{methode}

\courssub{Tableau récapitulatif : tonalités, procédés et effets dans la pièce}

\begin{popfigure}
\begin{center}
\begin{tikzpicture}[
    cell/.style={rectangle, draw=popline, minimum height=1.1cm, align=center, text width=4.8cm, font=\small},
    head/.style={cell, fill=popblue, text=white, font=\small\bfseries},
    rowc/.style={cell, fill=popblueL!30},
    rowd/.style={cell, fill=poporangeL!30},
    rows/.style={cell, fill=popgreenL!30},
]
% En-têtes
\node[head] (t1) at (0,0) {Tonalité};
\node[head] (t2) at (5.0,0) {Procédés principaux};
\node[head] (t3) at (10.0,0) {Effet recherché / Exemples dans \emph{Le Lion et la perle}};

% Ligne Comique
\node[rowc] (c1) at (0,-1.1) {\textbf{Comique}};
\node[rowc] (c2) at (5.0,-1.1) {Quiproquo, calembour, gestuelle burlesque, comique de répétition, inversion des rôles};
\node[rowc] (c3) at (10.0,-1.1) {Détente, plaisir du spectacle, caractérisation. \textit{Ex: Danse de l'arpenteur (Lakunle), scène de la lutte (Baroka), bavardages de Sadiku.}};

% Ligne Dramatique
\node[rowd] (d1) at (0,-2.2) {\textbf{Dramatique / Pathétique}};
\node[rowd] (d2) at (5.0,-2.2) {Champ lexical de la peur/piège, silences, didascalies lentes, monologue intérieur, tension croissante, ironie dramatique};
\node[rowd] (d3) at (10.0,-2.2) {Empathie, angoisse, prise de conscience. \textit{Ex: Sidi dans la chambre de Baroka (scène de séduction), découverte de la supercherie, dilemme final au carrefour.}};

% Ligne Satirique
\node[rows] (s1) at (0,-3.3) {\textbf{Satirique}};
\node[rows] (s2) at (5.0,-3.3) {Ironie, caricature, hyperbole, antiphrase, paradoxe, décalage langue/réalité, mise en abyme (théâtre dans le théâtre)};
\node[rows] (s3) at (10.0,-3.3) {Réflexion critique, correction des mœurs, dénonciation. \textit{Ex: Refus de la dot par Lakunle, collection de femmes de Baroka, photos de Sidi dans le magazine, le projet de timbre-poste.}};
\end{tikzpicture}
\end{center}
\legende{Tableau de synthèse : les trois tonalités majeures, leurs procédés privilégiés et leurs manifestations dans \emph{Le Lion et la perle}.}
\end{popfigure}

\courssub{Échelle des tonalités : de la comédie légère à la satire mordante}

Les scènes de la pièce ne sont pas monochromes. On peut visualiser la dominante tonale de chaque grande séquence sur un curseur allant du \textbf{divertissement pur} (comique léger) à la \textbf{charge critique violente} (satire mordante), en passant par l'\textbf{émotion tensionnelle} (dramatique).

\begin{popfigure}
\begin{center}
\begin{tikzpicture}[
    scale=0.9,
    every node/.style={font=\small},
    arrow/.style={->, >=Stealth, thick, popblue},
    tick/.style={thick, popdark},
    label/.style={align=center, text width=2.8cm},
    scenelabel/.style={align=center, text width=3cm, font=\tiny, popmuted}
]
% Axe principal
\draw[thick, popdark] (0,0) -- (14,0);
% Graduations et labels de tonalité
\foreach \x/\txt in {0/Comique léger, 3.5/Comique satirique, 7/Satire équilibrée, 10.5/Satire mordante, 14/Dramatique intense} {
    \draw[tick] (\x, -0.15) -- (\x, 0.15);
    \node[label, anchor=north] at (\x, -0.4) {\txt};
}
% Flèches de sens
\draw[arrow] (0.2, 0.5) -- (6.8, 0.5) node[midway, above, font=\tiny] {Montée de la critique sociale};
\draw[arrow] (13.8, 0.5) -- (7.2, 0.5) node[midway, above, font=\tiny] {Montée de la tension existentielle};

% Positionnement des scènes clés (cercles colorés)
% 1. Danse de l'arpenteur (Début Matin)
\fill[popgreen!80] (1, 0.8) circle (0.15);
\node[scenelabel, anchor=south, align=center] at (1, 1.0){1. Danse de l'arpenteur\\(Lakunle mime la modernité)};

% 2. Négociation dot / refus Lakunle (Matin)
\fill[popgreen!60!poporange] (3, 1.1) circle (0.15);
\node[scenelabel, anchor=south, align=center] at (3, 1.3){2. Refus de la dot\\(Idéologie vs Réalité)};

% 3. Sadiku apprend impuissance / Danse victoire (Midi)
\fill[poporange!80] (5.5, 0.8) circle (0.15);
\node[scenelabel, anchor=south, align=center] at (5.5, 1.0){3. Triomphe de Sadiku\\(Ironie dramatique)};

% 4. Scène de séduction Baroka / Sidi (Soir) -- POINT CENTRAL DRAMATIQUE
\fill[popblue!80] (8, -0.8) circle (0.15);
\node[scenelabel, anchor=north, align=center] at (8, -1.0){4. Séduction / Piège\\(Pathétique + Satire pouvoir)};

% 5. Lutte Baroka / Lutteur + Final (Matin suivant)
\fill[poppurple!80] (11, 0.8) circle (0.15);
\node[scenelabel, anchor=south, align=center] at (11, 1.0){5. Lutte \& Mariage\\(Satire puissance + Comique final)};

% 6. Lakunle seul / Départ (Fin)
\fill[poppink!80] (13, 0.5) circle (0.15);
\node[scenelabel, anchor=south, align=center] at (13, 0.7){6. Lakunle tourné en dérision\\(Satire finale de l'aliéné)};

% Légende couleurs
\begin{scope}[yshift=-3.5cm, xshift=1cm]
\node[label, anchor=west] at (0,0) {\textbf{Légende des scènes :}};
\foreach \y/\col/\txt in {0/popgreen!80/Comique dominant, 0.5/poporange/Satire comique, 1/popblue/Dramatique pathétique, 1.5/poppurple/Satire mordante, 2/poppink/Satire finale} {
    \fill[\col] (0, -\y) circle (0.1);
    \node[label, anchor=west, font=\tiny] at (0.3, -\y) {\txt};
}
\end{scope}
\end{tikzpicture}
\end{center}
\legende{Curseur des tonalités dominantes selon les grandes séquences de la pièce. La flèche bleue montre la bascule vers le dramatique/pathétique au centre (scène de séduction), la flèche orange la montée de la satire.}
\end{popfigure}

\begin{attention}[Une scène = plusieurs tonalités : dire laquelle DOMINE]
Ne dites jamais : « Cette scène est comique ET dramatique ET satirique. » Dites : « La tonalité \textbf{dominante} est \textbf{satirique} (critique du pouvoir de Baroka), \textbf{teintée} de \textbf{dramatique} (angoisse de Sidi) et d'une pointe de \textbf{comique} (gestes de Baroka). » Justifiez la dominance par l'\textbf{intention finale} : que veut faire ressentir/penser l'auteur \emph{in fine} ? Ici, la satire du pouvoir patriarcal l'emporte sur l'anecdote comique.
\end{attention}

\courssub{Exercice résolu : identifier la tonalité de trois extraits}

\begin{exoresolu}[Analyse de la tonalité dominante]
\textbf{Extrait 1 (Matin, p. 12-13) :} Lakunle : « \emph{Je ne paierai pas la dot. C'est une coutume barbare, dégradante, humiliante... Je veux t'épouser selon la loi chrétienne, par amour pur.} » Sidi : « \emph{Tu parles comme un livre. Ma dot doit être payée. Sinon, on dira que je ne vaux rien.} »

\textbf{Extrait 2 (Soir, p. 45-46) :} Baroka : « \emph{Viens, ma perle. Le lion ne mange pas la perle, il la porte au cou.} » (Il l'attire vers lui). Sidi : « \emph{Laissez-moi ! Vous avez menti ! Vous n'êtes pas impuissant !} » (Elle lutte, mais il est plus fort).

\textbf{Extrait 3 (Matin, p. 62) :} Lakunle (seul, après le départ de Sidi pour le mariage) : « \emph{Qu'importe ! Il y a d'autres filles... Je vais fonder un club des "Femmes Libérées". Nous danserons le "high-life" jusqu'à l'aube.} » (Il imite une danse moderne, ridicule).

\tcblower
\textbf{Solution détaillée :}

\textbf{Extrait 1 : Tonalité \textbf{SATIRIQUE} (dominante) teintée de COMique.}
\begin{itemize}
    \item \textit{Procédés :} Ironie (Lakunle donne des leçons de morale mais veut la femme « gratis »), Caricature du « moderne » (vocabulaire ampoulé « loi chrétienne », « amour pur » déconnecté du réel), Antiphrase (« coutume barbare » dit par celui qui profite de la coutume sans la payer), Décalage langue/réalité.
    \item \textit{Effet :} Rire critique. On rit \emph{de} Lakunle, pas \emph{avec} lui. Soyinka dénonce l'intellectuel aliéné qui nie les fondements économiques et symboliques de sa société.
\end{itemize}

\textbf{Extrait 2 : Tonalité \textbf{DRAMATIQUE} (dominante) teintée de SATIRIQUE.}
\begin{itemize}
    \item \textit{Procédés :} Champ lexical du piège/de la chasse (lion, perle, menti, impuissant), Didascalies de tension (lutte, attire, plus fort), Ironie dramatique (le spectateur sait la ruse, Sidi la découvre trop tard), Métaphore filée (Lion/Perle) qui fige le rapport de force.
    \item \textit{Effet :} Angoisse, pitié pour Sidi, horreur lucide pour la ruse de Baroka. La satire vise ici le \textbf{pouvoir masculin patriarcal} qui use de la ruse pour posséder le corps de la femme. Le pathétique sert la satire.
\end{itemize}

\textbf{Extrait 3 : Tonalité \textbf{COMIQUE} (dominante) teintée de SATIRIQUE.}
\begin{itemize}
    \item \textit{Procédés :} Comique de geste (danse ridicule seul), Comique de répétition (il recommence une stratégie qui a échoué), Ironie (il crée un club « Femmes Libérées» tout seul, sans femmes), Caricature du modernisme vide (« high-life »).
    \item \textit{Effet :} Rire franc, mais amer. Lakunle devient une figure burlesque, le clown de la modernité ratée. La satire se referme sur lui : l'aliéné reste seul face à son miroir déformant.
\end{itemize}
\end{exoresolu}

\courssub{Entraînement : Paronymie, Antonymie et Polysémie (rappel lexical)}

Le programme lie l'étude des tonalités à la précision lexicale (\textbf{paronymie}, \textbf{antonymie}, \textbf{polysémie}). Choisir le mot juste, c'est fixer la tonalité.

\begin{exemplebox}[Paronymie et Polysémie : pièges et nuances tonales]
\begin{itemize}
    \item \textbf{Paronymie (mots de forme proche, sens différents) :} \emph{Satirique} (adj, qui relève de la satire, critique) \textbf{/} \emph{Satyrique} (adj, relatif aux satyres, mythologie, lubricité). \textbf{Piège Bac :} « Le ton \textbf{satirique} de Soyinka » (critique) $\neq$ « Un dessin \textbf{satyrique} » (obscène).
    \item \textbf{Polysémie (un mot, plusieurs sens) :} \emph{Comique} (adj : qui fait rire / nom : genre théâtral / nom : artiste humoriste). \emph{Dramatique} (adj : intense, pathétique / nom : genre théâtral / nom : dramaturge $\rightarrow$ paronymie avec \emph{dramaturge}). \textit{Exemple :} « Une situation \textbf{comique} » (qui prête à rire) vs « Un auteur \textbf{comique} » (spécialiste du genre).
\end{itemize}
\end{exemplebox}

\begin{exemplebox}[Antonymie et structure du conflit]
La pièce se construit sur des \textbf{paires antonymiques} qui structurent les tonalités :
\begin{center}
\begin{tabularx}{\linewidth}{|l|X|X|}\hline
\textbf{Paire antonymique} & \textbf{Pôle A (Lakunle / Moderne)} & \textbf{Pôle B (Baroka / Tradition)} \\ \hline
\textit{Jeunesse / Vieillesse} & Jeune, neuf, futur & Vieux, ancien, passé \\ \hline
\textit{Modernité / Tradition} & École, livre, machine, loi & Coutume, oralité, force, ruse \\ \hline
\textit{Monogamie / Polygamie} & Un amour, une femme & Collection, pouvoir, lignage \\ \hline
\textit{Parole / Acte} & Discours théoriques, promesses & Actes concrets, résultats (lutte, séduction) \\ \hline
\end{tabularx}
\end{center}
Le \textbf{comique} naît souvent du choc de ces antonymes (le jeune qui parle comme un vieux livre). La \textbf{satire} naît de l'échec des deux pôles pris séparément. Le \textbf{dramatique} naît du déchirement de Sidi au milieu.
\end{exemplebox}

\courssub{Le savais-tu ? : La satire au Cameroun, une tradition vivante}

\begin{savaistu}[La satire camerounaise : de \emph{Le Popoli} aux réseaux sociaux]
Au Cameroun, la satire est un \textbf{genre majeur de l'espace public}, héritage direct de la tradition orale (le \emph{griot} qui « lave le linge sale » du chef) et de la presse d'opinion.
\begin{itemize}
    \item \textbf{La presse satirique :} \emph{Le Popoli} (créé en 1991), \emph{Le Messager Popoli}, \emph{Nyemba}. Leurs « Unes » utilisent la \textbf{caricature} (dessin de presse grossissant les traits des hommes politiques), le \textbf{titre antiphrastique} (\emph{« Le Président travaille dur... à se reposer »}), l'\textbf{ironie} (éditos faussement naïfs). C'est du \emph{Le Lion et la perle} en direct : on rit du pouvoir pour ne pas pleurer.
    \item \textbf{L'humour de scène :} Les troupes comme \emph{Les Guignols du Cameroun} (marionnettes), les one-man-shows (Dieudonné Mballa, Cabrel Nanjip) utilisent le \textbf{comique de situation} et le \textbf{comique de mots} (camfranglais, détournement de discours officiels) pour faire de la \textbf{satire sociale} (corruption, népotisme, vie chère, éducation).
    \item \textbf{Le lien avec Soyinka :} Comme Baroka, les « grands » sont tournés en dérision par leur « bedaine », leur « polygamie » (politique), leurs « timbres » (projets pharaoniques inachevés). Comme Lakunle, les « modernes » donneurs de leçons (ONG, experts, diaspora) sont raillés pour leur déconnexion du terrain.
\end{itemize}
\textbf{En classe :} Apportez une « Une » de \emph{Popoli} ou un extrait de sketch. Identifiez : 1. La cible. 2. Le procédé (caricature, ironie...). 3. La tonalité (satirique dominant). C'est la même méthode que pour \emph{Le Lion et la perle} !
\end{savaistu}

\begin{exercice}[Entraînement à l'identification tonale]
Relisez les extraits suivants de \emph{Le Lion et la perle} (ou consultez votre manuel) et pour chacun :
\begin{enumerate}
    \item Identifiez la \textbf{tonalité dominante} (Comique, Dramatique, Satirique).
    \item Citez \textbf{deux procédés précis} (lexique, figure, situation, didascalie) qui justifient votre choix.
    \item Expliquez en une phrase l'\textbf{intention de l'auteur} (ce qu'il veut nous faire penser/ressentir).
\end{enumerate}
\textbf{Extrait A :} La scène où Lakunle essaie d'enseigner la « danse moderne » aux écoliers (Matin, début).
\textbf{Extrait B :} Le monologue de Baroka face au lutteur : « \emph{Je suis le lion... Je ne vieillis pas.} » (Matin, fin).
\textbf{Extrait C :} Le dialogue final entre Lakunle et Sadiku : Lakunle veut payer la dot \emph{après} le mariage pour « régulariser ».
\end{exercice}


\begin{corrige}[Exercice n° 15 - Identification tonale]
\textbf{Extrait A :} \textbf{Comique} (dominant) / Satirique. \textit{Procédés :} Comique de geste (mimique ridicule), Comique de mots (vocabulaire technique inadapté « \emph{ballroom dancing} »), Caricature du moderne. \textit{Intention :} Ridiculiser l'importation servile de modèles culturels étrangers.
\textbf{Extrait B :} \textbf{Satirique} (dominant) / Dramatique (tension sur le vieillissement). \textit{Procédés :} Métaphore identitaire (Lion), Hyperbole (« Je ne vieillis pas »), Ironie (il vient de tricher pour prouver sa force). \textit{Intention :} Dénoncer le pouvoir personnel qui nie le temps et la loi pour s'éterniser.
\textbf{Extrait C :} \textbf{Satirique} (dominant) / Comique. \textit{Procédés :} Ironie situationnelle (Lakunle paie ce qu'il a refusé), Antiphrase (« régulariser »), Comique de répétition (il change d'avis par intérêt). \textit{Intention :} Montrer l'incohérence de l'aliéné : ses principes s'effondrent dès que son intérêt personnel (avoir Sidi) est en jeu.
\end{corrige}

\begin{aretenir}[Bilan : Les trois tonalités, trois lunettes sur le monde]
\begin{itemize}
    \item \textbf{Comique :} « Regarde comme c'est drôle / absurde / décalé. » $\rightarrow$ \textbf{Détente, caractérisation.}
    \item \textbf{Dramatique :} « Sens comme c'est grave / douloureux / urgent. » $\rightarrow$ \textbf{Empathie, tension, humanité.}
    \item \textbf{Satirique :} « Comprends comme c'est faux / injuste / ridicule / critique. » $\rightarrow$ \textbf{Réflexion, correction, lucidité.}
\end{itemize}
Dans \emph{Le Lion et la perle}, \textbf{Soyinka les tresse} : le comique attire le public, le dramatique l'attache aux personnages, la satire lui ouvre les yeux sur sa propre société. C'est la marque des grands classiques.
\end{aretenir}

\coursec{Reformuler les idées essentielles}

Après avoir analysé les \cle{tonalités} — comique, dramatique, satirique — qui colorent \emph{Le Lion et la Perle} et révèlent les intentions de Wole Soyinka, nous devons maintenant maîtriser l'outil indispensable pour rendre compte de tout texte : la \cle{reformulation}. Comprendre le ton d'un texte, c'en saisir l'âme ; le reformuler, c'en restituer le squelette avec ses propres mots. C'est la compétence centrale de la contraction de texte, épreuve reine du Probatoire technique.

\courssub{Principes et enjeux de la reformulation}

\begin{definition}[Reformulation]
Reformuler, c'est \textbf{exprimer la même idée avec d'autres mots et d'autres structures grammaticales}, en conservant rigoureusement le sens, la portée et la nuance de l'énoncé original, sans y ajouter d'interprétation personnelle ni d'exemple extérieur.
\end{definition}

La reformulation n'est pas une simple substitution lexicale. C'est une \cle{transposition syntaxique et sémantique} qui prouve que l'on a compris la logique interne du texte. Au Probatoire, la contraction de texte (environ 80 mots pour un texte de 400 mots, soit le cinquième environ) sanctionne toute phrase recopiée : c'est le \cle{plagiat}, lourdement pénalisé. À l'inverse, une reformulation fidèle mais allégée témoigne d'une maîtrise de la langue et de la pensée de l'auteur.

\begin{exemplebox}[Exemple chiffré : gain de concision]
Phrase originale (11 mots) : \emph{« Les jeunes délaissent massivement les filières agricoles »} \\
Reformulation (7 mots) : \emph{« La jeunesse abandonne en grand nombre l'agriculture »} \\
\cle{Outils utilisés} : synonymie (\emph{délaissent} $\rightarrow$ \emph{abandonne} ; \emph{filières agricoles} $\rightarrow$ \emph{agriculture}), changement de catégorie (\emph{jeunes} adjectif $\rightarrow$ \emph{jeunesse} nom), compression (\emph{massivement} $\rightarrow$ \emph{en grand nombre}).
\end{exemplebox}

\courssub{Les outils de la reformulation : la boîte à outils du résumeur}

Pour reformuler sans trahir, on actionne cinq leviers principaux. Le tableau ci-dessous synthétise les correspondances types.

\begin{center}
\rowcolors{1}{popcream}{white}
\begin{tabularx}{\linewidth}{|c|X|X|l|}
\hline
\rowcolor{popblueL} \textbf{N°} & \textbf{Phrase originale} & \textbf{Reformulation} & \textbf{Outil principal} \\
\hline
1 & \emph{« Car le sol est fertile, nous récoltons »} & \emph{« Puisque la terre est riche, la moisson abonde »} & Synonymie + connecteur \\
\hline
2 & \emph{« Il refuse la modernité, mais il n'est pas sot »} & \emph{« Bien qu'il rejette le progrès, il demeure lucide »} & Antonymie + concession \\
\hline
3 & \emph{« La décision du chef surprend le village »} & \emph{« Le village est surpris par la décision du chef »} & Voix passive \\
\hline
4 & \emph{« L'ignorance de Lakunle entraîne son échec »} & \emph{« Lakunle échoue par ignorance »} & Verbalisation + compression \\
\hline
5 & \emph{« Sidi choisit Baroka parce qu'il est puissant »} & \emph{« La puissance de Baroka détermine le choix de Sidi »} & Nominalisation + inversion \\
\hline
\end{tabularx}
\end{center}

\textbf{Détail des leviers.}\par
\begin{enumerate}
    \item \textbf{Synonymie et paraphrase lexicale} : remplacer un mot par son équivalent contextuel (\emph{délaisser} $\rightarrow$ \emph{abandonner}, \emph{fertile} $\rightarrow$ \emph{riche}). Attention aux \cle{faux amis} et à la \cle{paronymie} (ex. \emph{préjudice} / \emph{préjugé}) étudiée dans cette unité.
    \item \textbf{Antonymie avec négation} : dire le contraire du contraire. \emph{« Il n'est pas sot »} $\rightarrow$ \emph{« Il est intelligent »} (attention à la nuance !) ou \emph{« Il n'ignore pas »} $\rightarrow$ \emph{« Il sait »} (litote).
    \item \textbf{Conversion grammaticale (transcatégorisation)} : nominalisation (\emph{il décide} $\rightarrow$ \emph{sa décision}), verbalisation (\emph{son arrivée} $\rightarrow$ \emph{il arrive}), adjectivalisation (\emph{la puissance} $\rightarrow$ \emph{puissant}).
    \item \textbf{Changement de voix} : actif $\leftrightarrow$ passif. \emph{« Le lion mange la perle »} $\rightarrow$ \emph{« La perle est mangée par le lion »} (utile pour varier le sujet de la phrase).
    \item \textbf{Reformulation des connecteurs logiques} : \emph{car / parce que} $\rightarrow$ \emph{puisque / comme} (cause) ; \emph{mais} $\rightarrow$ \emph{cependant / néanmoins / pourtant} (opposition) ; \emph{donc} $\rightarrow$ \emph{aussi / par conséquent / dès lors} (conséquence) ; \emph{si} $\rightarrow$ \emph{dans le cas où / à condition que} (condition).
\end{enumerate}

\courssub{Méthodologie pas à pas : de la phrase source à la phrase cible}

\begin{methode}[Reformuler efficacement en 4 étapes]
\textbf{Étape 1 : Analyser la structure.} Identifier le \cle{sujet}, le \cle{verbe}, les \cle{compléments} et le \cle{connecteur} logique. Repérer l'idée défendue dans la phrase.

\textbf{Étape 2 : Choisir les leviers.} Décider quels outils (synonymie, voix, nominalisation, connecteur) permettront de raccourcir sans perdre la nuance (ton satirique, dramatique, etc.).

\textbf{Étape 3 : Produire la phrase cible.} Écrire la nouvelle phrase en appliquant les transformations choisies. Vérifier l'accord des temps et des personnes.

\textbf{Étape 4 : Valider par la rétro-traduction.} Relire l'original et la reformulation : \textbf{le sens est-il identique ? La pensée de l'auteur est-elle respectée ?} Si oui, garder ; sinon, recommencer à l'étape 2.
\end{methode}

\begin{attention}[Piège mortel : le plagiat par « copier-coller » masqué]
Ne recopiez \textbf{jamais} un segment de plus de 3 à 4 mots consécutifs du texte original. Changer un seul mot (\emph{délaissent} $\rightarrow$ \emph{quittent}) en gardant la même syntaxe est considéré comme du plagiat à l'examen et coûte de très nombreux points. La reformulation exige une \textbf{restructuration syntaxique profonde} : changez à la fois les mots ET la construction de la phrase.
\end{attention}

\courssub{Entraînement collectif : trois phrases-clés inspirées du \emph{Lion et la Perle}}

Appliquons la méthode à des phrases représentatives des lectures méthodiques 4 et 5 (scènes de la séduction, du mime de la danse du voyageur, et du dénouement). Le texte de Soyinka n'étant pas libre de droit, nous travaillons sur des phrases \emph{reconstituées} au plus près du style et des arguments de la pièce.

\begin{exoresolu}[Reformulation de trois énoncés argumentatifs]
\textbf{Énoncé 1 (Satire de la modernité superficielle) :} \\
\emph{« Lakunle, dans son arrogance, croit que la simple imitation des manières européennes suffit à faire de lui un homme moderne. »} (22 mots)

\textbf{Énoncé 2 (Pragmatisme de Baroka / tonalité dramatique) :} \\
\emph{« Le Bale use de ruse et de tradition pour neutraliser le progrès technique qui menace son autorité ancestrale. »} (19 mots)

\textbf{Énoncé 3 (Choix final de Sidi / ironie) :} \\
\emph{« Sidi préfère l'époux qui honore la coutume à celui qui la méprise sous prétexte de civilisation. »} (18 mots)

\tcblower
\textbf{Solution détaillée :}

\textbf{1.} \emph{« L'arrogance de Lakunle le leurre : singer l'Occident ne fait pas la modernité. »} (14 mots) \\
\textit{Outils :} nominalisation (\emph{croit que} $\rightarrow$ \emph{le leurre}), compression (\emph{simple imitation des manières européennes} $\rightarrow$ \emph{singer l'Occident}), changement de structure (subordonnée $\rightarrow$ infinitive), conservation de la tonalité satirique (\emph{leurre}, \emph{singer}).

\textbf{2.} \emph{« Par ruse et tradition, Baroka contrarie le progrès technique menaçant son autorité. »} (12 mots) \\
\textit{Outils :} inversion (circonstance placée en tête), compression (\emph{use de ruse} $\rightarrow$ \emph{par ruse}), synonymie (\emph{neutraliser} $\rightarrow$ \emph{contrarier}), suppression du relatif (\emph{qui menace} $\rightarrow$ \emph{menaçant}, participe présent).

\textbf{3.} \emph{« Sidi choisit l'époux respectueux des coutumes, non celui qui les insulte au nom du progrès. »} (16 mots) \\
\textit{Outils :} antonymie (\emph{honore} / \emph{méprise} $\rightarrow$ \emph{respectueux} / \emph{insulte}), nominalisation (\emph{honore la coutume} $\rightarrow$ \emph{respectueux des coutumes}), reformulation du connecteur (\emph{sous prétexte de} $\rightarrow$ \emph{au nom de}), conservation de l'opposition binaire.
\end{exoresolu}

\courssub{Production écrite : contraction d'un texte argumentatif sur le sport scolaire au Cameroun}

Voici un texte source (environ 210 mots). L'objectif : produire un résumé de \textbf{80 mots $\pm$ 10\,\%} (soit 72 à 88 mots), en respectant la structure argumentative (thèse -- arguments -- conclusion).

\begin{exemplebox}[Texte source : « Le sport scolaire, levier oublié de l'éducation camerounaise »]
Le sport scolaire au Cameroun traverse une crise silencieuse mais profonde. Pourtant, les vertus pédagogiques de l'activité physique ne sont plus à démontrer : elle forge la discipline, l'esprit d'équipe et la résilience, qualités indispensables à tout technicien. Malheureusement, les infrastructures manquent cruellement dans la majorité des établissements publics. Les terrains sont souvent inexistants ou impraticables en saison des pluies, et le matériel vétuste décourage les initiatives. Par ailleurs, l'emploi du temps surchargé relègue l'EPS au rang d'option négligeable, voire de variable d'ajustement. Les professeurs d'EPS, peu nombreux et mal outillés, peinent à dispenser un enseignement de qualité. Certains parents d'élèves, obsédés par la seule réussite académique, voient même le sport comme une perte de temps. Face à ce constat, l'État doit réhabiliter le sport comme discipline à part entière, doter les lycées techniques de gymnases et de stades fonctionnels, et intégrer l'EPS dans les coefficients du Probatoire. Car un technicien en bonne santé, équilibré et solidaire est un atout majeur pour l'industrialisation du pays. Négliger le sport, c'est hypothéquer l'avenir de notre jeunesse technique.
\end{exemplebox}

\begin{methode}[Rédaction du résumé (80 mots) : application de la méthode]
\textbf{1. Segmentation argumentative :} thèse (crise du sport scolaire) -- argument 1 (vertus) -- argument 2 (manque d'infrastructures) -- argument 3 (place horaire et parents) -- proposition (réhabilitation, coefficients) -- conclusion (enjeu national).

\textbf{2. Sélection des idées forces (noyaux) :} crise silencieuse / vertus (discipline, équipe, résilience) / infrastructures manquantes / horaire réduit et parents hostiles / solutions (réhabilitation, coefficients) / enjeu (technicien performant).

\textbf{3. Reformulation et chaînage (connecteurs) :} \emph{Malgré ses vertus pédagogiques (discipline, esprit d'équipe), le sport scolaire camerounais est en crise : infrastructures absentes, horaire réduit, hostilité parentale. Pour former des techniciens performants, l'État doit le réhabiliter (infrastructures, coefficient au Probatoire), car négliger le sport compromet l'industrialisation.}

\textbf{4. Comptage et ajustement :} 42 mots. \emph{Trop court.} Il faut développer les arguments clés.

\textbf{5. Version finale (82 mots) :} \\
\emph{« Bien que forgeant discipline, esprit d'équipe et résilience — atouts majeurs pour le technicien —, le sport scolaire camerounais est en crise : infrastructures défaillantes, horaire marginalisé, désintérêt parental. Pour garantir des techniciens équilibrés et solidaires, moteurs de l'industrialisation, l'État doit impérativement réhabiliter l'EPS : doter les lycées techniques d'équipements fonctionnels et l'intégrer aux coefficients du Probatoire. Négliger cette discipline, c'est sacrifier l'avenir de la jeunesse technique. »}
\end{methode}

\begin{popfigure}
\centering
\begin{tikzpicture}[
    box/.style={draw=popline, rounded corners=4pt, minimum height=4.5cm, text width=6.8cm, align=justify, font=\small, inner sep=6pt},
    title/.style={font=\bfseries\small, text=white, rounded corners=4pt, minimum height=0.7cm, text width=6.8cm, align=center}
]
% Colonne gauche : texte source
\node[title, fill=popblue] (tA) at (0,0) {Texte source ($\approx$ 210 mots)};
\node[box, anchor=north, align=center] (cA) at (0,-0.45){
\textbullet\ \textbf{Constat :} crise silencieuse du sport scolaire.\\
\textbullet\ \textbf{Vertus :} discipline, équipe, résilience.\\
\textbullet\ \textbf{Causes :} infrastructures absentes, matériel vétuste.\\
\textbullet\ \textbf{Organisation :} EPS négligée, horaires surchargés.\\
\textbullet\ \textbf{Acteurs :} professeurs mal outillés ; parents hostiles.\\
\textbullet\ \textbf{Solutions :} gymnases, stades, coefficient au Probatoire.\\
\textbullet\ \textbf{Enjeu :} technicien sain = industrialisation.
};
% Colonne droite : résumé
\node[title, fill=popgreen] (tR) at (9.2,0) {Résumé ($\approx$ 80 mots)};
\node[box, anchor=north] (cR) at (9.2,-0.45) {
\emph{« Bien que forgeant discipline, esprit d'équipe et résilience, le sport scolaire camerounais est en crise : infrastructures défaillantes, horaire marginalisé, désintérêt parental. Pour garantir des techniciens équilibrés, moteurs de l'industrialisation, l'État doit réhabiliter l'EPS : équipements fonctionnels et coefficient au Probatoire. Négliger cette discipline, c'est sacrifier l'avenir de la jeunesse technique. »}
};
% Flèche centrale
\draw[->, very thick, poporange] (cA.east) -- (cR.west)
    node[midway, above, font=\small\itshape, text=popdark, align=center] {Sélection\\Reformulation\\Compression};
\end{tikzpicture}
\legende{Du texte source au résumé : sélection des noyaux argumentatifs, reformulation syntaxique et lexicale, compression quantitative (environ 40\,\% du volume initial).}
\end{popfigure}

\courssub{Vérification finale : la check-list du résumeur}

\begin{aretenir}[Check-list avant de rendre la copie]
\begin{enumerate}
    \item \textbf{Fidélité} : la thèse de l'auteur est-elle respectée ? (Pas d'inversion de sens, pas d'ajout d'opinion personnelle.)
    \item \textbf{Reformulation} : aucune séquence de 4 mots identiques au texte ? Voix, catégories grammaticales et connecteurs changés ?
    \item \textbf{Concision} : nombre de mots dans la fourchette $\pm$ 10\,\% ? (Repère : une ligne manuscrite $\approx$ 8 à 10 mots selon l'écriture.)
    \item \textbf{Cohérence} : les connecteurs logiques (\emph{bien que, car, pour, donc}) assurent-ils la fluidité ?
    \item \textbf{Langue} : orthographe, accords, ponctuation impeccables ? Pas de registre familier, pas de sigle non défini.
\end{enumerate}
\end{aretenir}

\courssub{Compétence transversale : du résumé au compte rendu professionnel}

\begin{savaistu}[La reformulation, compétence technique et administrative]
La capacité à \cle{reformuler les idées essentielles} dépasse largement l'examen de Français. Dans tous les métiers techniques (BTP, industrie, informatique, secrétariat, gestion), le \textbf{compte rendu de réunion}, la \textbf{note de synthèse}, le \textbf{rapport d'intervention} ou le \textbf{cahier de chantier} exigent exactement la même gymnastique intellectuelle : écouter ou lire une masse d'informations (souvent orales, confuses, redondantes), en extraire la substantifique moelle (décisions, actions, responsables, délais), et la restituer par écrit de façon \textbf{claire, concise, structurée et neutre}. L'élève de Première technique qui maîtrise la contraction de texte s'entraîne, sans le savoir, à la rédaction professionnelle que son futur employeur attendra de lui dès le stage ou le premier emploi. \cle{Reformuler, c'est professionnaliser sa pensée.}
\end{savaistu}

\begin{exercice}[Entraînement individuel]
À partir du texte ci-dessous (extrait d'un rapport technique fictif, environ 110 mots), rédigez un résumé de \textbf{50 mots $\pm$ 10\,\%} en appliquant la méthode en 4 étapes.

\emph{« Le système de distribution d'eau potable de la zone Nord présente des pertes physiques estimées à 35\,\% du volume produit. Ces fuites, dues à la vétusté des canalisations en fonte grise posées avant 1980, engendrent un surcoût énergétique majeur pour le pompage et une intermittence de la desserte pour les abonnés. Le remplacement par du PVC haute densité, bien que coûteux à l'investissement (2,5 milliards FCFA), garantit un retour sur investissement en 4 ans grâce aux économies d'électricité et à la réduction des interventions d'urgence. La direction recommande un phasage sur 3 ans, priorisant les tronçons à forte pression. »}
\end{exercice}

\begin{corrige}[Exercice n°15]
\textbf{Proposition de résumé (52 mots) :}

\emph{« Le réseau d'eau de la zone Nord perd 35\,\% de sa production par des fuites sur canalisations vétustes (fonte grise d'avant 1980), causant surcoût énergétique et desserte intermittente. Le renouvellement en PVC haute densité (2,5 milliards FCFA), rentable en 4 ans, est recommandé par phases sur 3 ans, en priorité sur les tronçons à forte pression. »}

\textbf{Démarche suivie :} segmentation (constat / causes / solution / recommandation), sélection des noyaux chiffrés (35\,\%, 2,5 milliards, 4 ans, 3 ans), reformulation (\emph{vétusté des canalisations} $\rightarrow$ \emph{canalisations vétustes} ; \emph{phasage} $\rightarrow$ \emph{par phases}), chaînage par la relative et l'apposition.
\end{corrige}

\coursec{Paronymie, antonymie et exposé oral}

Dans la section précédente, nous avons appris à \cle{reformuler} les idées essentielles d'un texte : dire la même chose avec d'autres mots, sans trahir la pensée de l'auteur. Mais reformuler est un exercice délicat : il suffit parfois d'une seule lettre mal placée pour changer complètement le sens d'une phrase. Écrire « le gouvernement a \emph{effectué} des réformes » au lieu de « le gouvernement a \emph{affecté} des fonds aux réformes », ce n'est pas du tout la même chose ! Cette dernière section de la leçon vous donne deux outils de vocabulaire indispensables — la \cle{paronymie} et l'\cle{antonymie} — puis vous prépare à une compétence orale capitale : l'\cle{exposé}. Nous terminerons par un atelier pratique et par le compte-rendu collectif de l'unité.

\courssub{La paronymie : des mots voisins, des sens différents}

\textbf{Intuition.} Imaginez deux personnes qui se ressemblent beaucoup physiquement mais qui n'ont ni le même nom, ni le même métier, ni le même caractère. En français, certains mots sont comme ces « faux jumeaux » : ils se prononcent presque pareil ou s'écrivent presque pareil, mais ils ne signifient pas la même chose. Ce sont les paronymes.

\begin{definition}[Les paronymes]
Les \cle{paronymes} sont des mots de \textbf{prononciation} ou d'\textbf{orthographe voisines} mais de \textbf{sens différents}. Ils ont souvent une racine commune, ce qui explique leur ressemblance, mais leurs emplois ne sont jamais interchangeables.
\end{definition}

\begin{exemplebox}[Paires de paronymes à connaître]
\begin{itemize}
\item \textbf{affecter / effectuer} : \emph{affecter} = destiner, attribuer (« On affecte des fonds à la santé ») ; \emph{effectuer} = réaliser, accomplir (« Le mécanicien effectue une réparation »).
\item \textbf{abusif / abusant} : \emph{abusif} = qui constitue un abus (« un prix abusif ») ; \emph{abusant} = qui abuse, qui trompe (terme rare et soutenu).
\item \textbf{collision / collusion} : \emph{collision} = choc brutal entre deux véhicules ; \emph{collusion} = entente secrète et malhonnête (« une collusion entre deux candidats »).
\item \textbf{éminent / imminent} : \emph{éminent} = très distingué, remarquable (« un éminent professeur ») ; \emph{imminent} = qui va arriver très vite (« un danger imminent »).
\end{itemize}
\end{exemplebox}

\begin{attention}[Ne pas confondre paronymes et synonymes]
Les \cle{synonymes} ont des sens proches et peuvent se remplacer ; les \cle{paronymes} ont des \textbf{formes} proches mais des sens différents et ne peuvent \textbf{jamais} se remplacer. « Éminent » et « imminent » ne veulent pas dire la même chose : écrire « l'examen est éminent » au lieu de « l'examen est imminent » est une faute grave qui fait perdre des points à l'examen.
\end{attention}

\textbf{Le piège de la reformulation.} Quand vous résumez un texte, vous êtes tenté de remplacer les mots de l'auteur par des mots « qui ressemblent ». C'est la source la plus fréquente de \cle{contresens}. Si le texte dit « les deux partis sont en collusion » et que vous écrivez « les deux partis sont en collision », vous avez transformé une entente secrète en accident de la route : le sens est détruit. Règle d'or : en cas de doute sur un mot, \textbf{gardez le mot du texte} ou choisissez un synonyme sûr.

\begin{exemplebox}[Dix paires de paronymes fréquents au Baccalauréat technique]
\begin{center}
\begin{tabularx}{\linewidth}{|l|X|}
\hline
\rowcolor{popblueL} \textbf{Paire} & \textbf{Phrase d'emploi} \\
\hline
affecter / effectuer & L'État affecte un budget à l'école ; l'élève effectue ses révisions. \\
\hline
collision / collusion & La collision a bloqué la route ; la collusion a faussé le concours. \\
\hline
éminent / imminent & Un éminent écrivain ; un orage imminent. \\
\hline
infecter / infester & La plaie s'infecte ; les moustiques infestent la maison. \\
\hline
éruption / irruption & L'éruption du volcan ; l'irruption des élèves dans la salle. \\
\hline
différer / diffamer & Différer une réunion (reporter) ; diffamer un collègue (nuire à sa réputation). \\
\hline
censé / sensé & Il est censé venir (supposé) ; un discours sensé (qui a du sens). \\
\hline
avarie / avarice & Une avarie de moteur (panne) ; l'avarice du personnage (cupidité). \\
\hline
inclure / incliner & Le prix inclut le transport ; il incline la tête. \\
\hline
fortuite / fortifiante & Une rencontre fortuite (par hasard) ; une lecture fortifiante (qui donne des forces). \\
\hline
\end{tabularx}
\end{center}
\end{exemplebox}

\courssub{L'antonymie : reformuler par le contraire}

\textbf{Intuition.} Pour dire « l'élève refuse de tricher », on peut aussi dire « l'élève \emph{n'accepte pas} de tricher ». On a remplacé un verbe par la \textbf{négation de son contraire}. C'est une technique de reformulation très puissante, qui repose sur l'antonymie.

\begin{definition}[L'antonymie]
L'\cle{antonymie} est la relation qui unit deux mots de \textbf{sens contraires}. Deux antonymes appartiennent à la même classe grammaticale (deux adjectifs, deux verbes, deux noms) et s'opposent sur un même domaine de sens.
\end{definition}

\textbf{Les trois types d'antonymes.} Tous les contraires ne s'opposent pas de la même manière. On distingue :

\begin{itemize}
\item les antonymes \cle{contradictoires} : il n'y a \textbf{pas de milieu possible}. On est \emph{vivant} ou \emph{mort}, \emph{présent} ou \emph{absent}. Nier l'un, c'est affirmer l'autre : « il n'est pas vivant » signifie « il est mort » ;
\item les antonymes \cle{contraires graduables} : il existe des \textbf{degrés intermédiaires}. Entre \emph{chaud} et \emph{froid}, il y a \emph{tiède}, \emph{brûlant}, \emph{glacial}. « Ce n'est pas chaud » ne veut pas dire « c'est froid » ;
\item les antonymes \cle{réciproques} : les deux mots décrivent \textbf{la même action vue de deux côtés}. \emph{Acheter} et \emph{vendre} : si le client achète, le commerçant vend. Autres exemples : prêter/emprunter, donner/recevoir, époux/épouse.
\end{itemize}

\begin{popfigure}
\begin{center}
\begin{tikzpicture}[
  grow=down,
  level distance=1.5cm,
  every node/.style={draw, rounded corners=2pt, fill=popblueL, font=\small, align=center, inner sep=3pt},
  level 1/.style={sibling distance=4.2cm},
  level 2/.style={sibling distance=2.6cm},
  edge from parent/.style={draw=popdark, -{Stealth[length=2mm]}}
]
\node[fill=poporangeL, font=\small\bfseries] {Relations sémantiques}
  child { node[align=center]{Synonymie\\ (sens proche)} }
  child { node[align=center]{Antonymie\\ (sens contraire)}
    child { node[align=center]{Contradictoires\\ vivant / mort} }
    child { node[align=center]{Graduables\\ chaud / froid} }
    child { node[align=center]{Réciproques\\ acheter / vendre} }
  }
  child { node[align=center]{Paronymie\\ (forme proche)\\ éminent / imminent} }
  child { node[align=center]{Homonymie\\ (même forme)\\ vol / vol} };
\end{tikzpicture}
\end{center}
\legende{Arbre des principales relations sémantiques entre les mots. La paronymie (formes voisines, sens différents) ne doit jamais être confondue avec la synonymie (sens proches) ni avec l'homonymie (même forme, sens différents : \emph{vol} d'un oiseau / \emph{vol} dans un magasin).}
\end{popfigure}

\textbf{Utilité pour la contraction de texte.} Reformuler par la négation de l'antonyme permet de varier les tournures sans changer le sens :

\begin{center}
\begin{tabularx}{\linewidth}{|X|X|}
\hline
\rowcolor{popblueL} \textbf{Phrase du texte} & \textbf{Reformulation par antonymie} \\
\hline
Il refuse de participer. & Il n'accepte pas de participer. \\
\hline
Le projet a réussi. & Le projet n'a pas échoué. \\
\hline
La salle était pleine. & La salle n'était pas vide. \\
\hline
Il se souvient de tout. & Il n'oublie rien. \\
\hline
\end{tabularx}
\end{center}

\begin{attention}[Graduables : la négation ne suffit pas toujours]
Avec les antonymes \textbf{graduables}, la négation du contraire n'est pas équivalente : « le texte n'est pas mauvais » ne signifie pas « le texte est bon ». Réservez la reformulation par négation aux antonymes \textbf{contradictoires}, où l'équivalence est totale.
\end{attention}

\begin{exoresolu}[Reformuler sans contresens]
Énoncé : reformulez chaque phrase en utilisant la négation d'un antonyme, puis indiquez le type d'antonyme utilisé. a) « Le malade est vivant. » b) « Le candidat a été admis au concours. » c) « Le marchand vend ses produits. »
\tcblower
Solution détaillée :
\begin{itemize}
\item a) « Le malade n'est pas mort. » Antonymes \textbf{contradictoires} (vivant/mort) : l'équivalence est parfaite, il n'y a pas d'état intermédiaire.
\item b) « Le candidat n'a pas été recalé (refusé) au concours. » Antonymes \textbf{contradictoires} (admis/recalé) : reformulation correcte.
\item c) « Le marchand ne garde pas ses produits » serait un contresens ! Avec les antonymes \textbf{réciproques} (vendre/acheter), la négation ne fonctionne pas : on reformule plutôt en changeant de point de vue : « Les clients achètent leurs produits au marchand. »
\end{itemize}
\end{exoresolu}

\courssub{L'exposé : préparation et présentation}

\textbf{Intuition.} Vous savez déjà faire un exposé sans le savoir : chaque fois que vous racontez un match de football ou expliquez une règle de jeu à des amis, vous organisez des idées, vous parlez debout devant un groupe, vous captez l'attention. L'exposé scolaire repose sur les mêmes gestes, mais préparés méthodiquement.

\begin{definition}[L'exposé]
L'\cle{exposé} est une prise de parole \textbf{organisée} et \textbf{limitée dans le temps}, destinée à informer ou à convaincre un auditoire sur un sujet précis. Il comprend trois moments : l'\textbf{introduction} (accroche et annonce du plan), le \textbf{développement} (les parties annoncées) et la \textbf{conclusion} (bilan et ouverture).
\end{definition}

\textbf{La préparation : les trois outils de l'orateur.}
\begin{enumerate}
\item \textbf{Le plan} : après la recherche documentaire, rédigez un plan en deux ou trois parties équilibrées. Chaque partie contient une idée principale et un exemple tiré de l'œuvre étudiée.
\item \textbf{Les fiches} : ne rédigez jamais votre exposé en entier ! Notez sur de petites fiches numérotées les \textbf{mots-clés}, les citations courtes et les exemples. Les fiches évitent la lecture monotone.
\item \textbf{Les supports} : un tableau, un schéma ou une citation écrite au tableau peuvent appuyer votre propos et aider l'auditoire à suivre.
\end{enumerate}

\textbf{La présentation : le triangle de la communication.} Toute communication orale repose sur trois pôles : l'\textbf{orateur} (celui qui parle), le \textbf{message} (ce qu'il dit) et l'\textbf{auditoire} (ceux qui écoutent). Un exposé réussi maintient l'équilibre entre ces trois pôles : l'orateur adapte son message à l'auditoire, et l'auditoire renvoie des signes (regards, hochements de tête) que l'orateur doit observer.

\begin{popfigure}
\begin{center}
\begin{tikzpicture}[scale=1.1]
\coordinate (O) at (0,3.2);
\coordinate (M) at (-3,0);
\coordinate (A) at (3,0);
\draw[very thick, popblue] (O) -- (M) -- (A) -- cycle;
\node[draw, fill=poporangeL, rounded corners=3pt, font=\small\bfseries, align=center] at (O) {Orateur\\ \scriptsize posture, voix, regard};
\node[draw, fill=popgreenL, rounded corners=3pt, font=\small\bfseries, align=center] at (M) {Message\\ \scriptsize plan, idées, exemples};
\node[draw, fill=poppinkL, rounded corners=3pt, font=\small\bfseries, align=center] at (A) {Auditoire\\ \scriptsize écoute, réactions};
\draw[-{Stealth[length=2.5mm]}, thick, poppurple] (-1.9,0.9) -- (-0.9,2.1) node[midway, left, font=\scriptsize, align=center]{adapter\\ le propos};
\draw[-{Stealth[length=2.5mm]}, thick, poppurple] (1.9,0.9) -- (0.9,2.1) node[midway, right, font=\scriptsize, align=center]{capter\\ l'attention};
\draw[{Stealth[length=2.5mm]}-{Stealth[length=2.5mm]}, thick, poppurple] (-2.2,-0.45) -- (2.2,-0.45) node[midway, below, font=\scriptsize]{réactions et feed-back};
\end{tikzpicture}
\end{center}
\legende{Le triangle de la communication orale : l'orateur adapte son message à l'auditoire et observe ses réactions pour ajuster son débit, sa voix et ses explications.}
\end{popfigure}

\begin{methode}[Pour un exposé réussi]
\begin{enumerate}
\item \textbf{Annoncer le plan} dès l'introduction : « Nous verrons d'abord..., ensuite..., enfin... ». L'auditoire sait où vous l'emmenez.
\item \textbf{Parler lentement} et articuler : le stress accélère le débit ; imposez-vous des pauses après chaque idée importante.
\item \textbf{Regarder l'auditoire} : balayez la salle du regard, ne fixez ni vos fiches, ni le plafond, ni le correcteur.
\item \textbf{Soigner la posture} : debout, droit, les deux pieds au sol ; les gestes accompagnent la parole sans être excessifs.
\item \textbf{Gérer le temps} : répétez avec un chronomètre ; prévoyez ce que vous pouvez raccourcir si vous êtes en retard.
\item \textbf{Conclure en rappelant l'essentiel} : reformulez en une phrase l'idée principale, puis remerciez l'auditoire et proposez des questions.
\end{enumerate}
\end{methode}

\begin{savaistu}[L'exposé, un entraînement pour la vie]
L'exposé oral prépare directement à l'\textbf{épreuve orale du Baccalauréat technique}, mais aussi aux \textbf{soutenances professionnelles} : présenter un projet à un chef d'atelier, défendre un devis devant un client, expliquer une panne à un supérieur... Dans la vie active au Cameroun comme ailleurs, celui qui sait parler clairement devant un groupe gagne la confiance des autres. Chaque exposé en classe de Première est donc une répétition de votre avenir professionnel.
\end{savaistu}

\courssub{Atelier : exposé de 5 minutes sur la satire dans \emph{Le Lion et la perle}}

\begin{experience}[Atelier en binômes : préparer et présenter un exposé]
\textbf{Matériel} : l'œuvre \emph{Le Lion et la perle} de Wole Soyinka, des fiches bristol, un chronomètre, la grille d'évaluation de l'exposé.\par
\textbf{Protocole} :
\begin{enumerate}
\item Par binôme, relisez les passages satiriques étudiés lors des lectures méthodiques : le personnage de Lakunle (l'intellectuel ridiculisé), les coutumes moquées, les jeux de langage comiques.
\item Établissez un plan en deux parties, par exemple : I. Les cibles de la satire (personnages et comportements moqués) ; II. Les procédés satiriques (exagération, ironie, caricature).
\item Rédigez vos fiches : mots-clés, deux citations courtes, un exemple par partie.
\item Répartissez la parole : chaque membre du binôme présente une partie ; l'un introduit, l'autre conclut.
\item Répétez avec le chronomètre : 5 minutes maximum.
\end{enumerate}
\textbf{Observations attendues} : pendant les passages à l'oral, les autres binômes remplissent la grille d'évaluation (respect du plan, clarté de la voix, regard, gestion du temps, pertinence des exemples).\par
\textbf{Interprétation} : la comparaison des grilles permet à chaque binôme d'identifier ses points forts et ses axes de progrès avant l'examen.
\end{experience}

\begin{exercice}[Entraînement au vocabulaire]
\begin{enumerate}
\item Choisissez le bon paronyme : a) Le chauffeur a (effectué / affecté) une réparation. b) Une (collision / collusion) entre deux camions a bloqué l'axe Yaoundé--Douala. c) Le départ est (éminent / imminent).
\item Reformulez par la négation d'un antonyme contradictoire : a) « Le texte est cohérent. » b) « La réponse est exacte. »
\item Classez les paires suivantes selon le type d'antonymie : grand/petit ; légal/illégal ; enseigner/apprendre.
\end{enumerate}
\end{exercice}

\begin{corrige}[Exercice : entraînement au vocabulaire]
\begin{enumerate}
\item a) \textbf{effectué} (réaliser) ; b) \textbf{collision} (choc) ; c) \textbf{imminent} (qui va arriver).
\item a) « Le texte n'est pas incohérent. » b) « La réponse n'est pas fausse (inexacte). »
\item grand/petit : contraires \textbf{graduables} (il existe des tailles intermédiaires) ; légal/illégal : \textbf{contradictoires} (pas de milieu) ; enseigner/apprendre : \textbf{réciproques} (la même situation vue des deux côtés).
\end{enumerate}
\end{corrige}

\courssub{Compte-rendu et remédiation}

\textbf{Correction collective du résumé type.} Le professeur projette ou dicte le résumé type du texte travaillé dans l'unité. Chaque élève compare sa production au modèle en suivant trois questions : ai-je conservé \textbf{toutes les idées essentielles} ? ai-je respecté la \textbf{structure argumentative} (thèse, arguments, conclusion) ? ai-je évité les \textbf{contresens} liés aux paronymes et les citations trop longues ? Les écarts sont notés en marge, puis les formulations sont améliorées collectivement au tableau.

\begin{aretenir}[Grille d'auto-évaluation de l'unité]
\begin{center}
\begin{tabularx}{\linewidth}{|X|c|c|c|}
\hline
\rowcolor{popblueL} \textbf{Compétence} & \textbf{Acquis} & \textbf{En cours} & \textbf{À revoir} \\
\hline
Identifier la thèse et les arguments d'un texte & & & \\
\hline
Supprimer les éléments secondaires (exemples, répétitions) & & & \\
\hline
Reformuler sans contresens (paronymes maîtrisés) & & & \\
\hline
Utiliser l'antonymie pour reformuler par la négation & & & \\
\hline
Préparer et présenter un exposé en respectant le temps & & & \\
\hline
\end{tabularx}
\end{center}
Cochez chaque case honnêtement : toute compétence « à revoir » fait l'objet d'une remédiation individuelle (exercices supplémentaires, relecture de la leçon, entretien avec le professeur) avant l'évaluation sommative.
\end{aretenir}

\begin{attention}[Les trois erreurs qui coûtent cher à l'examen]
\begin{enumerate}
\item Confondre un paronyme avec un synonyme et produire un contresens dans le résumé ;
\item Recopier des phrases entières du texte au lieu de reformuler avec ses propres mots ;
\item Lire ses fiches pendant l'exposé au lieu de regarder l'auditoire et de gérer son temps.
\end{enumerate}
\end{attention}

Vous disposez désormais de tous les outils de cette unité : repérer la structure d'un texte argumentatif, le réduire, reformuler ses idées essentielles avec un vocabulaire précis, et restituer votre travail à l'oral. Ces compétences, entraînées sur \emph{Le Lion et la perle}, seront mobilisées dans toutes les épreuves de Français du Probatoire et du Baccalauréat technique — et bien au-delà, dans votre vie professionnelle.

\newpage

\coursec{Activité d'intégration}

\courssub{Objectifs de l'activité}

À l'issue de cette activité d'intégration, l'élève sera capable de :
\begin{itemize}
\item repérer dans un texte argumentatif de presse la \cle{thèse}, les \cle{arguments} et les \cle{exemples} ;
\item identifier la \cle{tonalité} dominante d'un éditorial (comique, dramatique, satirique) et la justifier par des indices précis ;
\item produire un \cle{résumé} de 100 mots exactement en appliquant les quatre techniques de réduction (suppression, généralisation, condensation, substitution) et en reformulant les idées essentielles avec ses propres mots ;
\item préparer et présenter un mini-\cle{exposé} oral de 3 minutes ;
\item évaluer le travail d'un autre groupe à l'aide d'une grille critériée et participer à une remédiation collective.
\end{itemize}

Cette activité mobilise les savoir-faire officiels : appliquer les notions de l'unité à une situation concrète de la vie courante, et rédiger puis justifier une démarche, une réponse ou une conclusion.

\courssub{Situation}

Le club de presse du lycée technique de Nkol-Eton (Yaoundé) prépare son émission hebdomadaire « La parole aux jeunes », diffusée lors de la pause méridienne. Le thème retenu cette semaine est un débat qui agite la capitale : \emph{Faut-il interdire les motos-taxis dans le centre-ville de Yaoundé ?} Chaque groupe d'élèves doit résumer l'éditorial ci-dessous, publié dans un quotidien camerounais, puis le présenter oralement. Le résumé retenu sera lu au micro de l'émission.

\textbf{Texte de départ (400 mots) :}

\begin{quote}
\textbf{Les motos-taxis, ce fléau à deux roues}

Voilà des années que le centre-ville de Yaoundé est devenu un théâtre permanent de l'anarchie. Dès six heures du matin, des centaines de motos-taxis envahissent les rues, slaloment entre les voitures, grillent les feux rouges et transforment nos trottoirs en pistes de course. Le spectacle est quotidien : klaxons assourdissants, pots d'échappement crachant une fumée noire, conducteurs casqués qui hèlent les passants à tout-va. Notre belle capitale mérite-t-elle vraiment ce chaos ?

Les partisans des motos-taxis avancent un argument qui n'est pas négligeable : ces engins créent des milliers d'emplois pour une jeunesse désœuvrée. À Yaoundé, on estime à près de vingt mille le nombre de jeunes qui vivent de ce métier. Supprimer les motos, disent-ils, c'est jeter ces jeunes dans la délinquance. L'argument est séduisant, mais il ne résiste pas à l'examen.

D'abord, l'insécurité routière causée par les motos-taxis est effrayante. Selon le service des urgences de l'hôpital central, un accident de la circulation sur trois dans le centre-ville implique une moto. Bras cassés, crânes fracassés, familles endeuillées : voilà le prix que nous payons chaque semaine pour cette prétendue « solution » au chômage. Aucun emploi ne justifie un tel bilan humain.

Ensuite, l'image de la ville en pâtit gravement. Quel investisseur étranger, quel touriste, peut se sentir à l'aise dans un centre-ville où régnerait la loi du plus rapide ? Douala, notre rival économique, a su réglementer sévèrement la circulation des motos dans son centre des affaires, et le résultat est visible : des rues plus calmes, des commerces plus attractifs. Yaoundé, capitale politique, ne peut pas rester la ville du désordre.

Certes, interdire ne suffit pas. Il faut accompagner la mesure : créer des navettes de bus modernes, réserver des zones de stationnement aux motos à la périphérie, former les jeunes conducteurs à d'autres métiers. Une interdiction brutale et sans alternative serait injuste et inefficace.

En définitive, l'interdiction des motos-taxis dans le centre-ville de Yaoundé s'impose comme une nécessité de salubrité publique, à condition qu'elle s'accompagne de mesures sociales sérieuses. Le gouvernement municipal doit cesser de tergiverser : la sécurité et la dignité des habitants de la capitale ne sont pas négociables. Il est temps de rendre nos rues aux piétons.
\end{quote}

\textbf{Données de la tâche :} texte de départ : 400 mots ; résumé attendu : \cle{100 mots exactement} (soit le quart du texte) ; groupes de 4 élèves ; exposé oral : 3 minutes par porte-parole ; évaluation croisée sur grille de 20 points.

\courssub{Consignes}

\begin{enumerate}
\item Lisez silencieusement l'éditorial, puis dégagez le \cle{thème} et la \cle{problématique} du texte en une phrase chacun.
\item Relevez la \cle{thèse} de l'éditorialiste : où est-elle énoncée ? Recopiez la phrase qui l'exprime le plus clairement.
\item Construisez le tableau de la structure argumentative : pour chaque argument, indiquez l'exemple ou la donnée chiffrée qui l'illustre. Repérez aussi l'argument adverse (la \cle{concession}) et la manière dont l'auteur le réfute.
\item Quelle est la \cle{tonalité} dominante de l'éditorial ? Justifiez votre réponse par trois indices précis (champ lexical, procédé de style, marque d'énonciation). Relevez au passage deux paires de \cle{paronymes} ou d'\cle{antonymes} présents dans le texte.
\item Rédigez en groupe un résumé de \textbf{100 mots exactement}, en appliquant les quatre techniques de réduction et en reformulant toutes les idées : aucune phrase ne doit être recopiée du texte. Comptez les mots et inscrivez le total à la fin.
\item Désignez un porte-parole et préparez le mini-exposé de 3 minutes : plan d'annonce, présentation du thème, de la thèse, des arguments et de la conclusion du résumé.
\item Pendant les exposés des autres groupes, remplissez la grille d'évaluation croisée, puis participez à la remédiation collective sur les erreurs relevées.
\end{enumerate}

\courssub{Ressources à mobiliser}

\begin{aretenir}[Outils de la leçon à réactiver]
\begin{itemize}
\item \textbf{Structure argumentative :} thèse (idée défendue), arguments (raisons), exemples (preuves concrètes), concession-réfutation, connecteurs logiques (\emph{d'abord, ensuite, certes, en définitive}).
\item \textbf{Quatre techniques de réduction :} suppression (éléments secondaires), généralisation (terme générique pour une énumération), condensation (une phrase pour deux), substitution (un mot pour un groupe de mots).
\item \textbf{Règles du résumé :} respecter la proportion (ici 1/4), reformuler sans recopier, rester fidèle à la pensée de l'auteur, ne pas ajouter d'opinion personnelle, compter les mots.
\item \textbf{Tonalités :} comique (rire), dramatique (émotion, gravité), satirique (dénonciation par l'ironie et l'exagération).
\item \textbf{Lexique :} paronymie (mots proches par la forme : \emph{effrayant/effarant}), antonymie (mots de sens contraire : \emph{ordre/désordre}).
\item \textbf{Exposé oral :} introduction (accroche, sujet, plan), développement, conclusion ; voix claire, regard vers l'auditoire, gestion du temps.
\end{itemize}
\end{aretenir}

\begin{attention}[Pièges fréquents au Probatoire]
\begin{itemize}
\item Recopier des phrases entières du texte : le résumé doit être une \emph{reformulation} personnelle.
\item Dépasser ou sous-estimer le nombre de mots imposé : comptez mot par mot et inscrivez le total.
\item Ajouter son propre avis (« je pense que... ») : le résumé restitue la pensée de l'auteur, pas la vôtre.
\item Confondre des paronymes (\emph{effrayant/effarant}, \emph{interdiction/intercession}) : vérifiez chaque mot au dictionnaire mentalement.
\end{itemize}
\end{attention}

\courssub{Démarche et corrigé commenté}

\begin{exoresolu}[Résolution guidée de la tâche]
\textbf{Étape 1 — Thème et problématique.} Le thème est la circulation des motos-taxis dans le centre-ville de Yaoundé. La problématique peut se formuler ainsi : \emph{les motos-taxis doivent-ils être interdits dans le centre-ville de Yaoundé, et à quelles conditions ?}

\tcblower

\textbf{Étape 2 — La thèse.} Elle est énoncée explicitement dans le dernier paragraphe : « l'interdiction des motos-taxis dans le centre-ville de Yaoundé s'impose comme une nécessité de salubrité publique, à condition qu'elle s'accompagne de mesures sociales sérieuses ». La thèse est donc \emph{nuancée} : oui à l'interdiction, mais avec accompagnement social.

\textbf{Étape 3 — Structure argumentative.}

\begin{center}
\begin{tabularx}{\linewidth}{|l|X|X|}\hline
\rowcolor{popblueL} \textbf{Élément} & \textbf{Contenu} & \textbf{Exemple / preuve} \\ \hline
Concession & Les motos créent des emplois & 20 000 jeunes vivent de ce métier \\ \hline
Réfutation & L'argument « ne résiste pas à l'examen » & Aucun emploi ne justifie un tel bilan humain \\ \hline
Argument 1 & Insécurité routière effrayante & 1 accident sur 3 implique une moto (hôpital central) \\ \hline
Argument 2 & Image dégradée de la ville & Douala a réglementé : rues plus calmes \\ \hline
Condition & Accompagner l'interdiction & Navettes de bus, parkings en périphérie, reconversion \\ \hline
\end{tabularx}
\end{center}

\textbf{Étape 4 — Tonalité et lexique.} La tonalité dominante est \cle{satirique} avec des accents dramatiques : champ lexical de la dénonciation (« fléau », « anarchie », « chaos »), exagération ironique (« pistes de course », « la loi du plus rapide »), question rhétorique (« Notre belle capitale mérite-t-elle vraiment ce chaos ? ») ; le passage sur les « bras cassés, crânes fracassés » relève du dramatique. Paronymes à distinguer : \emph{effrayant/effarant} ; \emph{interdiction/intercession}. Antonymes présents dans le texte : \emph{ordre/désordre}, \emph{calme/chaos}.

\textbf{Étape 5 — Modèle de résumé (100 mots).}

\begin{quote}
L'éditorialiste dénonce l'anarchie créée par les motos-taxis au centre de Yaoundé. À l'argument de l'emploi des jeunes, il oppose d'abord l'insécurité routière : un accident sur trois y implique une moto. Il fait ensuite valoir que ce désordre nuit à l'image de la capitale, contrairement à Douala qui a réglementé la circulation. Il préconise donc l'interdiction des motos-taxis au centre-ville, mais accompagnée de mesures sociales : transports en commun, stationnements en périphérie et reconversion des conducteurs. Il conclut que la municipalité doit agir vite pour rendre les rues aux piétons. (100 mots)
\end{quote}

Techniques utilisées : suppression des détails pittoresques (klaxons, fumée), généralisation (« mesures sociales » pour la liste), condensation des deux arguments en deux phrases, substitution (« ce désordre » pour toute la description).

\textbf{Étape 6 — Exposé oral.} Plan du porte-parole : accroche (le débat actuel), présentation du texte, annonce du plan (thèse, arguments, tonalité), lecture commentée du résumé, conclusion.

\textbf{Étape 7 — Grille d'évaluation croisée (20 points).}

\begin{center}
\begin{tabularx}{\linewidth}{|l|X|c|}\hline
\rowcolor{popblueL} \textbf{Critère} & \textbf{Attendus} & \textbf{Points} \\ \hline
Fidélité & Idées essentielles conservées, aucune opinion ajoutée & /5 \\ \hline
Réduction & 100 mots exactement, proportion respectée & /5 \\ \hline
Reformulation & Aucune phrase recopiée, vocabulaire correct (pas de confusion de paronymes) & /5 \\ \hline
Qualité orale & Voix, regard, plan respecté, 3 minutes & /5 \\ \hline
\end{tabularx}
\end{center}

Remédiation collective : corriger au tableau les résumés qui recopient des phrases, dépassent 100 mots ou confondent des paronymes.
\end{exoresolu}

\courssub{Lien avec la lecture méthodique du \emph{Lion et la perle}}

Les lectures méthodiques n\textsuperscript{o} 4 et 5 de l'œuvre au programme entretiennent un rapport direct avec cette activité. Comme dans l'éditorial, il s'agit d'abord de \emph{suivre un échange argumenté} : dans la pièce de Wole Soyinka, les personnages défendent des positions opposées — Baroka, le Lion, incarne la tradition et la ruse, tandis que Lakunle défend une modernité scolaire et bavarde. Identifier qui dit quoi, pour quelle raison et avec quelles preuves, c'est exactement repérer thèse, arguments et exemples. Ensuite, la pièce joue sur les \cle{tonalités} : le comique naît des excès de langue de Lakunle et des situations de dupe ; le dramatique affleure dans le sort de Sidi, la « perle » du village ; le satirique perce dans la critique des prétentions modernistes comme des manipulations traditionnelles. Enfin, raconter en quelques phrases ce qui s'est passé dans une scène, sans recopier les répliques, c'est déjà pratiquer la \cle{reformulation des idées essentielles} : la contraction de texte prolonge naturellement la lecture méthodique.

\courssub{Bilan de l'activité}

Cette activité a permis de vérifier que les savoirs de l'unité fonctionnent ensemble dans une tâche authentique. L'élève a d'abord \emph{lu} un texte argumentatif de presse pour en dégager la structure (thèse, concession, arguments, exemples), compétence directement liée aux lectures méthodiques du \emph{Lion et la perle}, où il s'agit aussi de suivre les positions et les émotions des personnages. Il a ensuite \emph{réduit} le texte au quart de sa longueur en distinguant l'essentiel de l'accessoire, puis \emph{reformulé} sans recopier, ce qui exige un vocabulaire précis — d'où l'importance de la paronymie et de l'antonymie. Enfin, la présentation orale a montré qu'un bon résumé ne vaut que s'il est clairement communiqué : l'exposé prolonge l'écrit. L'évaluation croisée et la remédiation ont mis en évidence trois erreurs fréquentes à éviter au Probatoire : recopier le texte au lieu de reformuler, dépasser le nombre de mots imposé, et confondre des mots de forme proche. Retenons qu'un résumé réussi est fidèle, bref, personnel et lisible à l'oral.

\newpage

\coursec{Méthodes et exercices résolus}

\courssub{Méthodes et exercices résolus}

\begin{methode}[Identifier les caractéristiques du texte argumentatif]
\textbf{Objectif :} Repérer les marques de l'argumentation pour comprendre la visée du texte.
\begin{enumerate}
    \item \textbf{Trouver la thèse :} Chercher l'idée principale défendue (souvent en début ou en fin de texte, ou après un connecteur logique comme \emph{donc}, \emph{en conclusion}).
    \item \textbf{Repérer les arguments :} Identifier les idées qui soutiennent la thèse (souvent introduites par \emph{car}, \emph{puisque}, \emph{en effet}, \emph{premièrement}).
    \item \textbf{Analyser les connecteurs logiques :} Souligner les liens (cause, conséquence, opposition, concession) qui structurent le raisonnement.
    \item \textbf{Identifier les procédés d'ancrage énonciatif :} Repérer les marques d'énonciation (\emph{je}, \emph{nous}, présent de l'énonciation) et les modalisateurs (adverbes : \emph{certainement} ; verbes : \emph{devoir}, \emph{falloir} ; adjectifs : \emph{indéniable}) qui montrent l'engagement de l'auteur.
    \item \textbf{Classer le type d'argumentation :} Déductive (thèse $\to$ arguments), inductive (arguments $\to$ thèse), dialectique (thèse / antithèse / synthèse).
\end{enumerate}
\cle{Thèse} \cle{Argument} \cle{Connecteur logique} \cle{Modalisateur} \cle{Ancrage énonciatif}
\end{methode}

\begin{exoresolu}[Analyse d'un paragraphe argumentatif]
\textbf{Énoncé :} Lisez l'extrait suivant et répondez aux questions.
\begin{quote}
\emph{« L'école technique ne doit pas être considérée comme une voie de garage. En effet, elle forme des techniciens compétents, indispensables au développement industriel du pays. De plus, les débouchés y sont réels et variés. Certes, certains préjugés persistent, mais ils s'effondrent face à la réalité du terrain. Par conséquent, valoriser cette filière est un impératif national. »}
\end{quote}
\begin{enumerate}
    \item Quelle est la thèse défendue ?
    \item Relevez deux arguments et les connecteurs qui les introduisent.
    \item Identifiez un modalisateur et un procédé d'ancrage énonciatif.
    \item Quel est le type de structure argumentative ?
\end{enumerate}
\tcblower
\textbf{Solution détaillée :}
\begin{enumerate}
    \item \textbf{Thèse :} « L'école technique ne doit pas être considérée comme une voie de garage » (phrase 1), reformulée en conclusion : « Valoriser cette filière est un impératif national ». C'est l'opinion que l'auteur veut faire admettre.
    \item \textbf{Arguments :}
        \begin{itemize}
            \item Argument 1 : « elle forme des techniciens compétents, indispensables au développement industriel du pays ». Connecteur : \textbf{« En effet »} (explication, justification).
            \item Argument 2 : « les débouchés y sont réels et variés ». Connecteur : \textbf{« De plus »} (addition, renforcement).
        \end{itemize}
    \item \textbf{Modalisateur :} « \textbf{impératif} » (adjectif exprimant la nécessité, l'obligation forte) ou « \textbf{doit} » (verbe modal de valeur déontique). \textbf{Ancrage énonciatif :} l'énonciation est ici \textbf{dépersonnalisée} (ni \emph{je} ni \emph{nous} explicites) : l'auteur objectivise son propos par l'emploi du \textbf{présent de vérité générale} (« forme », « sont », « persistent », « s'effondrent »), ce qui donne à son opinion une valeur de fait.
    \item \textbf{Structure :} \textbf{déductive} (thèse annoncée d'emblée $\to$ arguments $\to$ conclusion qui reformule la thèse). On note une concession (« Certes... ») immédiatement réfutée (« mais... »), procédé qui renforce la thèse en écartant l'objection adverse.
\end{enumerate}
\textbf{Conclusion :} Ce paragraphe est un texte argumentatif typique à structure déductive, utilisant l'addition d'arguments et la concession réfutée pour imposer la thèse.
\end{exoresolu}

\begin{methode}[Appliquer les techniques de réduction (contraction de texte)]
\textbf{Objectif :} Réduire un texte au tiers ou au quart de sa longueur tout en conservant les idées essentielles et la logique argumentative.
\begin{enumerate}
    \item \textbf{Lecture active :} Lire deux fois. Première lecture : compréhension globale. Deuxième lecture : découpage en séquences logiques (idées forces).
    \item \textbf{Supprimer le superflu :} Biffer les exemples, les détails, les illustrations, les répétitions, les périphrases, les adjectifs qualificatifs non essentiels, les connecteurs redondants.
    \item \textbf{Regrouper et généraliser :} Remplacer une énumération par un hyperonyme (ex. : \emph{chiens, chats, oiseaux} $\to$ \emph{animaux domestiques}). Fusionner les phrases voisines.
    \item \textbf{Reformuler (ne jamais copier-coller) :} Changer le vocabulaire (synonymes), la syntaxe (passer du discours direct au discours indirect, nominaliser des verbes, subordonner). Conserver les termes techniques incontournables.
    \item \textbf{Respecter la structure :} Conserver l'enchaînement thèse / arguments / conclusion. Garder les connecteurs logiques nécessaires à la cohérence.
    \item \textbf{Contrôler le volume :} Compter les mots (texte original et contraction). Vérifier l'exactitude du sens et la correction linguistique.
\end{enumerate}
\cle{Suppression} \cle{Généralisation} \cle{Reformulation} \cle{Nominalisation} \cle{Subordination}
\end{methode}

\begin{exoresolu}[Contraction d'un texte argumentatif (cible : le tiers)]
\textbf{Énoncé :} Contractez le texte suivant en conservant les idées essentielles. Le texte fait 95 mots. Visez 30 à 35 mots.
\begin{quote}
\emph{« Le développement durable est un concept qui répond aux besoins du présent sans compromettre la capacité des générations futures à répondre aux leurs. Il repose sur trois piliers indissociables : l'efficacité économique, l'équité sociale et la préservation de l'environnement. Par exemple, une entreprise qui pollue pour maximiser son profit sacrifie l'avenir. Au contraire, investir dans les énergies propres crée des emplois verts et protège la planète. Ainsi, ce modèle est la seule voie viable pour l'humanité. »}
\end{quote}
\tcblower
\textbf{Solution détaillée :}

\textbf{Étape 1 : Découpage et sélection des idées forces.}
\begin{itemize}
    \item Idée 1 (définition) : le développement durable répond aux besoins du présent sans compromettre l'avenir.
    \item Idée 2 (piliers) : trois piliers : économie, social, environnement.
    \item Idée 3 (exemple et contraste) : polluer pour le profit, c'est sacrifier l'avenir ; les énergies propres créent des emplois et protègent la planète.
    \item Idée 4 (conclusion) : c'est la seule voie viable.
\end{itemize}
\textbf{Étape 2 : Suppression et regroupement.}
On supprime « concept qui », « indissociables », l'exemple développé (l'entreprise), « Au contraire », « Ainsi ». L'idée 3 est réduite à sa seule conclusion positive (les énergies propres).
\textbf{Étape 3 : Reformulation (nominalisation et subordination).}
\begin{quote}
\emph{« Le développement durable satisfait les besoins actuels sans hypothéquer l'avenir, en conciliant efficacité économique, équité sociale et protection de l'environnement ; l'investissement dans les énergies propres, créateur d'emplois verts, en fait la seule voie viable pour l'humanité. »}
\end{quote}
\textbf{Comptage :} 38 mots (légèrement au-dessus de la cible, acceptable si la consigne tolère une marge ; on peut encore supprimer « pour l'humanité » pour atteindre 35 mots). Le sens est respecté, la structure conservée (définition $\to$ piliers $\to$ solution $\to$ conclusion).
\end{exoresolu}

\begin{methode}[Réaliser une lecture méthodique : \emph{Le Lion et la Perle} (Wole Soyinka) -- lectures n° 4 et 5]
\textbf{Objectif :} Analyser un extrait théâtral pour en dégager le sens, les enjeux dramatiques et les procédés d'écriture.
\begin{enumerate}
    \item \textbf{Contextualiser :} Situer l'extrait dans la pièce (moment de l'intrigue : exposition, nœud, péripéties, dénouement ; la pièce est organisée en trois parties : \emph{Matin}, \emph{Midi}, \emph{Nuit}). Identifier les personnages présents.
    \item \textbf{Analyser la situation de communication :} Qui parle ? À qui ? Dans quel but ? (séduction, manipulation, confrontation, confidence).
    \item \textbf{Étudier le registre tonal :} Repérer les marques du \textbf{comique} (quiproquo, jeu de mots, situation burlesque, caricature), du \textbf{dramatique} (menace, tension, destin funeste, pathos) ou du \textbf{satirique} (ironie, critique des mœurs, opposition tradition / modernité, ridicule des puissants).
    \item \textbf{Analyser le langage théâtral :} didascalies (gestes, ton, déplacements), niveaux de langue (soutenu, familier), figures de style (métaphores, proverbes, images), rythme des échanges (stichomythie, monologue).
    \item \textbf{Dégager la portée :} Quel thème majeur est illustré ? (choc des cultures, condition féminine, pouvoir, tradition et modernité). Quelle évolution pour les personnages ?
\end{enumerate}
\cle{Didascalie} \cle{Stichomythie} \cle{Registre} \cle{Portée} \cle{Situation de communication}
\end{methode}

\begin{exoresolu}[Analyse d'un extrait du \emph{Lion et la Perle} (lecture n° 4 : la scène du magazine)]
\textbf{Énoncé :} Étudiez l'extrait suivant, inspiré de la partie \emph{Midi} (Lakunlé montre à Sidi le magazine qui publie ses photos).
\begin{quote}
\emph{LAKUNLÉ : (Feuilletant le magazine) Regarde ! Voici la perle du village. La première page... la deuxième... la troisième... Tout le magazine n'est qu'elle ! \\
SIDI : (Rayonnante) C'est moi ? Vraiment ? Donne-moi ! ... Oh ! Lakunlé, je suis belle ! Regarde mes yeux, ma bouche... Le Bale lui-même n'a pas tant de place ! \\
LAKUNLÉ : (Amer) Le Bale... Le Bale a trois lignes en page six. Toi, tu as le magazine entier. C'est la modernité, Sidi. L'image prime sur le pouvoir réel.}
\end{quote}
\begin{enumerate}
    \item Situez la scène et identifiez le registre dominant.
    \item Montrez comment le personnage de Sidi est caractérisé ici.
    \item Analysez la vision de Lakunlé sur la modernité (ironie, satire).
    \item Quel thème majeur émerge ?
\end{enumerate}
\tcblower
\textbf{Solution détaillée :}
\begin{enumerate}
    \item \textbf{Contextualisation :} Dans la partie \emph{Midi}, Lakunlé, l'instituteur moderniste, montre à Sidi le magazine étranger qui consacre un reportage photographique à sa beauté. \textbf{Registre dominant :} \textbf{satirique} (critique de la vanité et du pouvoir de l'image médiatique), teinté de \textbf{comique} (naïveté de Sidi, amertume de Lakunlé).
    \item \textbf{Caractérisation de Sidi :} \textbf{vanité, narcissisme.} La didascalie « rayonnante », les exclamations (« C'est moi ? », « je suis belle ! »), la focalisation sur son physique (« mes yeux, ma bouche ») la montrent éblouie par sa propre image. Elle se compare au Bale (le chef traditionnel) et se réjouit de le surpasser en visibilité médiatique : elle incarne le pouvoir nouveau de la séduction par l'image.
    \item \textbf{Vision de Lakunlé :} \textbf{ironie amère.} « Le Bale a trois lignes en page six. Toi, tu as le magazine entier » : le contraste chiffré montre que l'image médiatique (modernité) écrase l'autorité traditionnelle (pouvoir réel). « L'image prime sur le pouvoir réel » fonctionne comme une phrase-thèse satirique : Lakunlé, pourtant partisan de la modernité, en dénonce ici la superficialité.
    \item \textbf{Thème :} le \textbf{conflit tradition / modernité} (le Bale contre le magazine et la photographie), le \textbf{pouvoir de l'image et des médias}, la \textbf{beauté féminine} comme levier de pouvoir.
\end{enumerate}
\textbf{Conclusion :} Cet extrait utilise la satire pour révéler la fragilité du pouvoir traditionnel face à la séduction moderne de l'image, tout en moquant la candeur de Sidi.
\end{exoresolu}

\begin{methode}[Identifier et analyser les tonalités : comique, dramatique, satirique]
\textbf{Objectif :} Nommer le ton d'un texte et le justifier par des indices textuels précis.
\begin{enumerate}
    \item \textbf{Le comique :} vise à faire rire. \textbf{Procédés :} quiproquo, répétition, jeu de mots, situation burlesque (incongruité), caricature (exagération des traits), langage familier ou trivial, comique de geste (didascalies), comique de mots.
    \item \textbf{Le dramatique :} vise l'émotion, la tension, la pitié ou la terreur. \textbf{Procédés :} lexique du destin, de la mort, de la fatalité ; phrases nominales et exclamatives ; didascalies lourdes (silences, regards) ; pathos (souffrance morale ou physique) ; sentiment d'inéluctabilité du dénouement.
    \item \textbf{Le satirique :} vise à corriger par le rire (critique sociale ou morale). \textbf{Procédés :} \textbf{ironie} (dire le contraire de ce qu'on pense), \textbf{antiphrase}, \textbf{hyperbole}, \textbf{caricature} des défauts (avarice, vanité, hypocrisie), \textbf{comparaison dévalorisante}, \textbf{parodie} (imiter un style pour le tourner en ridicule). Cible : un individu, un groupe, une institution, la société.
    \item \textbf{Distinction clé :} le comique fait rire pour rire ; le satirique fait rire pour \textbf{dénoncer et corriger} ; le dramatique suscite tension et émotion, sans rire.
\end{enumerate}
\cle{Ironie} \cle{Antiphrase} \cle{Caricature} \cle{Pathos} \cle{Visée corrective}
\end{methode}

\begin{exoresolu}[Classification tonale et justification]
\textbf{Énoncé :} Pour chacun des courts extraits ci-dessous, identifiez la tonalité dominante (comique, dramatique, satirique) et justifiez par deux indices textuels précis.
\begin{enumerate}
    \item \emph{« Maître Corbeau, sur un arbre perché, / Tenait en son bec un fromage. / Maître Renard, par l'odeur alléché, / Lui tint à peu près ce langage : "Hé ! bonjour, Monsieur du Corbeau. / Que vous êtes joli ! que vous me semblez beau !" »} (La Fontaine, \emph{Le Corbeau et le Renard})
    \item \emph{« Il vit le fer s'enfoncer dans sa chair. Pas un cri. Les yeux fixes, il regarda le ciel se refermer sur lui. Le tambour cessa. Le silence était plus lourd que la pierre. »}
    \item \emph{« Notre honorable député, grand pourfendeur de la corruption, vient d'être épinglé pour détournement de fonds publics. Il a déclaré, la main sur le cœur : "Je n'ai fait que prélever ma prime de risque pour avoir risqué ma vie à écouter les doléances du peuple." »}
\end{enumerate}
\tcblower
\textbf{Solution détaillée :}
\begin{enumerate}
    \item \textbf{Tonalité : comique} (au service d'une visée morale, propre à la fable).
    \begin{itemize}
        \item Indice 1 : \textbf{anthropomorphisme} : des animaux titrés « Maître », « Monsieur », qui parlent un langage humain cérémonieux.
        \item Indice 2 : \textbf{flatterie outrancière} : « Que vous êtes joli ! que vous me semblez beau ! » : le Renard flatte le Corbeau pour le duper ; le lecteur, qui devine la ruse, sourit de la naïveté de la victime.
    \end{itemize}
    \item \textbf{Tonalité : dramatique.}
    \begin{itemize}
        \item Indice 1 : \textbf{lexique de la mort et de la souffrance physique} : « le fer s'enfoncer dans sa chair », « le ciel se refermer sur lui » (métaphore de la fin, de la tombe).
        \item Indice 2 : \textbf{rythme haché et poids du silence} : phrases brèves (« Pas un cri. », « Le tambour cessa. ») et comparaison hyperbolique finale (« Le silence était plus lourd que la pierre ») qui insistent sur la gravité de l'instant.
    \end{itemize}
    \item \textbf{Tonalité : satirique.}
    \begin{itemize}
        \item Indice 1 : \textbf{ironie et antiphrase} : « grand pourfendeur de la corruption » désigne un homme « épinglé pour détournement de fonds » ; le contraste violent entre l'image publique et la réalité dénonce l'hypocrisie.
        \item Indice 2 : \textbf{argument fallacieux et cynisme caricatural} : « prélever ma prime de risque pour avoir risqué ma vie à écouter les doléances » détourne le sens des mots (« prime de risque », « risquer sa vie ») pour justifier le vol : caricature du discours politique hypocrite.
    \end{itemize}
\end{enumerate}
\end{exoresolu}

\begin{methode}[Reformuler les idées essentielles (résumé, contraction)]
\textbf{Objectif :} Restituer le sens d'un texte avec ses propres mots, sans déformer la pensée de l'auteur.
\begin{enumerate}
    \item \textbf{Découper le texte en unités de sens :} une idée force par paragraphe ou par séquence logique.
    \item \textbf{Formuler l'idée force de chaque unité :} une phrase simple, à la troisième personne, sous forme impersonnelle (« L'auteur montre que... », « Le texte explique que... »).
    \item \textbf{Éliminer les détails :} pas d'exemples, pas de citations, pas de « par exemple ».
    \item \textbf{Articuler les idées :} utiliser des connecteurs logiques (\emph{donc, or, mais, par conséquent, en outre}) pour rendre visibles les liens (cause, conséquence, opposition, concession).
    \item \textbf{Vérifier la fidélité :} relire l'original et la reformulation. Le sens est-il identique ? Le point de vue de l'auteur est-il respecté (aucune opinion personnelle) ?
\end{enumerate}
\cle{Unité de sens} \cle{Idée force} \cle{Connecteur logique} \cle{Fidélité} \cle{Impersonnalisation}
\end{methode}

\begin{exoresolu}[Reformulation d'idées essentielles (synthèse de trois paragraphes)]
\textbf{Énoncé :} Voici trois paragraphes d'un article. Reformulez l'idée essentielle de chaque paragraphe en une seule phrase, puis reliez les trois phrases en un paragraphe cohérent.
\begin{quote}
\emph{§1 : L'urbanisation galopante en Afrique subsaharienne pose d'immenses défis infrastructuraux. Les villes croissent plus vite que la capacité des États à fournir eau, électricité et transports. \\
§2 : Face à ce déficit, les populations s'organisent en économie informelle. Elles créent des réseaux parallèles d'approvisionnement en eau, de transport (mototaxis) et d'énergie (générateurs). \\
§3 : Cette résilience citoyenne ne dispense pas les pouvoirs publics de leurs obligations. Au contraire, elle doit inspirer des politiques publiques intégrant le secteur informel comme partenaire.}
\end{quote}
\tcblower
\textbf{Solution détaillée :}

\textbf{Étape 1 : Idée force par paragraphe.}
\begin{itemize}
    \item §1 : L'urbanisation rapide en Afrique subsaharienne crée un déficit d'infrastructures (eau, électricité, transport) que les États ne parviennent pas à combler.
    \item §2 : Les populations pallient ce vide par l'économie informelle (réseaux parallèles d'eau, mototaxis, générateurs).
    \item §3 : Cette initiative citoyenne ne décharge pas l'État de ses devoirs ; elle doit au contraire être intégrée aux politiques publiques.
\end{itemize}
\textbf{Étape 2 : Synthèse articulée (reformulation unique).}
\begin{quote}
\emph{« Si l'urbanisation galopante en Afrique subsaharienne engendre un déficit d'infrastructures que les États peinent à résorber, la résilience des populations, qui s'organisent en réseaux informels parallèles pour accéder aux services de base, doit inciter les pouvoirs publics à intégrer ce secteur comme partenaire légitime des politiques urbaines, au lieu de s'en exonérer. »}
\end{quote}
\textbf{Vérification :} les trois idées sont présentes ; les liens logiques sont explicites (\emph{si... / ...doit inciter... au lieu de...}) ; le vocabulaire est reformulé (pas de copier-coller) ; le point de vue de l'auteur est conservé, sans opinion personnelle.
\end{exoresolu}

\begin{methode}[Maîtriser la paronymie et l'antonymie]
\textbf{Objectif :} Distinguer les paronymes (mots proches par la forme, différents par le sens) et utiliser les antonymes (mots de sens opposé) pour préciser une idée ou structurer une argumentation.
\begin{enumerate}
    \item \textbf{Paronymie (pièges à éviter) :} mots de forme proche mais de sens différent. \textbf{Méthode :} mémoriser les paires fréquentes et les associer chacune à une phrase repère.
    \begin{itemize}
        \item \emph{Inférer} (déduire, conclure) / \emph{infirmer} (contredire, remettre en cause).
        \item \emph{Éventualité} (possibilité, cas possible) / \emph{événement} (fait qui se produit).
        \item \emph{Conjecture} (hypothèse, supposition) / \emph{conjoncture} (situation économique, ensemble de circonstances).
        \item \emph{Préjudice} (dommage subi) / \emph{préjugé} (opinion préconçue, sans fondement).
        \item \emph{Attentif} (qui fait attention, soigneux) / \emph{attentatoire} (qui porte atteinte à quelque chose).
    \end{itemize}
    \item \textbf{Antonymie (outil argumentatif) :} mots de sens opposé. \textbf{Types :} complémentaires (\emph{mort / vivant}), graduels (\emph{chaud / froid}, avec des nuances intermédiaires : \emph{tiède}, \emph{frais}), réciproques (\emph{acheter / vendre}).
    \item \textbf{Usage dans l'exposé et l'argumentation :} créer des contrastes (antithèses), définir par la négative (« ce n'est pas X, c'est Y »), structurer un plan dialectique (thèse / antithèse).
\end{enumerate}
\cle{Paronyme} \cle{Antonyme} \cle{Antithèse} \cle{Antonymes graduels} \cle{Antonymes complémentaires}
\end{methode}

\begin{exoresolu}[Paronymie, antonymie et préparation d'un exposé oral]
\textbf{Énoncé :}
\begin{enumerate}
    \item Choisissez le bon paronyme entre parenthèses :
    \begin{enumerate}
        \item Le tribunal a reconnu l'existence d'un \emph{(préjudice / préjugé)} moral.
        \item C'est une simple \emph{(conjecture / conjoncture)} sans preuve formelle.
        \item Son attitude est \emph{(attentive / attentatoire)} à la dignité humaine.
    \end{enumerate}
    \item Remplacez le mot en gras par un terme de même famille de sens mais d'intensité moindre, pour nuancer la phrase : « La situation économique est \textbf{catastrophique}. »
    \item Vous devez présenter un exposé oral de 5 minutes sur « L'apport du secteur informel à l'économie urbaine ». Donnez la structure type (introduction, développement, conclusion) avec deux arguments pour le développement.
\end{enumerate}
\tcblower
\textbf{Solution détaillée :}
\begin{enumerate}
    \item \textbf{Paronymes :}
    \begin{enumerate}
        \item \textbf{Préjudice} (dommage subi). \emph{Préjugé} = opinion préconçue, non fondée.
        \item \textbf{Conjecture} (hypothèse, supposition). \emph{Conjoncture} = situation économique, ensemble des circonstances.
        \item \textbf{Attentatoire} (qui porte atteinte). \emph{Attentive} = qui fait attention.
    \end{enumerate}
    \item \textbf{Nuance (échelle d'intensité) :} « La situation économique est \textbf{préoccupante} » (ou \emph{difficile}, \emph{fragile}, \emph{mauvaise}). \emph{Catastrophique} marque le degré extrême ; choisir un terme plus faible permet de nuancer le propos, comme le permet l'antonymie graduelle.
    \item \textbf{Structure de l'exposé oral (5 minutes) :}
    \begin{itemize}
        \item \textbf{Introduction (1 min) :} accroche (un chiffre ou un fait marquant : le secteur informel occupe la majorité des actifs urbains en Afrique subsaharienne), définition du secteur informel, problématique (frein ou moteur de l'économie urbaine ?), annonce du plan.
        \item \textbf{Développement (3 min) :}
        \begin{itemize}
            \item \textbf{Argument 1 (les apports) :} le secteur informel pourvoit massivement à l'emploi (jeunes, femmes), fait preuve de flexibilité et répond à une demande que le secteur formel ne satisfait pas (transport, commerce de proximité, logement). Exemples : mototaxis, vendeurs ambulants.
            \item \textbf{Argument 2 (les limites et les pistes) :} précarité (absence de protection sociale), faible contribution fiscale, insécurité juridique. D'où la nécessité d'une transition : formalisation progressive, simplification administrative, accès au crédit.
        \end{itemize}
        \item \textbf{Conclusion (1 min) :} synthèse (l'informel n'est pas une marge, c'est le cœur battant de la ville), ouverture (des politiques inclusives : « passer de la tolérance à la reconnaissance »), remerciements et invitation aux questions.
    \end{itemize}
\end{enumerate}
\end{exoresolu}

\begin{attention}[Les 5 erreurs qui coûtent des points au Probatoire et au Baccalauréat technique]
\begin{enumerate}
    \item \textbf{Confondre résumé et commentaire (ou dissertation) :} dans une contraction ou un résumé, \textbf{aucune opinion personnelle}, aucune analyse de style, aucun « je pense que ». On restitue \textbf{uniquement} les idées de l'auteur, dans leur ordre. \emph{Sanction :} hors sujet, note très lourdement pénalisée.
    \item \textbf{Copier-coller des phrases du texte :} la contraction exige une \textbf{réécriture} (synonymes, changement de syntaxe, nominalisation). Recopier des segments entiers fait perdre tous les points de « reformulation ».
    \item \textbf{Oublier les connecteurs logiques (texte « télégramme ») :} une suite de phrases sans liens (\emph{donc, mais, cependant, par conséquent, en outre}) montre que la \textbf{structure argumentative} n'a pas été comprise ; le correcteur ne suit plus le raisonnement.
    \item \textbf{Confondre les paronymes dans la copie :} écrire « il a fait un \emph{préjugé} à la société » au lieu de « \emph{préjudice} », ou « la \emph{conjoncture} est hasardeuse » pour « \emph{conjecture} ». Ces confusions trahissent un lexique approximatif et coûtent des points de langue.
    \item \textbf{Identifier un registre tonal à l'envers :} dire qu'un texte est « comique » parce qu'il contient de l'ironie (c'est alors satirique), ou « dramatique » parce qu'un personnage crie (cela peut être comique ou pathétique). \textbf{Justifiez toujours par des procédés précis} (lexique, syntaxe, figures de style, didascalies), jamais par une simple impression.
\end{enumerate}
\end{attention}

\newpage

\coursec{Exercices}

\courssub{Je vérifie mes connaissances}

\begin{exercice}[QCM : le texte argumentatif]
Pour chaque question, entoure la (ou les) bonne(s) réponse(s).
\begin{enumerate}
\item Un texte argumentatif a pour but principal :
\begin{enumerate}
\item[a)] de raconter une histoire ;
\item[b)] de convaincre ou de persuader le lecteur ;
\item[c)] de décrire un paysage.
\end{enumerate}
\item La \cle{thèse} d'un texte argumentatif est :
\begin{enumerate}
\item[a)] l'ensemble des exemples du texte ;
\item[b)] l'opinion que l'auteur défend ;
\item[c)] le titre du texte.
\end{enumerate}
\item Un \cle{argument} est :
\begin{enumerate}
\item[a)] une raison avancée pour soutenir la thèse ;
\item[b)] une insulte adressée à l'adversaire ;
\item[c)] une figure de style.
\end{enumerate}
\item Dans un résumé-contraction, il est interdit :
\begin{enumerate}
\item[a)] de reformuler avec ses propres mots ;
\item[b)] d'ajouter son opinion personnelle ;
\item[c)] de supprimer les exemples.
\end{enumerate}
\end{enumerate}
\end{exercice}

\begin{exercice}[Vrai ou faux : justifie ta réponse]
Réponds par vrai ou faux, puis justifie en une phrase.
\begin{enumerate}
\item La contraction de texte consiste à recopier les phrases les plus importantes du texte.
\item La généralisation est une technique de réduction qui consiste à remplacer plusieurs termes précis par un terme plus général.
\item La tonalité comique d'un passage vise toujours à émouvoir le lecteur.
\item Deux paronymes sont des mots de sens opposé.
\item Dans un exposé oral, le plan doit être annoncé dès l'introduction.
\end{enumerate}
\end{exercice}

\begin{exercice}[Définitions]
Définis en deux ou trois phrases chacune des notions suivantes, en donnant un exemple pour chacune :
\begin{enumerate}
\item la \cle{paronymie} ;
\item l'\cle{antonymie} ;
\item la \cle{tonalité satirique} ;
\item la \cle{fusion} (technique de réduction).
\end{enumerate}
\end{exercice}

\begin{exercice}[Repérage rapide]
Voici quatre extraits brefs. Indique pour chacun la tonalité dominante (comique, dramatique ou satirique) et cite un procédé qui la manifeste.
\begin{enumerate}
\item « Le ministre arriva avec trois heures de retard, salua gravement la foule et déclara que la ponctualité était la politesse des rois. »
\item « Elle serra contre sa poitrine la lettre froissée : c'était la dernière que son fils lui avait écrite avant de disparaître. »
\item « Il trébucha sur son propre parapluie, se rattrapa au chapeau du voisin et finit assis dans le plat de foufou. »
\item « Nos vaillants députés, après d'âpres débats sur leurs indemnités, ajournèrent la discussion sur l'école du village. »
\end{enumerate}
\end{exercice}

\courssub{Je m'entraîne}

\begin{exercice}[Paronymes : le bon choix]
Complète les dix phrases suivantes en choisissant le bon paronyme dans chaque couple, puis reformule chaque phrase en utilisant un \cle{antonyme} du mot souligné (ou du mot choisi) sans changer le sens global.
\begin{enumerate}
\item Il faut \dots\ (effectuer / affecter) un contrôle avant de valider le bulletin.
\item Le directeur a été \dots\ (affecté / effectué) au lycée technique de Garoua.
\item Une tempête est \dots\ (éminente / imminente) : le ciel est tout noir.
\item C'est un professeur \dots\ (éminent / imminent), respecté de tous.
\item Le prix du sac \dots\ (inclut / incline) les frais de transport.
\item Le malade \dots\ (inclut / incline) la tête pour saluer le médecin.
\item Cette histoire est \dots\ (vraisemblable / véridique) : elle ressemble à la réalité.
\item Son témoignage est \dots\ (vraisemblable / véridique) : il dit exactement ce qui s'est passé.
\item Il a fait une \dots\ (allusion / illusion) discrète au départ de son ami.
\item Le mirage n'est qu'une \dots\ (allusion / illusion) d'optique.
\end{enumerate}
\end{exercice}

\begin{exercice}[Contraction guidée d'un paragraphe]
Voici un paragraphe de 80 mots. Réduis-le à environ 40 mots en appliquant successivement les quatre techniques de réduction : suppression, sélection, généralisation, fusion. Justifie chaque choix en une phrase.

\emph{« Les réseaux sociaux occupent une place énorme, considérable, gigantesque dans la vie des jeunes Camerounais. Le matin au réveil, ils consultent WhatsApp ; à la pause de dix heures, ils ouvrent Facebook ; le soir après les devoirs, ils regardent des vidéos sur TikTok. Ces plateformes, qui coûtent cher en forfaits Internet, permettent pourtant de s'informer, d'apprendre et de communiquer avec des amis éloignés, par exemple ceux qui étudient à l'étranger. »}
\end{exercice}

\begin{exercice}[Identifier la structure argumentative]
Lis le texte suivant, puis réponds aux questions.

\emph{« Le téléphone portable à l'école est un problème. Certes, il permet de joindre rapidement les parents en cas d'urgence. Mais en classe, il distrait les élèves : au lieu d'écouter l'explication du professeur de mathématiques, beaucoup tapotent sous la table. De plus, les notes chutent lorsque les nuits sont passées à discuter en ligne. Enfin, le portable favorise la tricherie pendant les évaluations. C'est pourquoi il faut l'interdire strictement dans les salles de classe. »}
\begin{enumerate}
\item Dégage la thèse de l'auteur en une phrase.
\item Releve les trois arguments qui soutiennent cette thèse.
\item Quelle concession l'auteur accorde-t-il à la thèse adverse ? Par quel connecteur est-elle introduite ?
\item Releve deux connecteurs logiques et précise leur fonction (addition, conséquence, opposition...).
\end{enumerate}
\end{exercice}

\begin{exercice}[Tonalités dans \emph{Le Lion et la perle}]
Dans la pièce de Wole Soyinka, relis la scène où Baroka, le Lion, feint d'être devenu impuissant pour tromper Sidi.
\begin{enumerate}
\item Quelle est la tonalité dominante de cette scène ?
\item Releve trois procédés (gestes, paroles, ironie, quiproquo) qui manifestent cette tonalité.
\item Explique en quelques lignes ce que cette scène révèle du personnage de Baroka.
\end{enumerate}
\end{exercice}

\courssub{Je me prépare au Bac}

\begin{exercice}[Type Baccalauréat : résumé d'un texte argumentatif]
Lis attentivement le texte suivant, puis traite les questions.

\textbf{Texte : Les réseaux sociaux et la jeunesse camerounaise}

\emph{De Yaoundé à Maroua, les écrans ont envahi le quotidien des jeunes. Dans les quartiers, dans les salles de classe et jusque dans les champs, le téléphone portable est devenu le compagnon inséparable de toute une génération. Faut-il s'en réjouir ou s'en inquiéter ? À nos yeux, les réseaux sociaux, utilisés sans discernement, menacent la réussite scolaire et l'équilibre des adolescents camerounais.}

\emph{D'abord, ils dévorent le temps d'étude. Combien d'élèves sacrifient leurs nuits à des conversations sans fin, puis s'endorment sur leur table au moment des leçons ? Une enquête menée dans plusieurs lycées de la capitale montre que les élèves les plus connectés sont aussi ceux dont les résultats baissent le plus vite. Le temps, cette richesse irremplaçable, s'évapore en « likes » et en commentaires.}

\emph{Ensuite, les réseaux sociaux exposent les jeunes à des modèles dangereux. La course aux apparences pousse certains adolescents à copier des comportements étrangers à nos valeurs : vulgarité, mépris du travail, culte de l'argent facile. Des faits divers récents l'ont montré : des élèves se sont laissé piéger par de fausses promesses de gains en ligne.}

\emph{Certes, il serait injuste de tout condamner. Ces plateformes permettent de s'informer, de suivre des cours gratuits, de garder le contact avec la famille restée au village. Un élève de Première peut y trouver des explications de mathématiques ou des conseils d'orientation. L'outil, en soi, n'est pas mauvais.}

\emph{Mais c'est précisément là que le problème se pose : un outil exige une main qui le maîtrise. Or nos jeunes sont livrés seuls à la tentation, sans guide ni règle. C'est pourquoi parents, enseignants et pouvoirs publics doivent unir leurs efforts pour éduquer à un usage raisonnable des réseaux sociaux. La jeunesse du Cameroun mérite mieux que d'être esclave de son écran : elle mérite d'en devenir le maître.}

\begin{enumerate}
\item Dégage la thèse de l'auteur en une phrase.
\item Releve les deux arguments principaux et la concession faite à la thèse adverse.
\item Résume ce texte en 90 mots (plus ou moins 10 \%), en respectant les règles de la contraction : reformulation personnelle, fidélité aux idées de l'auteur, absence de commentaire et de jugement personnel.
\item Indique le nombre exact de mots de ton résumé.
\end{enumerate}
\end{exercice}

\begin{exercice}[Type Probatoire : lecture méthodique de \emph{Le Lion et la perle}]
Lis l'extrait suivant (adapté de la pièce de Wole Soyinka), puis réponds aux questions.

\emph{Sidi, éblouie par ses photographies publiées dans le magazine, dédaigne Lakunle, l'instituteur qui veut l'épouser sans payer la dot. Baroka, le vieux chef du village, jaloux de la beauté de la jeune fille, fait répandre la rumeur de sa propre impuissance pour l'attirer dans son palais. Sidi, moqueuse, vient narguer le vieillard qu'elle croit vaincu. Mais le Lion n'a rien perdu de sa ruse : la perle du village tombe dans le piège, et c'est finalement le vieux chef traditionnel qu'elle épousera, abandonnant l'instituteur et ses discours sur le progrès.}

\begin{enumerate}
\item Quelle est la tonalité dominante de cet extrait ? Justifie ta réponse en relevant trois procédés (ironie, situation de quiproquo, mots comiques, retournement de situation...).
\item Dégage les idées essentielles de l'extrait en trois ou quatre phrases, sans recopier le texte.
\item Quel conflit de valeurs cet épisode met-il en scène ? Appuie ta réponse sur deux éléments précis de la pièce.
\item Reformule en une phrase personnelle le message que l'auteur semble adresser au spectateur à travers la défaite de Lakunle.
\end{enumerate}
\end{exercice}

\begin{exercice}[Exposé oral évalué : tradition et modernité]
Tu dois présenter devant la classe un exposé de cinq minutes sur le sujet suivant : \emph{« Tradition et modernité dans Le Lion et la perle »}. Prépare ton intervention par écrit, puis présente-la à l'oral.
\begin{enumerate}
\item Rédige une introduction qui accroche l'attention des auditeurs, présente le sujet et annonce clairement le plan en deux parties.
\item Développe deux arguments étayés chacun par un exemple précis tiré de la pièce (personnage, scène, réplique) : un argument sur la force de la tradition, un argument sur l'ambiguïté de la modernité.
\item Rédige une conclusion personnelle et justifiée : selon toi, la pièce invite-t-elle à choisir entre tradition et modernité, ou à les concilier ?
\item À l'oral, tu seras évalué sur : le respect du temps (5 minutes), la clarté du plan, la correction de la langue et la qualité du contact avec l'auditoire (regard, voix, gestuelle).
\end{enumerate}
\end{exercice}

\newpage

\coursec{Corrigés des exercices}

\coursec{Corrigés détaillés — Leçon 15 : Formuler les idées essentielles d'un texte à résumer}

\courssub{Bloc 1 : Le texte argumentatif et les techniques de contraction}

\begin{corrige}[Exercice 1 — QCM : le texte argumentatif]
\begin{enumerate}
\item \textbf{Bonne réponse : b)} de convaincre ou de persuader le lecteur. \cle{But du texte argumentatif} : faire adhérer le destinataire à une opinion ; raconter relève du \cle{texte narratif}, décrire du \cle{texte descriptif}.
\item \textbf{Bonne réponse : b)} l'opinion que l'auteur défend. La \cle{thèse} est le point de vue que l'auteur cherche à faire accepter ; les exemples servent à l'illustrer, le titre peut l'annoncer mais ne s'y confond pas.
\item \textbf{Bonne réponse : a)} une raison avancée pour soutenir la thèse. Un \cle{argument} est un raisonnement logique ; l'insulte est un procédé disqualifié (\cle{argument ad hominem}), et la figure de style n'est qu'un ornement éventuel.
\item \textbf{Bonne réponse : b)} d'ajouter son opinion personnelle. Le \cle{résumé-contraction} exige la \cle{fidélité} totale à l'auteur : \cle{reformuler} est obligatoire, \cle{supprimer les exemples} est une technique de réduction autorisée, mais tout commentaire personnel est interdit.
\end{enumerate}
\end{corrige}

\begin{corrige}[Exercice 2 — Vrai ou faux : contraction et tonalités]
\begin{enumerate}
\item \textbf{Faux.} La contraction interdit de recopier le texte : il faut \cle{reformuler} les idées essentielles avec ses propres mots, sinon on commet un \cle{plagiat}.
\item \textbf{Vrai.} La \cle{généralisation} remplace un ensemble de termes précis par un terme englobant : « WhatsApp, Facebook, TikTok » devient « les réseaux sociaux ».
\item \textbf{Faux.} La \cle{tonalité comique} vise à faire rire ou sourire ; c'est la \cle{tonalité dramatique} qui cherche à émouvoir (\cle{pitié}, \cle{terreur}).
\item \textbf{Faux.} Les \cle{paronymes} sont des mots de formes voisines mais de sens différents (\emph{effectuer}/\emph{affecter}) ; les mots de sens opposé sont des \cle{antonymes}.
\item \textbf{Vrai.} Annoncer le \cle{plan} dès l'introduction guide l'auditoire et structure l'écoute : c'est une règle fondamentale de l'\cle{exposé oral}.
\end{enumerate}
\end{corrige}

\begin{corrige}[Exercice 3 — Définitions clés de l'unité]
\begin{enumerate}
\item \textbf{La \cle{paronymie}} est la relation qui unit deux mots de formes (graphie ou prononciation) très proches mais de sens différents. Elle exige une grande vigilance pour éviter les \cle{confusions lexicales}. \emph{Exemple :} \emph{éminent} (remarquable) et \emph{imminent} (qui va arriver très vite).
\item \textbf{L'\cle{antonymie}} est la relation qui unit deux mots de sens opposés. Elle permet de \cle{nuancer} ou de \cle{construire un contraste} dans un texte. \emph{Exemple :} \emph{généreux} et \emph{avare}.
\item \textbf{La \cle{tonalité satirique}} est le registre d'un texte qui dénonce les défauts, les vices ou les ridicules d'une personne ou d'une société en se moquant d'eux. Elle associe le \cle{rire} et la \cle{critique}. \emph{Exemple :} un passage qui ridiculise un ministre vantant la ponctualité tout en arrivant en retard.
\item \textbf{La \cle{fusion}} est une technique de réduction qui consiste à regrouper en une seule phrase les idées de plusieurs phrases du texte de départ. Elle s'appuie souvent sur des \cle{connecteurs logiques} (cause, conséquence, opposition). \emph{Exemple :} « Il pleut. Je reste à la maison. » devient « \cle{Comme} il pleut, je reste à la maison. »
\end{enumerate}
\end{corrige}

\begin{corrige}[Exercice 4 — Repérage rapide des tonalités]
\begin{enumerate}
\item \textbf{Tonalité \cle{satirique}.} Le \cle{décalage ironique} entre les trois heures de retard du ministre et son éloge de la ponctualité dénonce l'hypocrisie des puissants (\cle{ironie de situation}).
\item \textbf{Tonalité \cle{dramatique}.} Le \cle{pathétique} du geste (« serra contre sa poitrine la lettre froissée ») et le mot « disparaître » suscitent l'\cle{émotion} et la \cle{pitié}.
\item \textbf{Tonalité \cle{comique}.} L'\cle{accumulation} de gaffes (trébucher, se rattraper au chapeau, tomber dans le foufou) relève de la \cle{farce} et provoque le \cle{rire}.
\item \textbf{Tonalité \cle{satirique}.} L'\cle{antithèse} entre « vaillants députés » et leur comportement égoïste (les indemnités avant l'école du village) \cle{raille} leurs priorités.
\end{enumerate}
\end{corrige}

\courssub{Bloc 2 : Paronymie, antonymie et emploi précis}

\begin{corrige}[Exercice 5 — Paronymes : le bon choix et antonymes]
\begin{enumerate}
\item Il faut \textbf{effectuer} un contrôle. \cle{Antonyme} : « Il ne faut pas \emph{négliger} le contrôle avant de valider le bulletin. »
\item Le directeur a été \textbf{affecté} au lycée technique de Garoua. \cle{Antonyme} : « Le directeur a été \emph{non affecté} / \emph{sans poste} » (sens opposé : il n'a pas reçu d'affectation).
\item Une tempête est \textbf{imminente}. \cle{Antonyme} : « Une tempête est \emph{improbable} : le ciel est dégagé. »
\item C'est un professeur \textbf{éminent}. \cle{Antonyme} : « C'est un professeur \emph{obscur} / \emph{médiocre}, que personne ne remarque. »
\item Le prix du sac \textbf{inclut} les frais de transport. \cle{Antonyme} : « Le prix du sac \emph{exclut} les frais de transport. »
\item Le malade \textbf{incline} la tête. \cle{Antonyme} : « Le malade \emph{redresse} la tête pour saluer le médecin. »
\item Cette histoire est \textbf{vraisemblable}. \cle{Antonyme} : « Cette histoire est \emph{invraisemblable} : elle ne ressemble en rien à la réalité. »
\item Son témoignage est \textbf{véridique}. \cle{Antonyme} : « Son témoignage est \emph{mensonger} : il déforme ce qui s'est passé. »
\item Il a fait une \textbf{allusion} discrète. \cle{Antonyme} : « Il a fait une \emph{mention explicite} du départ de son ami. »
\item Le mirage n'est qu'une \textbf{illusion} d'optique. \cle{Antonyme} : « Ce que nous voyons est une \emph{réalité}, non un mirage. »
\end{enumerate}
\begin{attention}[Pièges fréquents]
\emph{Effectuer} = réaliser, accomplir (une action) ; \emph{affecter} = assigner à un poste (administratif) ou toucher émotionnellement. Ne jamais les intervertir à l'écrit. \emph{Imminent} (temps) vs \emph{éminent} (qualité). \emph{Inclure} (contenir) vs \emph{incliner} (pencher). \emph{Vraisemblable} (apparence de vérité) vs \emph{véridique} (conforme à la vérité). \emph{Allusion} (indirect) vs \emph{illusion} (erreur des sens).
\end{attention}
\end{corrige}

\courssub{Bloc 3 : Contraction guidée et structure argumentative}

\begin{corrige}[Exercice 6 — Contraction guidée d'un paragraphe (texte sur les réseaux sociaux)]
\textbf{Résumé obtenu (40 mots) :} « Les réseaux sociaux occupent une place considérable dans la vie des jeunes Camerounais, qui les consultent à tout moment de la journée. Malgré leur coût, ces plateformes permettent de s'informer, d'apprendre et de communiquer avec des proches éloignés. »

\textbf{Justification des techniques employées :}
\begin{enumerate}
\item \textbf{\cle{Suppression} :} les trois synonymes « énorme, considérable, gigantesque » sont réduits à un seul adjectif, car la répétition n'ajoute aucune idée nouvelle (\cle{redondance}).
\item \textbf{\cle{Généralisation} :} « WhatsApp, Facebook, TikTok » et les trois moments de la journée deviennent « les réseaux sociaux » et « à tout moment de la journée » : un \cle{terme hyperonyme} remplace les énumérations (\cle{hyponymes}).
\item \textbf{\cle{Sélection} :} l'exemple « ceux qui étudient à l'étranger » est écarté : c'est une \cle{illustration secondaire}, non une idée essentielle.
\item \textbf{\cle{Fusion} :} « s'informer, apprendre, communiquer » sont regroupés dans une seule proposition introduite par « Malgré leur coût », qui conserve l'\cle{opposition} coût/avantages présente dans le texte (\cle{concession}).
\end{enumerate}
\begin{attention}[Point de vigilance]
L'idée de nuance (« qui coûtent cher... pourtant ») doit absolument être conservée : supprimer la concession déformerait la \cle{pensée de l'auteur}.
\end{attention}
\end{corrige}

\begin{corrige}[Exercice 7 — Identifier la structure argumentative (texte : portable en classe)]
\begin{enumerate}
\item \textbf{\cle{Thèse} :} le téléphone portable doit être strictement interdit dans les salles de classe car il nuit aux élèves.
\item \textbf{Les trois \cle{arguments} :}
\begin{enumerate}
\item[a)] il \cle{distrait} les élèves pendant les cours (ils tapotent au lieu d'écouter) ;
\item[b)] il fait \cle{chuter les résultats} scolaires (nuits passées en ligne) ;
\item[c)] il \cle{favorise la tricherie} pendant les évaluations.
\end{enumerate}
\item \textbf{\cle{Concession} :} le portable permet de joindre rapidement les parents en cas d'urgence. Elle est introduite par le connecteur d'admission « \textbf{Certes} », suivi de l'opposition « \textbf{Mais} » qui marque le \cle{retour à la thèse} de l'auteur (\cle{structure concessive}).
\item \textbf{\cle{Connecteurs} logiques :} « \textbf{De plus} » (\cle{addition} : il ajoute un deuxième argument) ; « \textbf{Enfin} » (\cle{énumération} : dernier argument) ; « \textbf{C'est pourquoi} » (\cle{conséquence}/\cle{conclusion} : il introduit la thèse finale).
\end{enumerate}
\begin{aretenir}[Structure type repérée]
\emph{Thèse annoncée — Concession (Certes... Mais) — Arguments — Thèse reformulée}. Savoir repérer le couple « Certes... Mais » est indispensable pour ne pas attribuer à l'auteur l'opinion adverse.
\end{aretenir}
\end{corrige}

\courssub{Bloc 4 : Lectures méthodiques de \emph{Le Lion et la perle} (Wole Soyinka)}

\begin{corrige}[Exercice 8 — Tonalités dans \emph{Le Lion et la perle} (Lecture n° 4 : scène de la fausse impuissance)]
\begin{enumerate}
\item \textbf{Tonalité dominante : \cle{comique}}, teintée d'\cle{ironie satirique} : le spectateur, qui connaît la ruse de Baroka, rit de la naïveté de Sidi.
\item \textbf{Trois \cle{procédés comiques/satiriques} :}
\begin{enumerate}
\item[a)] le \cle{quiproquo} : Sidi croit le vieux chef impuissant et vaincu, alors qu'il est parfaitement lucide et manipulateur ;
\item[b)] le \cle{jeu théâtral de Baroka} : il feint la faiblesse, gémit, se plaint de sa vieillesse avec une \cle{mauvaise foi} visible (\cle{comique de gestes} et de \cle{mots}) ;
\item[c)] l'\cle{ironie dramatique} : le public sait que la rumeur est un piège, ce qui rend les moqueries de Sidi risibles et \cle{retournables contre elle}.
\end{enumerate}
\item \textbf{Ce que la scène révèle :} Baroka apparaît comme un personnage \cle{rusé}, \cle{intelligent} et \cle{calculateur}, fidèle à son surnom de « \cle{Lion} ». Loin d'être un vieillard fini, il \cle{manipule l'information} et les \cle{apparences} pour atteindre ses fins. La scène montre que la \cle{tradition}, qu'il incarne, possède encore une \cle{redoutable vitalité} face à la modernité naïve incarnée par Sidi et Lakunle.
\end{enumerate}
\end{corrige}

\begin{corrige}[Exercice 10 — Type Probatoire : Lecture méthodique n° 5 (Dénouement)]
\begin{enumerate}
\item \textbf{Tonalité dominante : \cle{comique} et \cle{satirique}.} Trois procédés :
\begin{enumerate}
\item[a)] le \cle{quiproquo} : Sidi croit Baroka impuissant alors qu'il a tout manigancé ;
\item[b)] l'\cle{ironie} : c'est en venant narguer le Lion que la « perle » tombe dans son piège — la \cle{moqueuse devient la dupée} ;
\item[c)] le \cle{retournement de situation} (\cle{péripétie}) : le vieillard qu'on croyait vaincu triomphe finalement du jeune instituteur.
\end{enumerate}
\item \textbf{Idées essentielles (\cle{reformulées}) :} Fière de sa célébrité soudaine, Sidi rejette le mariage sans dot que lui propose Lakunle. Le chef Baroka, désireux de la conquérir, invente une rumeur sur sa propre faiblesse pour l'attirer chez lui. Trompée par cette ruse, Sidi finit par épouser le vieux chef, abandonnant l'instituteur moderniste.
\item \textbf{Conflit de valeurs :} la pièce oppose la \cle{tradition} (Baroka, la \cle{dot}, la \cle{polygamie}, la ruse du Lion) à la \cle{modernité} (Lakunle, le \cle{refus de la dot}, les discours sur le \cle{progrès}, le magazine). Éléments précis : le refus de Lakunle de payer la dot, qu'il juge arriérée, et le triomphe final de Baroka qui épouse Sidi selon les \cle{coutumes}.
\item \textbf{Message \cle{reformulé} :} à travers l'échec de Lakunle, Soyinka suggère qu'une modernité \cle{superficielle}, faite de slogans et de mépris des coutumes, ne peut l'emporter sur une tradition \cle{vivante} et \cle{intelligente}.
\end{enumerate}
\begin{attention}[Consignes d'examen]
\begin{itemize}
\item Distinguer le \cle{comique de situation} (quiproquo) de l'\cle{ironie} (écart entre ce que croit Sidi et ce que sait le spectateur).
\item Question 2 : aucune phrase ne doit être reprise mot à mot de l'extrait (\cle{reformulation} obligatoire).
\item Question 4 : exige une formulation \cle{personnelle}, pas une paraphrase.
\end{itemize}
\end{attention}
\end{corrige}

\courssub{Bloc 5 : Épreuves types Baccalauréat / Probatoire}

\begin{corrige}[Exercice 9 — Type Baccalauréat : Résumé d'un texte argumentatif (Réseaux sociaux et jeunesse)]
\textbf{1. \cle{Thèse} dégagée :} les réseaux sociaux, utilisés sans discernement, menacent la réussite scolaire et l'équilibre des jeunes Camerounais ; il faut donc \cle{éduquer} la jeunesse à un usage maîtrisé de ces outils.

\textbf{2. \cle{Arguments} et \cle{concession} :}
\begin{itemize}
\item \textbf{Argument 1 :} ils \cle{dévorent le temps d'étude} (nuits passées en ligne, sommeil en classe, baisse des résultats confirmée par une enquête).
\item \textbf{Argument 2 :} ils \cle{exposent les jeunes à des modèles dangereux} (culte de l'argent facile, vulgarité, arnaques en ligne).
\item \textbf{Concession :} ces plateformes rendent aussi des \cle{services} (information, cours gratuits, lien familial) ; l'outil en soi n'est pas mauvais.
\end{itemize}

\textbf{3. Résumé proposé (91 mots) :}

« Les écrans ont envahi la vie des jeunes Camerounais. Or, selon l'auteur, les réseaux sociaux mal utilisés compromettent leur réussite et leur équilibre. D'une part, ils absorbent le temps d'étude : les élèves très connectés voient leurs résultats baisser. D'autre part, ils diffusent des modèles néfastes — argent facile, vulgarité — et piègent certains adolescents. L'auteur reconnaît toutefois que ces outils permettent de s'informer, d'apprendre et de rester en contact avec les proches. Mais un outil exige d'être maîtrisé : parents, enseignants et pouvoirs publics doivent donc éduquer les jeunes à un usage raisonnable, afin qu'ils dominent l'écran au lieu d'en être esclaves. »

\textbf{4. \cle{Nombre de mots} : 91} (dans la fourchette autorisée de 81 à 99 mots).

\textbf{\cle{Barème indicatif} (sur 10) :}
\begin{itemize}
\item Thèse correctement dégagée : 2 pts
\item Arguments et concession relevés : 2 pts
\item Respect du nombre de mots ($\pm$ 10\%) : 1 pt
\item Fidélité aux idées de l'auteur sans avis personnel : 2 pts
\item Reformulation personnelle (aucune phrase recopiée) : 2 pts
\item Correction de la langue (orthographe, syntaxe, ponctuation) : 1 pt
\end{itemize}

\begin{attention}[Points de vigilance pour l'examen]
\begin{itemize}
\item Ne jamais écrire « je pense que... » (\cle{impersonnalisation}).
\item Ne pas recopier les exemples (l'enquête, les faits divers) mais les \cle{généraliser}.
\item Conserver la \cle{concession} et la \cle{conclusion} de l'auteur, sinon le résumé déforme sa pensée.
\item Compter les mots avec soin et l'indiquer en fin de copie.
\end{itemize}
\end{attention}
\end{corrige}

\begin{corrige}[Exercice 11 — Exposé oral évalué : Tradition et modernité dans \emph{Le Lion et la perle}]
\textbf{1. \cle{Introduction} (modèle) :} « Mesdames et Messieurs, chers camarades : peut-on moderniser un village sans tuer son âme ? C'est la question que pose Wole Soyinka dans \emph{Le Lion et la perle}, où s'affrontent Baroka, le chef traditionnel, et Lakunle, l'instituteur moderniste, pour conquérir Sidi, la perle du village. \cle{Plan annoncé} : nous verrons d'abord la force de la tradition incarnée par Baroka, puis l'ambiguïté de la modernité représentée par Lakunle. »

\textbf{2. \cle{Développement} (modèle) :}
\begin{itemize}
\item \textbf{Argument 1 — la force de la tradition :} Baroka, surnommé le Lion, incarne une tradition \cle{rusée} et \cle{vivante}. \emph{Exemple précis :} dans la scène de la fausse impuissance (Lecture 4), il manipule Sidi en feignant la faiblesse, prouvant que la tradition n'est pas une survivance figée mais une \cle{intelligence à l'œuvre}. Sa victoire finale montre sa capacité à \cle{s'adapter} (il utilise même le magazine, objet moderne, à son profit).
\item \textbf{Argument 2 — l'ambiguïté de la modernité :} Lakunle prêche le progrès mais apparaît \cle{ridicule} et \cle{incohérent}. \emph{Exemple précis :} il refuse de payer la dot au nom de la civilisation, mais ses discours pompeux et ses manières maladroites le rendent comique ; sa modernité reste \cle{verbale}, sans prise sur la réalité du village.
\end{itemize}

\textbf{3. \cle{Conclusion} (modèle) :} « La pièce n'invite pas à choisir brutalement entre tradition et modernité : Baroka lui-même adopte certains objets modernes, tandis que la modernité de Lakunle échoue parce qu'elle méprise les coutumes. Soyinka semble donc plaider pour une \cle{modernité enracinée}, capable de concilier le progrès et l'héritage culturel africain. »

\textbf{4. \cle{Conseils pour l'oral} :} chronométrer l'entraînement (environ 1 min d'introduction, 3 min de développement, 1 min de conclusion) ; regarder l'auditoire, pas ses notes ; parler lentement et articuler ; annoncer les transitions (« d'abord », « ensuite », « pour conclure »).

\textbf{\cle{Barème indicatif} (sur 10) :}
\begin{itemize}
\item Introduction avec accroche et plan annoncé : 2 pts
\item Deux arguments étayés par des \cle{exemples précis} (scène, réplique, personnage) : 4 pts
\item Conclusion \cle{personnelle et justifiée} par le contenu de la pièce : 2 pts
\item Prestation orale (temps, voix, regard, langue) : 2 pts
\end{itemize}

\begin{attention}[Erreurs pénalisant l'oral]
\begin{itemize}
\item Exemple non précis (« dans la pièce, on voit que... ») : ne vaut pas un exemple \cle{cité}.
\item Conclusion non justifiée (simple opinion) : hors sujet.
\item Dépassement du temps imparti : pénalité.
\item Lecture du texte : note éliminatoire pour la prestation.
\end{itemize}
\end{attention}
\end{corrige}

\newpage

\coursec{Fiche bilan}

\begin{popfigure}
\begin{center}\tcbox[colback=popblue,colframe=popblue,coltext=white,arc=9pt,boxsep=4pt,left=6pt,right=6pt]{\bfseries\small Idées essentielles d'un texte à résumer}\end{center}
\vspace{-2pt}
\begin{tcbraster}[raster columns=3,raster equal height=rows,raster column skip=6pt,raster row skip=6pt,enhanced,blankest]
\begin{tcolorbox}[enhanced,colback=white,colframe=poporange,boxrule=0.9pt,arc=3pt,title={Texte argumentatif},fonttitle=\bfseries\scriptsize,coltitle=white,colbacktitle=poporange,left=4pt,right=4pt,top=2pt,bottom=2pt,boxsep=1pt,toptitle=1.5pt,bottomtitle=1.5pt,halign=left]
\begin{itemize}[nosep,leftmargin=*,itemsep=1pt,font=\scriptsize]
\item thèse, arguments, exemples
\item visée : convaincre, persuader
\end{itemize}
\end{tcolorbox}
\begin{tcolorbox}[enhanced,colback=white,colframe=popgreen,boxrule=0.9pt,arc=3pt,title={Structure et réduction},fonttitle=\bfseries\scriptsize,coltitle=white,colbacktitle=popgreen,left=4pt,right=4pt,top=2pt,bottom=2pt,boxsep=1pt,toptitle=1.5pt,bottomtitle=1.5pt,halign=left]
\begin{itemize}[nosep,leftmargin=*,itemsep=1pt,font=\scriptsize]
\item repérer le plan
\item supprimer, généraliser, condenser
\end{itemize}
\end{tcolorbox}
\begin{tcolorbox}[enhanced,colback=white,colframe=poppurple,boxrule=0.9pt,arc=3pt,title={Reformulation},fonttitle=\bfseries\scriptsize,coltitle=white,colbacktitle=poppurple,left=4pt,right=4pt,top=2pt,bottom=2pt,boxsep=1pt,toptitle=1.5pt,bottomtitle=1.5pt,halign=left]
\begin{itemize}[nosep,leftmargin=*,itemsep=1pt,font=\scriptsize]
\item synonymes, nominalisation
\item fidélité sans copie
\end{itemize}
\end{tcolorbox}
\begin{tcolorbox}[enhanced,colback=white,colframe=poppink,boxrule=0.9pt,arc=3pt,title={Tonalités},fonttitle=\bfseries\scriptsize,coltitle=white,colbacktitle=poppink,left=4pt,right=4pt,top=2pt,bottom=2pt,boxsep=1pt,toptitle=1.5pt,bottomtitle=1.5pt,halign=left]
\begin{itemize}[nosep,leftmargin=*,itemsep=1pt,font=\scriptsize]
\item comique, dramatique, satirique
\end{itemize}
\end{tcolorbox}
\begin{tcolorbox}[enhanced,colback=white,colframe=popgold,boxrule=0.9pt,arc=3pt,title={Lecture méthodique},fonttitle=\bfseries\scriptsize,coltitle=white,colbacktitle=popgold,left=4pt,right=4pt,top=2pt,bottom=2pt,boxsep=1pt,toptitle=1.5pt,bottomtitle=1.5pt,halign=left]
\begin{itemize}[nosep,leftmargin=*,itemsep=1pt,font=\scriptsize]
\item \emph{Le Lion et la perle} n° 4 et 5
\end{itemize}
\end{tcolorbox}
\begin{tcolorbox}[enhanced,colback=white,colframe=popblue!60,boxrule=0.9pt,arc=3pt,title={Lexique et oral},fonttitle=\bfseries\scriptsize,coltitle=white,colbacktitle=popblue!60,left=4pt,right=4pt,top=2pt,bottom=2pt,boxsep=1pt,toptitle=1.5pt,bottomtitle=1.5pt,halign=left]
\begin{itemize}[nosep,leftmargin=*,itemsep=1pt,font=\scriptsize]
\item paronymie, antonymie
\item exposé : préparer, présenter
\end{itemize}
\end{tcolorbox}
\end{tcbraster}

\legende{Carte mentale de la leçon : formuler les idées essentielles d'un texte à résumer.}
\end{popfigure}

\begin{bilan}
\textbf{Définitions clés}
\begin{itemize}
\item Le \cle{texte argumentatif} : texte qui défend une \cle{thèse} (point de vue) à l'aide d'\cle{arguments} (raisons) illustrés par des \cle{exemples}, pour convaincre ou persuader le lecteur.
\item La \cle{contraction de texte} (résumé) : réduction fidèle d'un texte qui en conserve les idées essentielles, sans commentaire personnel ni copie de phrases.
\item L'\cle{idée essentielle} : idée principale d'un paragraphe ou d'une partie, sans laquelle le raisonnement s'effondre.
\item La \cle{reformulation} : dire autrement, avec ses propres mots, sans trahir le sens.
\item Les \cle{tonalités} : registre \cle{comique} (faire rire), \cle{dramatique} (émouvoir, conflit grave), \cle{satirique} (dénoncer par le rire et l'ironie).
\item La \cle{paronymie} : mots proches par la forme mais de sens différents (\emph{éminent} / \emph{imminent}). L'\cle{antonymie} : mots de sens contraires (\emph{généreux} / \emph{avare}).
\end{itemize}

\textbf{Repères à connaître par cœur}
\begin{center}
\begin{tabularx}{\linewidth}{|l|X|}\hline
\rowcolor{popblueL} \textbf{Notion} & \textbf{À retenir} \\ \hline
Étapes du résumé & Lire (2 fois) $\rightarrow$ dégager le thème et la thèse $\rightarrow$ repérer le plan $\rightarrow$ sélectionner les idées essentielles $\rightarrow$ reformuler $\rightarrow$ rédiger et vérifier. \\ \hline
Techniques de réduction & Supprimer (exemples, répétitions, détails), généraliser (un terme pour plusieurs), condenser (une phrase pour un paragraphe). \\ \hline
Reformulation & Synonymes, changement de construction, nominalisation (verbe $\rightarrow$ nom), voix passive/active. \\ \hline
Connecteurs logiques & D'abord, ensuite, enfin ; car, puisque ; mais, cependant ; donc, ainsi. \\ \hline
Exposé oral & Introduction (sujet + plan), développement (parties annoncées), conclusion (bilan + ouverture) ; regarder l'auditoire, parler clairement. \\ \hline
\end{tabularx}
\end{center}

\textbf{Méthodes express}
\begin{itemize}
\item Souligner les connecteurs : ils balisent la structure argumentative.
\item Une idée essentielle par paragraphe ; un exemple ne se résume pas, il se supprime.
\item Interdiction absolue : recopier des phrases du texte, donner son avis, dépasser le nombre de mots imposé.
\item Pour l'exposé : fiches courtes, entraînement à voix haute, gestion du temps.
\end{itemize}
\end{bilan}

\begin{aretenir}[Checklist avant le Bac]
\begin{itemize}
\item $\square$ Je sais définir le texte argumentatif (thèse, arguments, exemples).
\item $\square$ Je distingue convaincre (raison) et persuader (sentiments).
\item $\square$ Je repère la thèse et le plan d'un texte en deux lectures.
\item $\square$ Je maîtrise les trois techniques de réduction : supprimer, généraliser, condenser.
\item $\square$ Je reformule sans copier : synonymes, nominalisation, changement de construction.
\item $\square$ Je respecte la consigne de longueur et je ne donne jamais mon avis dans un résumé.
\item $\square$ Je reconnais les tonalités comique, dramatique et satirique dans un extrait.
\item $\square$ Je distingue paronymes et antonymes et je sais les employer correctement.
\item $\square$ Je connais les grandes étapes des lectures méthodiques n° 4 et 5 du \emph{Lion et la perle}.
\item $\square$ Je sais préparer et présenter un exposé : plan, fiches, posture, voix.
\item $\square$ Je relis ma copie : orthographe, accords, ponctuation, connecteurs logiques.
\end{itemize}
\end{aretenir}

\findecours

\end{document}
